% !TeX root = main-vnmm-mlpg.tex
%% 
%% Manuscript for Engineering Analysis with Boundary Elements (EABE)
%% Title: A Truly Cell-Free Vector Nodal Meshless Method via Meshless Local Petrov-Galerkin (VNMM-MLPG) for 2D Electromagnetic Problems
%% 
\documentclass[preprint,12pt]{elsarticle}

\usepackage[utf8]{inputenc}
\usepackage[T1]{fontenc}
\usepackage{lmodern}
\usepackage{amssymb}
\usepackage{amsmath}
\usepackage{amsfonts}
\usepackage{amsthm}
%\usepackage{lineno}
\usepackage{graphicx}
\usepackage{tabularx}
\usepackage{booktabs}
\usepackage{bm}
\usepackage{microtype}
\usepackage{float}
\usepackage{subcaption}
\usepackage[pdfpagelabels=false]{hyperref}

\hypersetup{
    colorlinks=true,
    linkcolor=blue,
    citecolor=blue,
    urlcolor=blue,
    pdfencoding=auto,
    hypertexnames=false
}

\newtheorem{theorem}{Theorem}
\newtheorem{proposition}{Proposition}
\newtheorem{remark}{Remark}
\newtheorem{lemma}{Lemma}

\renewcommand{\topfraction}{0.88}
\renewcommand{\bottomfraction}{0.80}
\renewcommand{\textfraction}{0.12}
\renewcommand{\floatpagefraction}{0.75}

\journal{Engineering Analysis with Boundary Elements}

\begin{document}

\begin{frontmatter}

    \title{A Truly Cell-Free Vector Nodal Meshless Method via Meshless Local Petrov--Galerkin (VNMM-MLPG) for 2D Electromagnetic Problems}

    \author[ufmg]{Renato C. Mesquita\corref{cor1}}
    \ead{renato@ufmg.br}
    \author[ufmg]{Williams L. de Nicomedes}
    \author[ifetcongonhas]{Bárbara M. F. Gonçalves}
    \author[ifmgsabara]{Luilly A. G. Ortiz}

    \cortext[cor1]{Corresponding author}

    \affiliation[ufmg]{organization={Department of Electrical Engineering, Universidade Federal de Minas Gerais},
        addressline={Av. Antônio Carlos, 6627},
        city={Belo Horizonte},
        postcode={31270-901},
        state={MG},
        country={Brazil}}

    \affiliation[ifetcongonhas]{organization={Mechanical Department, Instituto Federal de Educação, Ciência e Tecnologia de Minas Gerais, Campus Congonhas},
        addressline={Av. Michael Pereira de Souza, 3007},
        city={Congonhas},
        postcode={36417-050},
        state={MG},
        country={Brazil}}

    \affiliation[ifmgsabara]{organization={Industrial Control and Processes Department, Instituto Federal de Educação, Ciência e Tecnologia de Minas Gerais, Campus Sabará},
        addressline={Rodovia MGC-262, s/n, km 10},
        city={Sabará},
        postcode={34590-390},
        state={MG},
        country={Brazil}}

    \begin{abstract}
        The Vector Nodal Meshless Method (VNMM) approximates $H(\mathbf{curl})$ vector fields on unstructured point clouds using directional nodal degrees of freedom. However, previous implementations relied on background integration cells, introducing intra-cell stencil discontinuities, quadrature overhead, and compromised meshless autonomy. This paper presents a truly cell-free Vector Nodal Meshless Method via Meshless Local Petrov--Galerkin (VNMM-MLPG). Combining linear vector point interpolation with weak forms over autonomous circular subdomains decouples the spatial discretization from background grids. Because fixed linear stencils yield constant curl and divergence per subdomain, Stokes' and Gauss' theorems imply that axisymmetric weight test functions vanishing on the subdomain boundary (MLPG1) produce identically zero stiffness, motivating a Local Galerkin (MLPG6) approach. To eliminate the structural asymmetry and loss of coercivity of row-by-row MLPG assembly without lens-intersection quadratures, we formulate an overlapping \textit{Symmetric Subdomain Galerkin} scheme evaluated in closed form, yielding symmetric positive-semidefinite global operators. Furthermore, modal decomposition shows that single-disk integration leaves two harmonic zero-energy deformation modes---pure shear and hyperbolic---in the local div-curl kernel. We resolve this rank deficiency by projecting these modes onto directional degrees of freedom and introducing orthogonal energy stabilization grounded in Variational Multiscale (VMS) theory, paired with a resolution-scaled divergence penalty. Benchmarks spanning affine patch tests, manufactured solutions, and cavity eigenproblems on regular and perturbed clouds confirm exact linear reproduction, theoretical $L^2$ field and $H(\mathbf{curl})$ semi-norm convergence rates, effective shifting of spurious gradient modes outside the operational spectrum, and higher accuracy, operator sparsity, and assembly speed compared to cell-based quadrature.
    \end{abstract}

    \begin{keyword}
        Vector Nodal Meshless Method (VNMM) \sep Meshless Local Petrov--Galerkin (MLPG) \sep Point Interpolation Method (PIM) \sep Truly cell-free methods \sep Computational electromagnetics \sep Resonant cavities
    \end{keyword}

\end{frontmatter}

%\linenumbers

%% =========================================================================
%% SECTION 1: INTRODUCTION
%% =========================================================================
\section{Introduction}
\label{sec:introduction}

Computational electromagnetics (CEM) has traditionally relied on the finite element method (FEM) with mixed edge elements \cite{nedelec1980mixed, monk2003finite, jin2014finite} to resolve vector boundary-value and eigenvalue problems in $H(\mathbf{curl})$. Edge elements assign tangential field components to elemental edges, enforcing trace continuity while naturally accommodating jump discontinuities across dielectric interfaces and satisfying the exact sequence of the de Rham complex. Consequently, irrotational gradient fields reside in the exact discrete nullspace of the curl operator, $\mathrm{grad} \, W_h \subset \ker(\mathrm{curl}_h)$, preventing the emergence of non-physical spurious eigenmodes \cite{bossavit1998computational, boffi2006finite}. Nevertheless, edge elements are constrained by the necessity of generating and maintaining conforming simplex meshes (triangles in 2D, tetrahedra in 3D). Mesh generation remains the primary computational bottleneck in complex CAD geometries, multi-scale devices, moving boundary problems, and topology optimization routines.

To circumvent mesh generation bottlenecks, meshless and meshfree methods have garnered sustained interest across solid mechanics, fluid dynamics, and wave physics \cite{belytschko2023meshfree, liu2009mesh, fasshauer2007meshfree, atluri2004meshless}. In computational electromagnetics, early meshless adaptations employed scalar formulations for electric scalar potentials, magnetic vector potentials, or out-of-plane electric fields governed by scalar Helmholtz equations \cite{viana1999moving, ho2001moving, parreira2006element, nicomedes2012meshless}. Extending meshless concepts to vector fields governed by Maxwell's curl-curl equations proved substantially more challenging. Early vector adaptations attempted to mimic finite element edge conformity by assigning degrees of freedom to pairs of nodes forming virtual edges, as in the Edge Meshless Method (EMM) \cite{lima2017edge}. While effective at suppressing spurious modes, such schemes reintroduce topological connectivity between node pairs, partially compromising the autonomy of the point cloud. On the other hand, purely nodal collocated schemes with vector degrees of freedom typically induce severe spurious mode contamination and derivative aliasing due to the lack of discrete de Rham complex exactness on arbitrary point clouds \cite{bossavit1990spurious, boffi2006finite}.

The Vector Nodal Meshless Method (VNMM), introduced by Ortiz et al.~\cite{ortiz2021nodal, ortiz2023tese} and Gonçalves et al.~\cite{goncalves2023vector, goncalves2025tese}, provides a node-based vector interpolation strategy. Fundamentally, the VNMM can be recognized as the asymptotic limit of the EMM as the virtual edge length tends to zero: by collapsing the node pair into a single evaluation point $\mathbf{x}_I$, the edge circulation reduces to a local directional projection along a unit director $\vec{t}_I$, restoring point autonomy without pairwise connectivity. From an interpolation standpoint, the VNMM constructs shape functions via local polynomial point collocation on support stencils, conceptually generalizing the Point Interpolation Method (PIM) \cite{liu2001point, liu2005introduction} to vector fields endowed with directional degrees of freedom $e_I = \vec{E}(\mathbf{x}_I) \cdot \vec{t}_I$.

In a recent study \cite{stable_vnmm2d_2026}, the mathematical foundations of the 2D VNMM were systematically examined on unconstrained irregular point clouds. That study resolved several fundamental questions:
\begin{enumerate}
    \item \textbf{Incomplete Basis Stagnation:} The classical incomplete solenoidal vector basis ($\mathcal{L}^1$, consisting of 3 terms per stencil) was shown to suffer from a persistent $\mathcal{O}(1)$ error plateau in curl approximation due to dilatational derivative aliasing on irregular clouds. This was resolved by establishing the complete linear vector space $\mathcal{P}_1 \times \mathcal{P}_1$ (requiring a 6-node stencil, denoted $\mathcal{P}^1$), which restored full linear completeness and $\mathcal{O}(h)$ curl convergence.
    \item \textbf{Geometric Singularity Mitigation:} A scale-invariant determinant qualification criterion on a dimensionless placement matrix $\widehat{\mathbf{A}}$, coupled with adaptive $k$-d tree searches, was developed to automatically reject nearly singular node and director configurations without empirical parameter tuning.
    \item \textbf{Variational Div-Curl Stabilization:} Because dynamic support stencils across point clouds lack global trace conformity, discrete gradients leak into the positive curl spectrum. A regularized Ritz--Galerkin formulation with a volumetric div-curl penalty ($s_{\mathrm{div}}$) was introduced, successfully shifting spurious modes to higher frequencies while preserving physical cavity resonances.
\end{enumerate}

Despite these substantial advances, the formulation reported in \cite{stable_vnmm2d_2026} left open a fundamental deficiency: \textbf{the reliance on a Cartesian background integration grid}. As classified by Belytschko et al.~\cite{belytschko1994element} and Atluri and Shen \cite{atluri2002meshless}, methods like the Element-Free Galerkin (EFG) are not \textit{truly cell-free} because evaluation of the global weak form requires an auxiliary background mesh of quadrature cells covering the domain $\Omega$. In the VNMM context, this background scaffolding creates acute numerical challenges:
\begin{itemize}
    \item \textbf{Intra-Cell Stencil Discontinuities:} In cell-based integration, the 6-node support stencil is selected point-by-point at each Gauss integration point $P_g$. Because nearest-neighbor stencils switch abruptly within the cell interior, the resulting shape functions and derivatives are piecewise discontinuous across intra-cell boundaries. As demonstrated in \cite{stable_vnmm2d_2026}, under-integrating these cells (e.g., using $2 \times 2$ Gauss points) produces artificial numerical softening that conceals rank deficiencies in the discrete curl operator. Suppressing this artifact required a dense, resolution-matched quadrature ($3 \times 3$ Gauss points per cell, exceeding 3,600 quadrature evaluations for a 441-node grid).
    \item \textbf{Loss of True Meshless Autonomy:} Although background integration grids do not require the strict element quality metrics (such as shape regularity or aspect-ratio limits) demanded by finite element meshes, constructing background cells and trimming cut cells along complex, curved, or 3D boundaries still introduces substantial geometric pre-processing overhead and quadrature complexity, thereby compromising the cell-free autonomy of meshless methods.
\end{itemize}

Several alternative domain integration strategies have been investigated in the broader meshfree literature \cite{hillman2018performance}, but each presents specific drawbacks when applied to unstructured vector problems:
\begin{itemize}
    \item \textbf{Direct Nodal Integration (DNI) and Rank Instabilities:} The simplest meshless strategy evaluates the global weak form directly at the nodes. However, as demonstrated by Babu{\v{s}}ka et al.~\cite{babuska2008quadrature} and Hillman and Chen~\cite{hillman2018performance}, DNI underestimates the stored electromagnetic energy of low-order modes, leading to rank deficiencies (sawtooth zero-energy modes) and degraded convergence rates without extensive ad-hoc stabilization.
    \item \textbf{Conforming vs. Non-Conforming Nodal Smoothing (SCNI and SNNI):} Stabilized Conforming Nodal Integration (SCNI) \cite{chen2001stabilized} restores stability and patch test exactness by computing smoothed gradients via boundary integrals over conforming Voronoi dual cells, preserving matrix symmetry \cite{hillman2018performance}. However, constructing Voronoi diagrams on irregular point clouds is geometrically brittle, sensitive to boundary slivers, and computationally demanding in 3D. While Stabilized Non-Conforming Nodal Integration (SNNI) dispenses with Voronoi meshing by using non-conforming local boxes \cite{hillman2018performance}, the loss of partition conformity fails the patch test, necessitating Petrov--Galerkin variational consistency (VC) corrections \cite{chen2013arbitrary, hillman2018performance} that require auxiliary local equation solves.
    \item \textbf{Strong-Form Nodal Collocation:} Enforcing the governing PDE directly at the nodes avoids numerical quadrature \cite{hillman2018performance}. However, for second-order PDEs, collocation requires computing high-order derivatives of shape functions and demands at least quadratic polynomial completeness for convergence \cite{hillman2018performance}. More fundamentally, for vector electromagnetic problems governed by Maxwell's curl-curl equations, strong-form collocation lacks a discrete de Rham complex structure, triggering spectral corruption and unphysical gradient spurious modes \cite{bossavit1990spurious, boffi2006finite}.
    \item \textbf{Discontinuous Galerkin with Trace Jump Penalties:} One could view shifting stencils as non-conforming elements and introduce interface integrals penalizing tangential trace jumps, as in interior penalty discontinuous Galerkin (IPDG) formulations for electromagnetics \cite{houston2004mixed, hesthaven2004high}. However, because point-interpolation stencils lack well-defined geometric element boundaries or compact influence domains, identifying and tracking the complex internal discontinuity interfaces across an unstructured point cloud is geometrically ambiguous and computationally intractable.
    \item \textbf{Autonomous Local Weak Forms (MLPG and LPIM):} Formulations based on local subdomains, such as the Meshless Local Petrov--Galerkin (MLPG) method of Atluri and Zhu \cite{atluri1998new, atluri2002meshless} and the Local Point Interpolation Method (LPIM) of Liu and Gu \cite{liu2001local, liu2005introduction}, circumvent background cells entirely by establishing weak forms over autonomous circular disks or spheres $\Omega_s^I$ centered at each node.
\end{itemize}

While MLPG and LPIM have been widely applied to scalar mechanics, acoustics, and electromagnetics \cite{sladek2001local, atluri2002meshless, liu2001local,nicomedes2012meshless}, their extension to vector electromagnetic formulations in $H(\mathbf{curl})$ presents specific theoretical and algebraic challenges:
\begin{enumerate}
    \item How should vector test functions $\vec{W}_I$ be chosen on $\Omega_s^I$? In MLS-based MLPG1 \cite{atluri1998new, atluri2002meshless} and radial-basis LRPIM \cite{liu2005introduction}, localized weight functions vanishing on the subdomain boundary ($w(r_s) = 0$) are commonly used because the trial derivatives vary spatially across $\Omega_s^I$. For a fixed polynomial stencil ($\mathcal{P}^1$) \cite{liu2001local, stable_vnmm2d_2026}, however, the curl and divergence are spatially constant on $\Omega_s^I$, so test functions vanishing on $\partial\Omega_s^I$ cause the local stiffness bilinear forms to vanish identically by Stokes' and Gauss' theorems.
    \item How can global matrix symmetry and positive semidefiniteness be guaranteed without lens-intersection quadratures? In the literature on MLPG \cite{atluri2002meshless, nicomedes2012meshless} and LPIM \cite{liu2001local, liu2005introduction}, equations are assembled in a row-by-row Petrov--Galerkin manner (one test function per node). Because the integration domains differ across nodes ($\Omega_s^I \ne \Omega_s^J$), the resulting global stiffness matrix is intrinsically unsymmetric ($K_{Ij} \ne K_{Ji}$). For conservative electromagnetic cavities, this loss of discrete self-adjointness can induce unphysical complex eigenvalues \cite{saad2011numerical, boffi2013mixed} and causes a loss of discrete coercivity under refinement for general non-polynomial fields. Conversely, symmetric overlapping-sphere methods such as the Method of Finite Spheres \cite{de2000method, nicomedes2021acoustic} require specialized numerical cubature over lens-shaped intersections $\Omega_s^I \cap \Omega_s^J$.
    \item How should boundary conditions and intersections between local subdomains and physical domain boundaries be treated? When a local disk intersects the physical boundary $\partial\Omega$, the integration domain is truncated into a circular segment $\Omega_s^B = \Omega_s \cap \Omega$, posing the challenge of evaluating spatial integrals on truncated geometries in closed form while leveraging the directional Kronecker-delta property inherited from VNMM \cite{ortiz2021nodal, goncalves2023vector, stable_vnmm2d_2026} to avoid singular boundary penalty integrals.
    \item How can discrete mode softening and local operator rank deficiency be resolved? Evaluating linear vector fields on a single circular disk subdomain produces spatially constant curl and divergence, causing the local stiffness contributions $\mathbf{K}_{\mathrm{curl}}^I$ and $\mathbf{K}_{\mathrm{div}}^I$ to each possess rank 1. Consequently, their sum $\mathbf{K}_{\mathrm{curl}}^I + s_{\mathrm{div}}\mathbf{K}_{\mathrm{div}}^I$ has rank 2 over the four-dimensional translation-free linear deformation subspace (whereas the local mass matrix $\mathbf{M}^I$ has full rank 6), leaving two harmonic zero-energy linear deformation modes (pure shear and hyperbolic) unpenalized within each stencil. Unlike cell-based $3\times 3$ Gauss quadrature (which couples adjacent stencils within each cell), single-disk integration allows these modes to soften, causing error stagnation and sublinear convergence unless properly stabilized.
\end{enumerate}

In this paper, we address these challenges by formulating, implementing, and validating the \textbf{VNMM-MLPG}---a truly cell-free vector nodal meshless method for 2D electromagnetics. By combining the directional $\mathcal{P}^1$ Vector Nodal Meshless Method \cite{ortiz2021nodal, stable_vnmm2d_2026} with an overlapping \textbf{Symmetric Subdomain Galerkin Assembly} within the Local Galerkin (MLPG6) framework \cite{atluri2002meshless}, the main contributions of this work are:
\begin{itemize}
    \item \textbf{Closed-Form Analytical Discretization of VNMM Operators:} Exact closed-form factorization of all local curl, divergence, and mass matrices in terms of the directional inverse collocation matrix $\boldsymbol{\beta}^I = \mathbf{A}_I^{-1}$ and analytical geometric moments over interior disks and boundary-truncated circular segments, eliminating numerical quadrature error and lens-intersection cubatures.
    \item \textbf{Symmetric Subdomain Galerkin Assembly and Coercivity Analysis:} Formulation of an additive overlapping-subdomain Ritz--Galerkin energy assembly that guarantees strictly symmetric positive-semidefinite global operators $\mathbf{K}$ and $\mathbf{M}$, together with a direct numerical demonstration that classical row-by-row MLPG/LPIM assembly loses discrete coercivity under refinement for non-polynomial $H(\mathbf{curl})$ fields.
    \item \textbf{Modal Rank Decomposition and VMS Orthogonal Energy Stabilization:} Analytical projection of the two harmonic zero-energy deformation modes ($\mathbf{v}_{\mathrm{shear}}^I, \mathbf{v}_{\mathrm{hyp}}^I$) onto the VNMM directional degrees of freedom, and derivation of a consistent Variational Multiscale (VMS) stabilization operator ($\alpha_{\mathrm{stab}} = c_{\mathrm{stab}}(h_{\mathrm{nom}}/L)^2$) with superlinear penalty scaling ($s_{\mathrm{div}} \propto h_{\mathrm{nom}}^{-1.5}$). This restores full rank 4 over the local linear deformation subspace, eliminates mode softening, and achieves superlinear-to-quadratic convergence ($p = 2.15\text{--}2.39$ on regular grids and $p = 1.62$ under $25\%$ nodal jitter).
    \item \textbf{Unified Boundary Treatment on Truncated Subdomains:} Integration of VNMM's directional Kronecker-delta collocation \cite{ortiz2021nodal, stable_vnmm2d_2026} with truncated-disk MLPG6 weak forms, eliminating essential Dirichlet rows/columns without Nitsche or Lagrange multiplier boundary integrals while reducing inhomogeneous Neumann (validated via mixed MMS benchmarks) and Robin conditions to simple 1D boundary segment integrals.
    \item \textbf{Accuracy, Sparsity, and Assembly Speedup over Cell-Based Quadrature:} Comprehensive verification across patch tests, manufactured solutions, subdomain containment analysis (confirming the classical MLPG/LPIM scaling $\alpha_s = 0.80$ \cite{atluri2002meshless, liu2009mesh} within the 6-node vector stencil containment bound $\alpha_s \le 1.15$), and PEC cavity eigenproblems, achieving $2.0\times\text{--}3.7\times$ faster matrix assembly, $20\%\text{--}24\%$ sparser matrices, and higher spectral accuracy than cell-based VNMM-Gauss.
\end{itemize}

The remainder of this paper is structured as follows. Section~\ref{sec:formulation} details the mathematical formulation of VNMM-MLPG, proves the breakdown of classical weight test functions, derives the closed-form Local Galerkin (MLPG6) matrices, contrasts row-by-row assembly against the symmetric Subdomain Galerkin assembly, and develops the Dimensionless Orthogonal Energy Stabilization for local rank restoration. Section~\ref{sec:results} presents the validation roadmap and numerical results across six comprehensive stages: closed-form elemental verification, patch test satisfaction, MMS asymptotic convergence rates, theoretical determination of $\alpha_s$, resonant cavity spectra with stabilization, and computational benchmarking. Section~\ref{sec:discussion} discusses advantages, current limitations, and open challenges. Section~\ref{sec:conclusions} concludes the paper.

%% =========================================================================
%% SECTION 2: THEORETICAL FORMULATION OF VNMM-MLPG
%% =========================================================================
\section{Theoretical Formulation of the VNMM-MLPG}
\label{sec:formulation}

\subsection{Continuous Boundary-Value Problem}
\label{subsec:continuous_bvp}

Consider a bounded two-dimensional domain $\Omega \subset \mathbb{R}^2$ with a Lipschitz boundary $\partial\Omega$, occupied by a linear, isotropic, non-dispersive medium characterized by relative permeability $\mu_r$ and permittivity $\epsilon_r$. The time-harmonic electric field $\vec{E}(x, y) = [E_x, E_y]^T$ satisfies Maxwell's curl-curl equation. To shift gradient-type spurious modes outside the operational frequency band on non-conforming meshless spaces, we adopt the div-curl regularized vector Helmholtz equation \cite{stable_vnmm2d_2026}:
\begin{equation}
    \nabla \times \left(\frac{1}{\mu_r} \nabla \times \vec{E}\right) - \frac{s_{\mathrm{div}}}{\mu_r} \nabla(\nabla \cdot \vec{E}) - k_0^2 \epsilon_r \vec{E} = \vec{J}_s \quad \text{in } \Omega,
    \label{eq:helmholtz_pde}
\end{equation}
where $k_0 = \omega\sqrt{\mu_0\epsilon_0}$ is the free-space wavenumber, $s_{\mathrm{div}} \ge 0$ is the dimensionless divergence regularization parameter, and $\vec{J}_s$ is an applied electric current density. In two dimensions with in-plane vector fields $\vec{E} = [E_x, E_y]^T$, the differential operators are defined by:
\begin{equation}
    \nabla \times \vec{E} = \left(\frac{\partial E_y}{\partial x} - \frac{\partial E_x}{\partial y}\right) \hat{\mathbf{z}}, \qquad \nabla \cdot \vec{E} = \frac{\partial E_x}{\partial x} + \frac{\partial E_y}{\partial y}.
\end{equation}

To establish a well-posed boundary-value problem, the physical boundary is partitioned into pairwise disjoint segments $\partial\Omega = \Gamma_D \cup \Gamma_N \cup \Gamma_R$ (with $\Gamma_D \cap \Gamma_N = \Gamma_D \cap \Gamma_R = \Gamma_N \cap \Gamma_R = \emptyset$), where $\vec{n}$ and $\vec{t}$ denote the outward unit normal and unit tangent to $\partial\Omega$, respectively:
\begin{enumerate}
    \item \textbf{Dirichlet Boundary Conditions ($\Gamma_D$):} The tangential trace of the electric field is prescribed:
          \begin{equation}
              \vec{n} \times \vec{E} = \vec{g}_D \quad \iff \quad \vec{E} \cdot \vec{t} = g_D \quad \text{on } \Gamma_D.
              \label{eq:pec_bc}
          \end{equation}
          When $g_D = 0$, this represents the homogeneous Dirichlet condition on perfectly conducting (PEC) electric walls ($\Gamma_{\mathrm{PEC}}$). When $g_D \ne 0$, it models inhomogeneous Dirichlet excitations, such as prescribed incident wave profiles or driven waveguide apertures.
    \item \textbf{Neumann Boundary Conditions ($\Gamma_N$):} The natural boundary flux (tangential magnetic field / surface curl) is specified:
          \begin{equation}
              \frac{1}{\mu_r} \left(\vec{n} \times (\nabla \times \vec{E})\right) = \vec{g}_N \quad \text{on } \Gamma_N.
              \label{eq:neumann_bvp_pde}
          \end{equation}
          When $\vec{g}_N = \mathbf{0}$, this represents the homogeneous Neumann condition on perfect magnetic conductors (PMC) or electromagnetic symmetry planes. When $\vec{g}_N \ne \mathbf{0}$, it accounts for impressed surface magnetic current densities $\vec{M}_s = -\vec{g}_N$.
    \item \textbf{Robin / Absorbing Boundary Conditions ($\Gamma_R$):} To truncate open computational domains or model surface impedance, first-order absorbing boundary conditions (ABC) of the Silver--M{\"u}ller type are imposed:
          \begin{equation}
              \frac{1}{\mu_r} \left(\vec{n} \times (\nabla \times \vec{E})\right) + \gamma \vec{E}_t = \vec{g}_R \quad \text{on } \Gamma_R,
              \label{eq:robin_bvp_pde}
          \end{equation}
          where $\gamma = j k_0 \sqrt{\epsilon_r / \mu_r}$ is the surface admittance, $\vec{E}_t = (\vec{E} \cdot \vec{t})\vec{t}$ is the tangential electric field vector along the boundary unit tangent $\vec{t}$, and $\vec{g}_R$ represents incoming plane or guided waves.
\end{enumerate}

\subsection{Nodal Discretization and Local Support Geometries}
\label{subsec:nodal_discretization}

The domain $\overline{\Omega} = \Omega \cup \partial\Omega$ is discretized by a cloud of $N$ nodes $\mathcal{S} = \{P_I\}_{I=1}^N$ with Cartesian coordinates $\mathbf{x}_I = [x_I, y_I]^T$. Each node $P_I$ is endowed with an assigned unit director vector $\vec{t}_I = [t_{x,I}, t_{y,I}]^T$, defining a scalar degree of freedom (DoF):
\begin{equation}
    e_I = \vec{E}(\mathbf{x}_I) \cdot \vec{t}_I.
\end{equation}
Along physical boundaries with prescribed tangential electric fields, directors are oriented strictly tangent to the boundary ($\vec{t}_I \parallel \partial\Omega$). On perfectly conducting walls ($\mathbf{x}_I \in \Gamma_{\mathrm{PEC}}$), condition \eqref{eq:pec_bc} then translates directly into the homogeneous constraint $e_I = 0$, whereas on inhomogeneous Dirichlet boundaries ($\Gamma_D$), it permits direct collocation of the prescribed tangential profile ($e_I = g_D(\mathbf{x}_I)$). Likewise, along material interfaces between distinct media, aligning directors parallel to the interface natively preserves the continuity of the tangential electric field while naturally accommodating physical jump discontinuities in the normal component.

In the MLPG framework, two distinct local domains are defined for each node $P_I$ (Fig.~\ref{fig:schematic_mlpg}):
\begin{enumerate}
    \item \textbf{Trial Support Domain ($\Omega_u^I$):} The spatial neighborhood containing an admissible stencil of $K = 6$ nodes used to construct the local trial vector field $\vec{E}^h(\mathbf{x})$.
    \item \textbf{Local Integration Subdomain ($\Omega_s^I$):} A circular disk of radius $r_s^I$ centered at $\mathbf{x}_I$:
          \begin{equation}
              \Omega_s^I = \left\{ \mathbf{x} \in \mathbb{R}^2 : \|\mathbf{x} - \mathbf{x}_I\| \le r_s^I \right\}, \qquad r_s^I = \alpha_s h_I,
              \label{eq:subdomain_disk}
          \end{equation}
          where $h_I \equiv h_{\mathrm{local}}(\mathbf{x}_I)$ denotes the local nodal spacing (distance to nearest neighbor or mean distance to adjacent nodes, in direct correspondence with $h_{\mathrm{local}}$ in \cite{stable_vnmm2d_2026}), and $\alpha_s > 0$ is a dimensionless radius scaling parameter. Throughout this entire work, the single universal value $\alpha_s = 0.80$ is adopted, as justified by the theoretical geometric bounds and parametric sensitivity analysis in Section~\ref{subsec:stage4_alpha_s}.
\end{enumerate}

A representative nodal distribution generated on a rectangular cavity with a $25\%$ random perturbation (jitter) is depicted in Fig.~\ref{fig:schematic_mlpg}. Interior nodes $\mathbf{x}_k$ are assigned arbitrary, randomized unit director orientations $\vec{t}_k$, whereas boundary nodes are constrained strictly tangential to $\partial\Omega$ ($\vec{t}_k \parallel \partial\Omega$). Depending on the spatial position of the central evaluation node relative to the physical boundary, three fundamental topological configurations of local integration subdomains $\Omega_s$ and trial support stencils $\Omega_u$ arise:
\begin{itemize}
    \item \textbf{Case 1 (Fully Interior Subdomain, $\Omega_s^{I_1} \subset \Omega$):} For an interior node $\mathbf{x}_{I_1}$ located sufficiently far from the physical boundary ($\mathrm{dist}(\mathbf{x}_{I_1}, \partial\Omega) \ge r_s^{I_1}$), the local integration domain is an uncut, complete circular disk.
    \item \textbf{Case 2 (Truncated Interior Subdomain, $\Omega_s^{I_2} \cap \Omega$):} For an interior node $\mathbf{x}_{I_2}$ positioned near the boundary ($\mathrm{dist}(\mathbf{x}_{I_2}, \partial\Omega) < r_s^{I_2}$), the local disk crosses $\partial\Omega$. Because trial and test fields are only defined on the physical domain, the active integration domain is strictly truncated to $\Omega_s^{I_2} \cap \Omega$, forming a circular segment whose boundary $\partial(\Omega_s^{I_2} \cap \Omega)$ is bounded by an interior circular arc $L_{sI_2} = \partial\Omega_s^{I_2} \cap \Omega$ and a straight boundary intersection segment $\Gamma_{sI_2} = \Omega_s^{I_2} \cap \partial\Omega$ (the 1D segment resulting from intersecting the 2D local disk $\Omega_s^{I_2}$ with the domain boundary $\partial\Omega$).
    \item \textbf{Case 3 (Truncated Boundary Subdomain, $\Omega_s^B \cap \Omega$):} For a node $\mathbf{x}_B \in \partial\Omega$ situated directly on the boundary, the local disk is centered on $\partial\Omega$ and truncated into an exact half-disk $\Omega_s^B \cap \Omega$. Its boundary consists of a semicircular interior arc $L_{sB} = \partial\Omega_s^B \cap \Omega$ and the straight boundary diameter segment $\Gamma_{sB} = \Omega_s^B \cap \partial\Omega$ (the 1D intersection between the 2D disk $\Omega_s^B$ and $\partial\Omega$), which constitutes the active segment where physical boundary conditions are enforced.
\end{itemize}
As established in Section~\ref{subsubsec:truncated_subdomains_analytic}, the constant curl and divergence properties of complete $\mathcal{P}^1$ shape functions guarantee that integration over truncated subdomains (Cases~2 and 3) is computed in exact closed analytical form via active area scaling $\mathbf{K}_I^{\mathrm{trunc}} = \eta_I \mathbf{K}_I^{\mathrm{full}}$ without requiring background numerical quadrature.
For all three cases, Fig.~\ref{fig:schematic_mlpg} also illustrates the real 6-node trial support stencils $\Omega_u$ selected by the scale-invariant determinant qualification algorithm (Section~\ref{subsec:stencil_selection}). Each stencil is indicated by open markers encircling the six participating nodes and dotted lines connecting them to the central evaluation node.

\begin{figure}[!t]
    \centering
    \includegraphics[width=0.72\textwidth]{fig_mlpg_subdomains_schematic.png}
    \caption{Schematic of the VNMM-MLPG nodal cloud, local integration subdomains $\Omega_s$ (drawn with an enlarged radius factor $\alpha_s = 1.10$ for clear visual illustration of near-boundary Case~2 truncation), and 6-node trial support stencils $\Omega_u$.}
    \label{fig:schematic_mlpg}
\end{figure}

\subsection{VNMM Vector Shape Functions and Field Interpolation via $\mathcal{P}^1$}
\label{subsec:trial_space}

Following \cite{stable_vnmm2d_2026}, within the local evaluation domain of node $P_I$, the vector field $\vec{E}(\mathbf{x})$ is approximated in the complete linear polynomial vector space $\mathcal{P}^1 = \mathcal{P}_1 \times \mathcal{P}_1$. This space is spanned by the basis matrix $\mathbf{P}(\Delta\mathbf{x}) \in \mathbb{R}^{2 \times 6}$:
\begin{equation}
    \mathbf{P}(\Delta\mathbf{x}) = \begin{bmatrix}
        1 & 0 & \Delta x & \Delta y & 0        & 0        \\
        0 & 1 & 0        & 0        & \Delta x & \Delta y
    \end{bmatrix},
    \label{eq:P_basis_matrix}
\end{equation}
with local coordinates $\Delta\mathbf{x} = [\Delta x, \Delta y]^T = [x - x_I, y - y_I]^T$ centered at node $P_I$.
Each local vector shape function $\vec{N}_a^I(\mathbf{x})$ associated with the $a$-th node of the support stencil $\Omega_u^I$ ($a = 1, \dots, 6$) is expressed directly as a linear combination of these basis terms:
\begin{equation}
    \vec{N}_a^I(\mathbf{x}) = \mathbf{P}(\mathbf{x} - \mathbf{x}_I) \boldsymbol{\beta}_a^I,
    \label{eq:shape_functions}
\end{equation}
where $\boldsymbol{\beta}_a^I = [\beta_{1a}^I, \dots, \beta_{6a}^I]^T \in \mathbb{R}^6$ is the vector of polynomial interpolation coefficients for shape function $a$. In explicit component form matching \cite{stable_vnmm2d_2026}:
\begin{equation}
    \begin{aligned}
        \vec{N}_a^I(x, y) = & \beta_{1a}^I \begin{bmatrix} 1 \\ 0 \end{bmatrix} + \beta_{2a}^I \begin{bmatrix} 0 \\ 1 \end{bmatrix} + \beta_{3a}^I \begin{bmatrix} x - x_I \\ 0 \end{bmatrix}                \\
                            & + \beta_{4a}^I \begin{bmatrix} y - y_I \\ 0 \end{bmatrix} + \beta_{5a}^I \begin{bmatrix} 0 \\ x - x_I \end{bmatrix} + \beta_{6a}^I \begin{bmatrix} 0 \\ y - y_I \end{bmatrix}.
    \end{aligned}
    \label{eq:P1_shape_function_expansion_mlpg}
\end{equation}

Enforcing the vector Kronecker-delta projection property \cite{stable_vnmm2d_2026}:
\begin{equation}
    \vec{N}_a^I(\mathbf{x}_k) \cdot \vec{t}_k = \delta_{ak}, \qquad k = 1, \dots, 6,
    \label{eq:kronecker_projection}
\end{equation}
across all 6 stencil nodes yields the local linear placement system for $\boldsymbol{\beta}_a^I$:
\begin{equation}
    \mathbf{A}_I \boldsymbol{\beta}_a^I = \mathbf{L}_a,
    \label{eq:collocation_system}
\end{equation}
where $\mathbf{L}_a = [\delta_{1a}, \dots, \delta_{6a}]^T$ is the $a$-th canonical basis vector of $\mathbb{R}^6$, and the $k$-th row of the local moment placement matrix $\mathbf{A}_I \in \mathbb{R}^{6 \times 6}$ is given explicitly by:
\begin{equation}
    (\mathbf{A}_I)_{k, :} = \vec{t}_k^T \mathbf{P}(\mathbf{x}_k - \mathbf{x}_I) = \begin{bmatrix} t_{kx} & t_{ky} & \Delta x_k t_{kx} & \Delta y_k t_{kx} & \Delta x_k t_{ky} & \Delta y_k t_{ky} \end{bmatrix},
    \label{eq:matrix_A_row}
\end{equation}
following the formulation in \cite{stable_vnmm2d_2026}.
Once an admissible stencil $\Omega_u^I$ is selected such that $\mathbf{A}_I$ is non-singular (Section~\ref{subsec:stencil_selection}), grouping all six systems together yields the compact matrix formulation:
\begin{equation}
    \mathbf{A}_I \boldsymbol{\beta}^I = \mathbf{I}_{6 \times 6} \implies \boldsymbol{\beta}^I = \mathbf{A}_I^{-1},
    \label{eq:beta_matrix_inversion}
\end{equation}
where $\boldsymbol{\beta}^I = [\boldsymbol{\beta}_1^I, \dots, \boldsymbol{\beta}_6^I] \in \mathbb{R}^{6 \times 6}$ is the local coefficient matrix. The interpolated electric vector field associated with the stencil of node $P_I$ is then expressed directly as:
\begin{equation}
    \vec{E}^h(\mathbf{x}) = \sum_{a=1}^6 \vec{N}_a^I(\mathbf{x}) e_a = \mathbf{P}(\mathbf{x} - \mathbf{x}_I) \boldsymbol{\beta}^I \mathbf{e}_{\mathrm{stencil}},
    \label{eq:field_interpolation_mlpg}
\end{equation}
where $\mathbf{e}_{\mathrm{stencil}} \equiv \mathbf{L}_I \mathbf{e} = [e_1, \dots, e_6]^T \in \mathbb{R}^6$ collects the scalar directional degrees of freedom ($e_a = \vec{E}(\mathbf{x}_a) \cdot \vec{t}_a$) via the local Boolean gathering matrix $\mathbf{L}_I \in \mathbb{R}^{6 \times N}$.

Because the vector shape functions $\vec{N}_a^I$ are linear polynomials in Cartesian coordinates, their curl and divergence are strictly constant everywhere inside the local subdomain $\Omega_s^I$:
\begin{align}
    (\nabla \times \vec{N}_a^I)_z & = \frac{\partial N_{ay}^I}{\partial x} - \frac{\partial N_{ax}^I}{\partial y} = \beta_{5a}^I - \beta_{4a}^I \equiv c_{\mathrm{curl}, a}^I \in \mathbb{R}, \label{eq:curl_constant} \\
    \nabla \cdot \vec{N}_a^I      & = \frac{\partial N_{ax}^I}{\partial x} + \frac{\partial N_{ay}^I}{\partial y} = \beta_{3a}^I + \beta_{6a}^I \equiv c_{\mathrm{div}, a}^I \in \mathbb{R}. \label{eq:div_constant}
\end{align}
Collecting these six nodal operator coefficients into the discrete operator column vectors $\mathbf{c}_{\mathrm{curl}}^I \equiv \mathbf{v}_{\mathrm{curl}}^I = [c_{\mathrm{curl}, 1}^I, \dots, c_{\mathrm{curl}, 6}^I]^T = (\boldsymbol{\beta}_{5, :}^I - \boldsymbol{\beta}_{4, :}^I)^T \in \mathbb{R}^6$ and $\mathbf{c}_{\mathrm{div}}^I \equiv \mathbf{v}_{\mathrm{div}}^I = [c_{\mathrm{div}, 1}^I, \dots, c_{\mathrm{div}, 6}^I]^T = (\boldsymbol{\beta}_{3, :}^I + \boldsymbol{\beta}_{6, :}^I)^T \in \mathbb{R}^6$, the rotational and divergent fields of the interpolated vector field $\vec{E}^h(\mathbf{x})$ within $\Omega_s^I$ are evaluated directly as:
\begin{align}
    (\nabla \times \vec{E}^h)_z(\mathbf{x}) & = \sum_{a=1}^6 (\nabla \times \vec{N}_a^I)_z e_a = \sum_{a=1}^6 c_{\mathrm{curl}, a}^I e_a = (\mathbf{c}_{\mathrm{curl}}^I)^T \mathbf{e}_{\mathrm{stencil}} \in \mathbb{R}, \label{eq:field_curl_constant} \\
    \nabla \cdot \vec{E}^h(\mathbf{x})      & = \sum_{a=1}^6 (\nabla \cdot \vec{N}_a^I) e_a = \sum_{a=1}^6 c_{\mathrm{div}, a}^I e_a = (\mathbf{c}_{\mathrm{div}}^I)^T \mathbf{e}_{\mathrm{stencil}} \in \mathbb{R}. \label{eq:field_div_constant}
\end{align}
Equations~\eqref{eq:field_curl_constant}--\eqref{eq:field_div_constant} demonstrate that within each local subdomain $\Omega_s^I$, both the curl and divergence of the trial vector field are spatially constant, depending solely on the precomputed coefficient matrix $\boldsymbol{\beta}^I$ and the nodal degree-of-freedom vector $\mathbf{e}_{\mathrm{stencil}}$.

\paragraph{VNMM as a Directional Vector Point Interpolation Method (Vector PIM):}
Fundamentally, the kinematic construction of the VNMM represents a directional vector extension of the Point Interpolation Method (PIM) originally established by Liu and Gu for scalar continuum mechanics \cite{liu2001point, liu2005introduction}. In classical scalar PIM, shape functions satisfy the standard scalar Kronecker-delta property ($\phi_i(\mathbf{x}_j) = \delta_{ij}$). In the electromagnetic vector context, the VNMM constitutes a Vector PIM wherein the vector shape functions $\vec{N}_a^I(\mathbf{x})$ satisfy the directional projection Kronecker-delta property:
\begin{equation}
    \vec{N}_a^I(\mathbf{x}_k) \cdot \vec{t}_k = \delta_{ak}, \qquad a, k \in \{1, \dots, 6\}.
    \label{eq:vector_pim_kronecker}
\end{equation}
Unlike conventional scalar and vector meshless approximations based on Moving Least Squares (MLS) or radial basis functions (RBF)---which lack the Kronecker-delta property and require complex auxiliary mechanisms such as Lagrange multipliers, penalty functions, or boundary element coupling to enforce essential constraints \cite{atluri2002meshless}---the directional Vector PIM constructs true vector interpolants over scattered nodes by directly inverting the polynomial moment placement matrix $\mathbf{A}_I$. Because nodal degrees of freedom represent physical directional projections ($e_k = \vec{E}(\mathbf{x}_k) \cdot \vec{t}_k$), tangential Dirichlet and PEC boundary conditions can be satisfied directly at the discrete nodal level by simple collocation, providing a cell-free kinematic foundation for the MLPG framework.

\subsection{Scale-Invariant Stencil Selection Algorithm}
\label{subsec:stencil_selection}

On unstructured, non-uniform, or perturbed point clouds, the local moment matrix $\mathbf{A}_I \in \mathbb{R}^{6 \times 6}$ defined in \eqref{eq:matrix_A_row} can suffer from geometric or directional collinearity singularities ($\det(\mathbf{A}_I) = 0$), as formalized by the Null-Space Orthogonality Principle in \cite{stable_vnmm2d_2026}. To ensure that every node $P_I$ is assigned an algebraically well-conditioned trial support stencil $\Omega_u^I$, an automated selection algorithm is required.

\subsubsection{Dimensionless Matrix Reformulation and Quartic Invariance}
Because columns 1--2 of $\mathbf{A}_I$ are dimensionless ($\mathcal{O}(1)$) while columns 3--6 scale with physical distances ($\mathcal{O}(h_I)$), the determinant exhibits a physical quartic scaling: $\det(\mathbf{A}_I) \sim \mathcal{O}(h_I^4)$. Consequently, applying a fixed dimensional determinant threshold causes artificial threshold decay under grid refinement ($h_I \to 0$) \cite{stable_vnmm2d_2026}.
To decouple qualification from physical domain dimensions, local coordinates are normalized by the local nodal spacing $h_I \equiv h_{\mathrm{local}}(\mathbf{x}_I)$:
\begin{equation}
    \hat{x}_k = \frac{x_k - x_I}{h_I}, \qquad \hat{y}_k = \frac{y_k - y_I}{h_I}.
\end{equation}
The dimensional matrix factors exactly as $\mathbf{A}_I = \widehat{\mathbf{A}}_I \mathbf{D}_h$, with diagonal scaling matrix $\mathbf{D}_h = \operatorname{diag}(1, 1, h_I, h_I, h_I, h_I)$, establishing the determinant relationship:
\begin{equation}
    \det(\mathbf{A}_I) = h_I^4 \det(\widehat{\mathbf{A}}_I) \implies |\det(\widehat{\mathbf{A}}_I)| = \frac{|\det(\mathbf{A}_I)|}{h_I^4}.
    \label{eq:dimensionless_det_relation}
\end{equation}
The scale-invariant admissibility criterion is formulated directly on the dimensionless matrix $\widehat{\mathbf{A}}_I$:
\begin{equation}
    |\det(\widehat{\mathbf{A}}_I)| \ge \widehat{\mathrm{Tol}} > 0.
    \label{eq:scale_invariant_det_criterion}
\end{equation}
Following the calibrated baseline in \cite{stable_vnmm2d_2026}, we adopt $\widehat{\mathrm{Tol}} = 0.80$, which bounds the stencil condition number to $\kappa_2(\widehat{\mathbf{A}}) \le 5.9$ while maintaining compact nearest-neighbor supports ($K_{\mathrm{avg}} \approx 6.5\text{--}7.0$) without support dilation as $h \to 0$. In our numerical implementation on perturbed clouds, each local integration radius $r_s^I = \alpha_s h_I$ adapts to the individual nearest-neighbor distance $h_I$, whereas the dimensional determinant threshold $|\det(\mathbf{A}_I)| \ge \widehat{\mathrm{Tol}} \, \overline{h}^4$ employs the cloud-averaged nearest-neighbor spacing $\overline{h} = \frac{1}{N}\sum_{I=1}^N h_I$ so that localized node pairing under stochastic jitter does not artificially relax the sextet conditioning threshold.

\subsubsection{The 4-Phase Automated Selection Pipeline}
Following \cite{stable_vnmm2d_2026}, for each node $P_I$ ($I = 1, \dots, N$), the active stencil $\Omega_u^I$ is determined through the following 4-phase procedure:
\begin{enumerate}
    \item \textbf{Phase 1 (Spatial Neighborhood Query via $k$-d Tree):} An orthogonal $k$-d tree constructed over all nodal coordinates $\mathcal{X}$ queries the $K_{\mathrm{initial}} \ge 12$ nearest spatial neighbors of $\mathbf{x}_I$ with $\mathcal{O}(\log N)$ complexity. The local spacing $h_I \equiv h_{\mathrm{local}}(\mathbf{x}_I)$ is evaluated deterministically from this local cluster (as the nearest-neighbor distance or canonical 6-neighbor mean).
    \item \textbf{Phase 2 (Central Node Anchoring):} The evaluation node $\mathbf{x}_I$ is permanently anchored as the first member of the candidate stencil ($\mathbf{x}_{(1)} = \mathbf{x}_I$, so $\Delta\mathbf{x}_{(1)} = \mathbf{0}$). This anchors the constant polynomial terms directly to $\vec{E}(\mathbf{x}_I)$. The remaining $K-1$ candidate nodes are sorted in ascending order of Euclidean distance $\|\mathbf{x}_k - \mathbf{x}_I\|$.
    \item \textbf{Phase 3 (Proximity-Ordered Evaluation with Early Stopping):} Candidate 5-combinations are drawn from the sorted neighbor list in order of increasing mean distance. For each candidate sextet, normalized local coordinates are assembled into $\widehat{\mathbf{A}}_I$. The first combination satisfying the dimensionless threshold:
          \begin{equation}
              |\det(\widehat{\mathbf{A}}_I)| \ge \widehat{\mathrm{Tol}}
          \end{equation}
          is immediately accepted as the active trial support $\Omega_u^I$, terminating the search loop.
    \item \textbf{Phase 4 (Adaptive Pool Expansion):} If no combination among the current pool of $K$ neighbors satisfies the qualification threshold, the pool size is incrementally expanded ($K \leftarrow K + \Delta K$, with $\Delta K = 4$) and Phase~3 is repeated. In practice, early stopping accepts an admissible stencil within the first few evaluations ($\bar{N}_{\mathrm{tested}} \le 6$ on regular and moderately perturbed clouds), keeping the computational cost minimal.
\end{enumerate}

\paragraph{Characteristic Length Scales:}
In this cell-free formulation, spatial scales enter exclusively through two quantities:
the local nodal spacing $h_I \equiv h_{\mathrm{local}}(\mathbf{x}_I)$ (evaluated from nearest neighbors to define the local integration disk radius $r_s^I = \alpha_s h_I$ and normalize local coordinates in $\widehat{\mathbf{A}}_I$), and the nominal global discretization spacing $h_{\mathrm{nom}} = L/(N_x - 1)$ used to quantify asymptotic error decay rates $\mathcal{O}(h_{\mathrm{nom}}^p)$ under grid refinement.

\paragraph{Efficiency in Cell-Free MLPG vs. Cell-Based VNMM-Gauss:}
A major computational advantage of the MLPG formulation over the cell-based VNMM-Gauss in \cite{stable_vnmm2d_2026} is that support stencil selection is executed \textbf{strictly once per node} ($N$ total executions), rather than at every numerical quadrature point across a background cell grid ($N_g \approx 10 N_{\mathrm{DoF}}$). This eliminates tens of thousands of redundant $k$-d tree queries and combinatorial matrix evaluations (a nearly $9\times$ reduction in stencil evaluations, e.g., from $20{,}736$ in Gauss to $2{,}401$ in MLPG on $49\times 49$), keeping the total stencil preprocessing overhead below a fraction of a second across all benchmark configurations.

\subsection{Local Weak Form and Overlapping Subdomain Energy Formulation}
\label{subsec:local_weak_form}

To establish a rigorous variational foundation without background quadrature meshes, we consider the local subdomains $\Omega_s^I \cap \Omega$ ($I = 1, \dots, N$) forming an overlapping cover of the physical domain, $\bigcup_{I=1}^N (\Omega_s^I \cap \Omega) \supseteq \Omega$. Taking the inner product of the regularized vector Helmholtz equation \eqref{eq:helmholtz_pde} with a local vector test function $\vec{W}_I(\mathbf{x})$ over $\Omega_s^I \cap \Omega$ and applying Green's first vector identity (Stokes' circulation and Gauss' divergence theorems) yields the exact local weak identity:
\begin{equation}
    \begin{aligned}
        \iint_{\Omega_s^I \cap \Omega} \bigg[ \frac{1}{\mu_r} (\nabla \times \vec{W}_I)_z (\nabla \times \vec{E})_z & + \frac{s_{\mathrm{div}}}{\mu_r} (\nabla \cdot \vec{W}_I) (\nabla \cdot \vec{E}) - k_0^2 \epsilon_r \vec{W}_I \cdot \vec{E} \bigg] d\Omega \\
                                                                                                                    & = \iint_{\Omega_s^I \cap \Omega} \vec{W}_I \cdot \vec{J}_s \, d\Omega + \int_{L_{sI}} \vec{W}_I \cdot \vec{\tau}(\vec{E}) \, ds \\
                                                                                                                    & \quad + \int_{\Gamma_s^I} \vec{W}_I \cdot \vec{\tau}_{\mathrm{ext}} \, ds,
    \end{aligned}
    \label{eq:local_weak_form}
\end{equation}
where $\partial(\Omega_s^I \cap \Omega) = L_{sI} \cup \Gamma_s^I$ is partitioned into the internal circular arc $L_{sI} = \partial\Omega_s^I \cap \Omega$ and the physical boundary intersection segment $\Gamma_s^I = \Omega_s^I \cap \partial\Omega$ carrying prescribed natural boundary fluxes $\vec{\tau}_{\mathrm{ext}}$ (e.g., surface magnetic current densities on PMC walls or Robin radiation impedances, as detailed in Section~\ref{subsec:boundary_treatment}), and $\vec{\tau}(\vec{E}) = \mu_r^{-1}(\vec{n} \times (\nabla \times \vec{E})) + s_{\mathrm{div}}\mu_r^{-1}(\nabla \cdot \vec{E})\vec{n}$ is the natural electromagnetic flux along the outward normal $\vec{n}$ of $\partial(\Omega_s^I \cap \Omega)$.

In classical Petrov--Galerkin formulations (MLPG1 \cite{atluri1998new, atluri2002meshless}), the test function $\vec{W}_I$ is chosen to vanish along the internal arc $L_{sI}$ so as to eliminate the internal line integral $\int_{L_{sI}} \vec{W}_I \cdot \vec{\tau}(\vec{E}) \, ds$. However, as shown next in Section~\ref{subsec:mlpg6_selection}, when a fixed polynomial stencil ($\mathcal{P}^1$) is employed on $\Omega_s^I$, vanishing test functions cause the entire stiffness bilinear form to collapse identically to zero. Conversely, when non-vanishing polynomial test functions ($\vec{W}_a^I = \vec{N}_a^I$, MLPG6/LPIM \cite{liu2001local, liu2005introduction, nicomedes2012meshless}) are assembled in a classical row-by-row manner on overlapping disks $\Omega_s^I \ne \Omega_s^J$, the global matrix is unsymmetric; while symmetric overlapping-sphere approaches such as the Method of Finite Spheres \cite{de2000method, nicomedes2021acoustic} require complex numerical cubature over lens-shaped intersections $\Omega_s^I \cap \Omega_s^J$.

To resolve this variationally without evaluating internal curved line integrals over $L_{sI}$ or lens-intersection quadratures, we propose an \textbf{Overlapping Subdomain Ritz--Galerkin Energy Formulation}. Specifically, the global electromagnetic action functional is discretized as the additive sum of local subdomain energy functionals over the cover $\{\Omega_s^I \cap \Omega\}_{I=1}^N$:
\begin{equation}
    \Pi_h(\vec{E}^h) = \sum_{I=1}^N \left\{ \frac{1}{2} a_{\Omega_s^I}(\vec{E}_I^h, \vec{E}_I^h) - \frac{1}{2} k_0^2 m_{\Omega_s^I}(\vec{E}_I^h, \vec{E}_I^h) - \ell_{\Omega_s^I}(\vec{E}_I^h) \right\},
\end{equation}
where the local volume bilinear forms and consistent linear functional on each subdomain $\Omega_s^I \cap \Omega$ are defined directly in symmetric volume energy form without integration by parts across $L_{sI}$:
\begin{equation}
    \begin{aligned}
        a_{\Omega_s^I}(\vec{W}_I, \vec{E}) - k_0^2 m_{\Omega_s^I}(\vec{W}_I, \vec{E}) & \equiv \iint_{\Omega_s^I \cap \Omega} \bigg[ \frac{1}{\mu_r} (\nabla \times \vec{W}_I)_z (\nabla \times \vec{E})_z \\
                                                                                      & \qquad + \frac{s_{\mathrm{div}}}{\mu_r} (\nabla \cdot \vec{W}_I) (\nabla \cdot \vec{E}) - k_0^2 \epsilon_r \vec{W}_I \cdot \vec{E} \bigg] d\Omega \\
                                                                                      & = \ell_{\Omega_s^I}(\vec{W}_I).
    \end{aligned}
    \label{eq:volume_energy_form}
\end{equation}
For homogeneous eigenvalue problems in conservative cavities ($\ell_{\Omega_s^I} \equiv 0$), taking the first variation $\delta \Pi_h = 0$ with respect to the global nodal degrees of freedom directly yields a self-adjoint generalized eigenvalue problem governed solely by the symmetric volume integrals $a_{\Omega_s^I}$ and $m_{\Omega_s^I}$. Crucially, because both the global stiffness matrix $\mathbf{K}$ and the global mass matrix $\mathbf{M}$ are assembled over the exact same family of overlapping active subdomains $\{\Omega_s^I \cap \Omega\}_{I=1}^N$, the pointwise overlap multiplicity $M(\mathbf{x}) = \sum_{I=1}^N \chi_{\Omega_s^I \cap \Omega}(\mathbf{x})$ (where $\chi$ denotes the characteristic indicator function, with $M \approx 2\pi \alpha_s^2 \approx 4.02$ in the interior) scales both the numerator and denominator of the discrete Rayleigh quotient congruently:
\begin{equation}
    \lambda_h = \frac{\mathbf{e}^T \mathbf{K} \mathbf{e}}{\mathbf{e}^T \mathbf{M} \mathbf{e}} = \frac{\sum_{I=1}^N a_{\Omega_s^I}(\vec{E}_I^h, \vec{E}_I^h)}{\sum_{I=1}^N m_{\Omega_s^I}(\vec{E}_I^h, \vec{E}_I^h)},
    \label{eq:rayleigh_quotient_overlap}
\end{equation}
thereby cancelling out globally without requiring partition-of-unity (PU) weighting functions $\phi_I(\mathbf{x}) = w_I(\mathbf{x}) / \sum_J w_J(\mathbf{x})$ (which are non-polynomial rational functions that destroy polynomial reproduction and necessitate numerical cubatures over lens-shaped intersections $\Omega_s^I \cap \Omega_s^J$ \cite{de2000method, nicomedes2021acoustic}).

For inhomogeneous source and boundary-value problems (such as the Patch Test and Method of Manufactured Solutions), the proposed additive formulation likewise preserves strict variational consistency through a \textbf{subdomain-by-subdomain variational balance}. By Green's vector identity \eqref{eq:local_weak_form}, the continuous governing operator satisfies:
\begin{equation}
    \begin{aligned}
        a_{\Omega_s^I}(\vec{W}_I, \vec{E}^*) - k_0^2 m_{\Omega_s^I}(\vec{W}_I, \vec{E}^*) & = \iint_{\Omega_s^I \cap \Omega} \vec{W}_I \cdot \vec{J}_s \, d\Omega + \int_{\partial(\Omega_s^I \cap \Omega)} \vec{W}_I \cdot \vec{\tau}(\vec{E}^*) \, ds \\
                                                                                      & \equiv \ell_{\Omega_s^I}(\vec{W}_I)
    \end{aligned}
    \label{eq:local_variational_balance}
\end{equation}
identically on each individual subdomain $\Omega_s^I \cap \Omega$. When discretized with trial shape functions $\vec{W}_a^I = \vec{N}_a^I$ ($a=1,\dots,6$), both the local stiffness/mass operators ($\mathbf{K}_I^{\mathrm{local}}, \mathbf{M}_I^{\mathrm{local}}$) and the local load vector $\mathbf{f}_I^{\mathrm{local}} = [\ell_{\Omega_s^I}(\vec{N}_1^I), \dots, \ell_{\Omega_s^I}(\vec{N}_6^I)]^T$ are integrated over the \emph{exact same} active geometric domain $\Omega_s^I \cap \Omega$. Consequently, the local subdomain residual vector for the exact nodal coefficients $\mathbf{e}^*$ vanishes:
\begin{equation}
    \mathbf{r}_I(\mathbf{e}^*) \equiv (\mathbf{K}_I^{\mathrm{local}} - k_0^2 \mathbf{M}_I^{\mathrm{local}}) \mathbf{L}_I \mathbf{e}^* - \mathbf{f}_I^{\mathrm{local}} = \mathbf{0} \quad \forall I = 1, \dots, N,
    \label{eq:subdomain_residual_zero}
\end{equation}
holding to machine precision ($\approx 10^{-15}$) for affine fields in the Patch Test, and bounded by the optimal local truncation error $\mathcal{O}(h_I^2 S_0^I)$ for general smooth fields in MMS. Upon assembling the global discrete system $\mathbf{K}\mathbf{e} - k_0^2 \mathbf{M}\mathbf{e} = \mathbf{F}$ via Boolean scatter matrices $\mathbf{L}_I^T$:
\begin{equation}
    \begin{aligned}
        \mathbf{R}(\mathbf{e}^*) & \equiv (\mathbf{K} - k_0^2 \mathbf{M})\mathbf{e}^* - \mathbf{F} \\
                                 & = \sum_{I=1}^N \mathbf{L}_I^T \underbrace{\left[ (\mathbf{K}_I^{\mathrm{local}} - k_0^2 \mathbf{M}_I^{\mathrm{local}}) \mathbf{L}_I \mathbf{e}^* - \mathbf{f}_I^{\mathrm{local}} \right]}_{\mathbf{r}_I(\mathbf{e}^*) \equiv \mathbf{0}} = \sum_{I=1}^N \mathbf{L}_I^T (\mathbf{0}) \equiv \mathbf{0}.
    \end{aligned}
    \label{eq:global_residual_sum}
\end{equation}
Because the local equilibrium \eqref{eq:subdomain_residual_zero} is satisfied on each disk independently, the global residual \eqref{eq:global_residual_sum} vanishes identically regardless of the pointwise overlap multiplicity $M(\mathbf{x})$. In verification benchmarks where $\vec{E}^*$ is prescribed analytically (Sections~\ref{subsec:stage2_patch_test} and \ref{subsec:stage3_mms}), $\ell_{\Omega_s^I}$ is evaluated directly via the consistent identity \eqref{eq:local_variational_balance}; whereas in general source-driven problems where $\vec{E}^*$ is unknown a priori, the overlapping disks act as a partition-free local integration cover in which opposing interior boundary fluxes along $L_{sI}$ cancel asymptotically in the domain interior as $h_{\mathrm{nom}} \to 0$ while physical boundary fluxes $\vec{\tau}_{\mathrm{ext}}$ are enforced along $\Gamma_s^I = \Omega_s^I \cap \partial\Omega$.

\subsection{Selection of the MLPG6 (Local Galerkin) Test Space}
\label{subsec:mlpg6_selection}

A central question in local meshless methods concerns the choice of the vector test space $\vec{W}_I$. In MLS-based MLPG1 \cite{atluri1998new, atluri2002meshless} and radial-basis LRPIM \cite{liu2005introduction}, localized radial weight functions that vanish continuously on the subdomain boundary ($w_I(r_s^I) = 0$) are routinely employed because the MLS/RBF trial derivatives vary spatially within $\Omega_s^I$. In the present vector PIM setting with a fixed linear polynomial stencil ($\mathcal{P}^1$) per subdomain \cite{liu2001local, stable_vnmm2d_2026}, however, using vanishing test functions ($\vec{W}_I|_{\partial\Omega_s^I} = \mathbf{0}$) causes the stiffness integral to vanish identically (analogous to the vanishing of constant-gradient domain integrals for compactly supported weights in scalar LPIM \cite{liu2001local}). Because the curl and divergence of linear vector trial fields ($\mathcal{P}^1$) are spatially constant within $\Omega_s^I$ (Eqs.~\eqref{eq:field_curl_constant}--\eqref{eq:field_div_constant}), Stokes' circulation theorem and Gauss' divergence theorem convert both volume stiffness integrals directly into closed contour integrals of the test function along $\partial\Omega_s^I$. Since $\vec{W}_I$ vanishes on the boundary, both integrals vanish identically, causing the entire stiffness matrix to collapse to zero ($\mathbf{K} \equiv \mathbf{0}$). Consequently, weight-based MLPG1 testing is unsuitable for fixed-stencil first-order vector approximations.

To resolve this limitation and preserve the mathematical structure of the continuous Maxwell operators, we adopt the \textbf{MLPG6 (Local Bubnov--Galerkin)} framework \cite{atluri2002meshless}, choosing the vector test functions directly from the space of \textbf{trial shape functions themselves} within each local subdomain $\Omega_s^I$:
\begin{equation}
    \vec{W}_a^I(\mathbf{x}) = \vec{N}_a^I(\mathbf{x}) \quad \forall \mathbf{x} \in \Omega_s^I, \quad a = 1, \dots, 6.
    \label{eq:mlpg6_choice}
\end{equation}
This Local Galerkin choice confers three key advantages:
\begin{enumerate}
    \item \textbf{Exact Symmetry and Positive-Definiteness:} Testing with the trial basis ($\vec{W}_a^I = \vec{N}_a^I$) ensures that the local stiffness matrix $\mathbf{K}_I^{\mathrm{local}}$ and mass matrix $\mathbf{M}_I^{\mathrm{local}}$ are formed by inner products of the basis functions with themselves. This guarantees that every local matrix is strictly symmetric and positive-(semi)definite, naturally preserving the self-adjointness of the vector curl-curl operator and preventing unphysical complex eigenvalues.
    \item \textbf{Non-Trivial Stiffness Transmission:} Because $(\nabla \times \vec{N}_a^I)_z$ and $\nabla \cdot \vec{N}_a^I$ are non-zero constants, their inner products against themselves are strictly positive, completely resolving the null-stiffness collapse of MLPG1.
    \item \textbf{Exact Closed-Form Integration:} Integrands reduce to low-order Cartesian polynomials, enabling exact closed-form analytical integration over both circular and truncated subdomains without numerical quadrature errors (derived in Section~\ref{subsec:analytic_formulas} and \ref{app:truncated_geometry}).
\end{enumerate}

\subsection{Closed-Form Analytical Discretization Formulas}
\label{subsec:analytic_formulas}

Let $\Omega_s^I$ be a circular disk of radius $r_s^I$. Its area and polar second moment of area are:
\begin{equation}
    \begin{aligned}
        S_0 & = \iint_{\Omega_s^I} d\Omega = \pi (r_s^I)^2,                                                                \\
        S_2 & = \iint_{\Omega_s^I} (x - x_I)^2 d\Omega = \iint_{\Omega_s^I} (y - y_I)^2 d\Omega = \frac{\pi (r_s^I)^4}{4}.
    \end{aligned}
    \label{eq:disk_moments}
\end{equation}
All odd moments vanish identically by radial symmetry:
\begin{equation}
    \iint_{\Omega_s^I} (x - x_I) \, d\Omega = \iint_{\Omega_s^I} (y - y_I) \, d\Omega = \iint_{\Omega_s^I} (x - x_I)(y - y_I) \, d\Omega \equiv 0.
\end{equation}

Within each subdomain $\Omega_s^I$, the complete local trial basis consists of the 6 vector shape functions $\vec{N}_a^I(\mathbf{x})$ ($a=1,\dots,6$). Substituting \eqref{eq:shape_functions} into the local weak form \eqref{eq:local_weak_form}, the local stiffness matrix $\mathbf{K}_I^{\mathrm{local}} \in \mathbb{R}^{6 \times 6}$ is partitioned into curl and divergence contributions:
\begin{equation}
    \mathbf{K}_I^{\mathrm{local}} = \mathbf{K}_{I, \mathrm{curl}}^{\mathrm{local}} + s_{\mathrm{div}} \mathbf{K}_{I, \mathrm{div}}^{\mathrm{local}},
\end{equation}
where, using the constant vectors $\mathbf{c}_{\mathrm{curl}}^I, \mathbf{c}_{\mathrm{div}}^I \in \mathbb{R}^6$ defined in \eqref{eq:curl_constant}--\eqref{eq:div_constant}:
\begin{align}
    \mathbf{K}_{I, \mathrm{curl}}^{\mathrm{local}} & = \frac{1}{\mu_r} \left( \mathbf{c}_{\mathrm{curl}}^I (\mathbf{c}_{\mathrm{curl}}^I)^T \right) \pi (r_s^I)^2, \label{eq:k_curl_analytic} \\
    \mathbf{K}_{I, \mathrm{div}}^{\mathrm{local}}  & = \frac{1}{\mu_r} \left( \mathbf{c}_{\mathrm{div}}^I (\mathbf{c}_{\mathrm{div}}^I)^T \right) \pi (r_s^I)^2. \label{eq:k_div_analytic}
\end{align}
Both matrices are outer products of vectors with themselves, hence manifestly symmetric and positive-semidefinite.

For the local mass matrix $\mathbf{M}_I^{\mathrm{local}} \in \mathbb{R}^{6 \times 6}$:
\begin{equation}
    M_{I, ab}^{\mathrm{local}} = \iint_{\Omega_s^I} \epsilon_r (\vec{N}_a^I \cdot \vec{N}_b^I) \, d\Omega.
\end{equation}
Expanding the dot product using the polynomial basis representation \eqref{eq:P1_shape_function_expansion_mlpg}:
\begin{equation}
    \begin{aligned}
        \vec{N}_a^I(\mathbf{x}) \cdot \vec{N}_b^I(\mathbf{x}) = & (\boldsymbol{\beta}_{1:2, a}^I)^T \boldsymbol{\beta}_{1:2, b}^I                                                                \\
                                                                & + (\boldsymbol{\beta}_{3:6, a}^I)^T \boldsymbol{\beta}_{3:6, b}^I \left[ (x - x_I)^2 + (y - y_I)^2 \right] + \text{odd terms},
    \end{aligned}
\end{equation}
where $\boldsymbol{\beta}_{1:2, a}^I = [\beta_{1a}^I, \beta_{2a}^I]^T$ and $\boldsymbol{\beta}_{3:6, a}^I = [\beta_{3a}^I, \beta_{4a}^I, \beta_{5a}^I, \beta_{6a}^I]^T$.
Integrating over the disk yields the closed-form mass matrix:
\begin{equation}
    \mathbf{M}_I^{\mathrm{local}} = \epsilon_r \left( \mathbf{C}_0^I \pi (r_s^I)^2 + \mathbf{C}_1^I \frac{\pi (r_s^I)^4}{4} \right),
    \label{eq:m_local_analytic}
\end{equation}
where $\mathbf{C}_0^I = (\boldsymbol{\beta}_{1:2}^I)^T \boldsymbol{\beta}_{1:2}^I \in \mathbb{R}^{6 \times 6}$ and $\mathbf{C}_1^I = (\boldsymbol{\beta}_{3:6}^I)^T \boldsymbol{\beta}_{3:6}^I \in \mathbb{R}^{6 \times 6}$. Both $\mathbf{C}_0^I$ and $\mathbf{C}_1^I$ are Gram matrices of real submatrices, guaranteeing that $\mathbf{M}_I^{\mathrm{local}}$ is strictly symmetric and positive-definite.

\subsubsection{Closed-Form Discretization for Truncated Subdomains (Cases 2 and 3)}
\label{subsubsec:truncated_subdomains_analytic}

When a local subdomain $\Omega_s^I$ crosses the physical domain boundary $\partial\Omega$---occurring either for a boundary node $\mathbf{x}_B \in \partial\Omega$ (Case~3) or for an interior node $\mathbf{x}_{I_2}$ located close to $\partial\Omega$ ($\mathrm{dist}(\mathbf{x}_{I_2}, \partial\Omega) < r_s^{I_2}$, Case~2), as illustrated in Fig.~\ref{fig:schematic_mlpg}---the physically admissible integration domain is strictly truncated to the intersection:
\begin{equation}
    \Omega_s^{I, \mathrm{active}} \equiv \Omega_s^I \cap \Omega.
    \label{eq:truncated_subdomain_def}
\end{equation}
Because the complete polynomial space $\mathcal{P}^1$ is employed, weak-form integrals over truncated geometries can be evaluated in closed form without numerical domain triangulation or background cells.

Recall that for the trial stencil of node $P_I$, both the curl $(\nabla \times \vec{N}_a^I)_z = c_{\mathrm{curl}, a}^I$ and the divergence $\nabla \cdot \vec{N}_a^I = c_{\mathrm{div}, a}^I$ are strictly spatial constants on $\Omega_s^I$ (Eqs.~\eqref{eq:field_curl_constant}--\eqref{eq:field_div_constant}). Consequently, the entire stiffness integrand factors outside the domain integral:
\begin{align}
    \mathbf{K}_{I, \mathrm{curl}}^{\mathrm{trunc}} & = \frac{1}{\mu_r} \left( \mathbf{c}_{\mathrm{curl}}^I (\mathbf{c}_{\mathrm{curl}}^I)^T \right) \iint_{\Omega_s^I \cap \Omega} d\Omega = \frac{1}{\mu_r} \left( \mathbf{c}_{\mathrm{curl}}^I (\mathbf{c}_{\mathrm{curl}}^I)^T \right) A_{\mathrm{active}}^I, \label{eq:k_curl_truncated}                          \\
    \mathbf{K}_{I, \mathrm{div}}^{\mathrm{trunc}}  & = \frac{s_{\mathrm{div}}}{\mu_r} \left( \mathbf{c}_{\mathrm{div}}^I (\mathbf{c}_{\mathrm{div}}^I)^T \right) \iint_{\Omega_s^I \cap \Omega} d\Omega = \frac{s_{\mathrm{div}}}{\mu_r} \left( \mathbf{c}_{\mathrm{div}}^I (\mathbf{c}_{\mathrm{div}}^I)^T \right) A_{\mathrm{active}}^I, \label{eq:k_div_truncated}
\end{align}
where $A_{\mathrm{active}}^I = \iint_{\Omega_s^I \cap \Omega} d\Omega$ denotes the active physical area. Introducing the dimensionless active area fraction:
\begin{equation}
    \eta_I \equiv \frac{A_{\mathrm{active}}^I}{\pi (r_s^I)^2} = \frac{\mathrm{Area}(\Omega_s^I \cap \Omega)}{\pi (r_s^I)^2} \in (0, 1],
    \label{eq:area_fraction_eta}
\end{equation}
\textbf{the exact stiffness matrix for any truncated subdomain is obtained by simple geometric scaling of the full-disk analytical stiffness matrix}:
\begin{equation}
    \mathbf{K}_I^{\mathrm{trunc}} = \eta_I \mathbf{K}_I^{\mathrm{full}}.
    \label{eq:k_truncated_scaled}
\end{equation}
In meshless literature and computational practice, three distinct strategies can be pursued to handle such truncated integration domains:
\begin{enumerate}
    \item \textbf{Adaptive Radius Truncation for Near-Boundary Nodes:}
          In classical scalar and vector MLPG formulations (e.g., Atluri and Shen \cite{atluri2002meshless}), a widely used strategy for interior nodes close to $\partial\Omega$ is to locally shrink the subdomain radius $r_s^I$ such that the disk is strictly tangent to the boundary:
          \begin{equation}
              r_s^I = \min\left(r_{s, 0}^I, \, \mathrm{dist}(\mathbf{x}_I, \partial\Omega)\right),
              \label{eq:radius_shrink_tangent}
          \end{equation}
          where $r_{s, 0}^I$ is the nominal subdomain radius. While this ensures $\Omega_s^I \subset \bar{\Omega}$ and eliminates Case~2 truncations altogether for interior nodes, it creates non-uniform, node-dependent integration sizes, can lead to vanishingly small subdomains for nodes arbitrarily close to $\partial\Omega$, and still leaves boundary nodes (Case~3) truncated.

    \item \textbf{Numerical Quadrature and Coordinate Mapping over Truncated Intersections:}
          A second alternative in the literature is to evaluate the local integrals over $\Omega_s^I \cap \Omega$ numerically. While crude filtering of polar Gauss points through boundary bounding boxes exhibits slow $\mathcal{O}(h_{\mathrm{quad}})$ convergence across chord discontinuities, Atluri \cite{atluri2004meshless} demonstrated that this limitation can be resolved by mapping the truncated circular domains onto a regular reference domain, such as a unit circle or standard parent domain, via coordinate transformations. Standard Gauss quadrature can then be applied smoothly on the mapped domain without boundary clipping errors or background cells, although evaluating the coordinate transformation and its associated Jacobian determinant for various chord angles and corner intersections introduces non-trivial algorithmic complexity.

    \item \textbf{Exact Closed-Form Active Area Scaling and Immunity to Small-Cut Ill-Conditioning (Present Formulation):}
          An important mathematical consequence of employing the complete polynomial space $\mathcal{P}^1$ in VNMM is that both numerical integration and coordinate mappings are bypassed. As established in Eqs.~\eqref{eq:k_curl_truncated}--\eqref{eq:k_truncated_scaled}, the stiffness integrand factors outside the domain integral, allowing the exact stiffness matrix on any truncated domain to be computed by directly scaling the full-disk analytical matrix by the active area fraction $\eta_I$.
          For piecewise planar boundaries (such as polygonal and rectangular cavities), the active area $A_{\mathrm{active}}^I$ and fraction $\eta_I$ for all boundary configurations, including flat walls ($\eta = 1/2$), corner vertices ($\eta = 1/4$), boundary nodes near corners ($\eta \in [1/4, 1/2)$), and interior nodes intersecting one or two walls, are evaluated in exact closed analytical form as detailed in \ref{app:truncated_geometry}.
          Furthermore, unlike unfitted finite element methods (such as CutFEM \cite{hansbo2002unfitted}) or cell-based meshless schemes where background cells intersected by $\partial\Omega$ can have arbitrarily small volume fractions ($\eta \to 0$, causing severe matrix ill-conditioning and requiring ghost-penalty stabilization), every integration disk $\Omega_s^I$ in VNMM-MLPG is centered directly at a physical node $\mathbf{x}_I \in \overline{\Omega}$. Consequently, on any convex polygonal domain with interior corner angles $\omega$, the active area fraction is strictly bounded below by $\eta_I \ge \omega/(2\pi)$ ($\eta_I \ge 1/2$ along flat walls and $\eta_I \ge 1/4$ at $90^\circ$ corners), rendering pathological small-cut ill-conditioning ($\eta_I \to 0$) geometrically impossible.
\end{enumerate}

For the local mass matrix $\mathbf{M}_I^{\mathrm{local}}$, substituting the complete $\mathcal{P}^1$ shape function expansion yields the exact moment-based integration over any truncated subdomain:
\begin{equation}
    \mathbf{M}_I^{\mathrm{trunc}} = \epsilon_r \left( \mathbf{P}_0^I S_0^I + \mathbf{P}_x^I S_{1x}^I + \mathbf{P}_y^I S_{1y}^I + \mathbf{P}_{xx}^I S_{2xx}^I + \mathbf{P}_{yy}^I S_{2yy}^I + \mathbf{P}_{xy}^I S_{2xy}^I \right),
    \label{eq:m_truncated_moments}
\end{equation}
where $\mathbf{P}_0^I, \dots, \mathbf{P}_{xy}^I \in \mathbb{R}^{6 \times 6}$ are the constant coefficient matrices formed by outer products of the placement matrix rows $\boldsymbol{\beta}^I$ (obtained from the expansion \eqref{eq:P1_shape_function_expansion_mlpg}), and $S_0^I, \dots, S_{2xy}^I$ are the geometric moments of $\Omega_s^I \cap \Omega$ relative to $\mathbf{x}_I$:
\begin{align}
    S_0^I     & = A_{\mathrm{active}}^I = \iint_{\Omega_s^I \cap \Omega} d\Omega,                                                                 \\
    S_{1x}^I  & = \iint_{\Omega_s^I \cap \Omega} (x - x_I) \, d\Omega, \quad S_{1y}^I = \iint_{\Omega_s^I \cap \Omega} (y - y_I) \, d\Omega,      \\
    S_{2xx}^I & = \iint_{\Omega_s^I \cap \Omega} (x - x_I)^2 \, d\Omega, \quad S_{2yy}^I = \iint_{\Omega_s^I \cap \Omega} (y - y_I)^2 \, d\Omega, \\
    S_{2xy}^I & = \iint_{\Omega_s^I \cap \Omega} (x - x_I)(y - y_I) \, d\Omega.
\end{align}
On an uncut full disk, reflection symmetry ensures $S_{1x}^I = S_{1y}^I = S_{2xy}^I = 0$ and $S_{2xx}^I = S_{2yy}^I = \frac{\pi}{4}(r_s^I)^4$, recovering \eqref{eq:m_local_analytic}. On truncated subdomains, normal reflection symmetry is broken, introducing non-vanishing linear moments (e.g., $S_{1y}^I = \frac{2}{3}(r_s^I)^3$ on flat walls and $S_{1x}^I = S_{1y}^I = \frac{1}{3}(r_s^I)^3$ on corner vertices). In our implementation, these moments are evaluated in closed form (as derived in \ref{app:truncated_geometry}), ensuring that all numerical tests and assemblies across full and truncated domains employ the exact mass matrix integration.

\subsubsection{Classical Row-by-Row MLPG Assembly and Matrix Symmetry Analysis}
\label{subsubsec:row_mlpg_symmetry}

In the classical MLPG paradigm originally conceived by Atluri and Zhu \cite{atluri1998new, atluri2002meshless}, each node $P_I$ is associated with a \textit{single} local test function. In our vector Local Galerkin setting, the test function for node $I$ is selected as the specific shape function corresponding to node $I$ itself within its local stencil:
\begin{equation}
    \vec{W}_I(\mathbf{x}) = \vec{N}_{i_0}^I(\mathbf{x}) \quad \forall \mathbf{x} \in \Omega_s^I,
    \label{eq:test_function_row}
\end{equation}
where $i_0 \in \{1, \dots, 6\}$ is the local stencil index satisfying $\mathbf{x}_{i_0} = \mathbf{x}_I$. Notice first that in classical MLPG6/LPIM \cite{atluri2002meshless, liu2001local}, the internal contour integral $-\int_{L_{sI}} \vec{W}_I \cdot \vec{\tau}(\vec{E}^h) \, ds$ is retained on the left-hand side of \eqref{eq:local_weak_form}. However, if $\vec{E}^h$ is evaluated using the single fixed linear stencil $\Omega_u^I$ throughout $\Omega_s^I$, its second-order derivatives vanish identically ($\nabla \times (\mu_r^{-1}\nabla \times \vec{N}_k^I) - s_{\mathrm{div}}\mu_r^{-1}\nabla(\nabla \cdot \vec{N}_k^I) \equiv \mathbf{0}$), so Green's identity implies that the volume stiffness integral and the boundary flux integral over $\partial\Omega_s^I$ cancel out identically. Avoiding this collapse in classical LPIM requires evaluating $\vec{\tau}(\vec{E}^h)$ along $L_{sI}$ via numerical Gauss quadrature that samples shifting neighbor stencils $\Omega_u(\mathbf{x}_g) \ne \Omega_u^I$, which forfeits closed-form integration and reintroduces $k$-d tree searches along curved arcs. Conversely, testing the pure volume energy form \eqref{eq:volume_energy_form} with $\vec{W}_I$ over $\Omega_s^I$ produces the $I$-th algebraic row equation of the global linear system in closed form:
\begin{equation}
    \sum_{k=1}^6 K_{I, k}^{\mathrm{row}} e_k = F_I,
\end{equation}
where the non-zero row coefficients are given analytically by:
\begin{equation}
    K_{I, k}^{\mathrm{row}} = \left[ \frac{1}{\mu_r} c_{\mathrm{curl}, i_0}^I c_{\mathrm{curl}, k}^I + \frac{s_{\mathrm{div}}}{\mu_r} c_{\mathrm{div}, i_0}^I c_{\mathrm{div}, k}^I \right] \pi (r_s^I)^2.
    \label{eq:k_row_analytic}
\end{equation}

\begin{remark}[Inherent Matrix Unsymmetry in Classical MLPG and LPIM]
    A fundamental property of classical MLPG \cite{atluri2002meshless} and the Local Point Interpolation Method (LPIM) \cite{liu2001local, liu2005introduction} that warrants careful mathematical emphasis is that \textbf{the global stiffness matrix $\mathbf{K}$ and mass matrix $\mathbf{M}$ assembled row-by-row are intrinsically unsymmetric} ($K_{Ij} \ne K_{Ji}$).
    This unsymmetry stems from two compounding sources:
    \begin{enumerate}
        \item \textbf{Geometric Disparity of Integration Subdomains ($\Omega_s^I \ne \Omega_s^J$):}
              The matrix element $K_{Ij}$ in row $I$ is evaluated by integrating over the local subdomain $\Omega_s^I$ centered at node $\mathbf{x}_I$:
              \begin{equation}
                  K_{Ij} = \iint_{\Omega_s^I} \left[ \frac{1}{\mu_r} (\nabla \times \vec{W}_I)_z (\nabla \times \vec{N}_j^I)_z + \frac{s_{\mathrm{div}}}{\mu_r} (\nabla \cdot \vec{W}_I) (\nabla \cdot \vec{N}_j^I) \right] d\Omega,
              \end{equation}
              whereas the transposed element $K_{Ji}$ in row $J$ originates from the weak form of node $J$ and is integrated over the distinct subdomain $\Omega_s^J$ centered at node $\mathbf{x}_J$:
              \begin{equation}
                  K_{Ji} = \iint_{\Omega_s^J} \left[ \frac{1}{\mu_r} (\nabla \times \vec{W}_J)_z (\nabla \times \vec{N}_i^J)_z + \frac{s_{\mathrm{div}}}{\mu_r} (\nabla \cdot \vec{W}_J) (\nabla \cdot \vec{N}_i^J) \right] d\Omega.
              \end{equation}
              Because $\Omega_s^I \ne \Omega_s^J$ (having different centers $\mathbf{x}_I \ne \mathbf{x}_J$ and potentially distinct radii $r_s^I \ne r_s^J$), the integration domains differ entirely.
        \item \textbf{Disparity of Support Stencils and Shape Function Expressions ($\Omega_u^I \ne \Omega_u^J$):}
              Furthermore, the support stencil $\Omega_u^I$ selected for node $P_I$ through the scale-invariant qualification algorithm (Section~\ref{subsec:stencil_selection}) does not coincide with the stencil $\Omega_u^J$ selected for node $P_J$ ($\Omega_u^I \ne \Omega_u^J$). Consequently, the local placement matrices differ ($\mathbf{A}_I \ne \mathbf{A}_J$), resulting in distinct polynomial coefficient matrices $\boldsymbol{\beta}^I \ne \boldsymbol{\beta}^J$ in \eqref{eq:beta_matrix_inversion}. Even within the overlapping spatial region $\Omega_s^I \cap \Omega_s^J$, the shape function $\vec{N}_j^I(\mathbf{x})$ defined over stencil $\Omega_u^I$ and the shape function $\vec{N}_i^J(\mathbf{x})$ defined over stencil $\Omega_u^J$ possess completely different polynomial coefficient vectors and analytical expressions.
    \end{enumerate}
    Consequently, $K_{Ij} \ne K_{Ji}$ in general, rendering the global system matrix non-Hermitian \cite{atluri2002meshless, liu2001local}. In solid mechanics literature, this asymmetry has historically been handled either by employing unsymmetric solvers or through ad-hoc algebraic averaging $\frac{1}{2}(\mathbf{K} + \mathbf{K}^T)$, which lacks variational consistency.

    For affine polynomial fields (such as the Galerkin Patch Test in Section~\ref{subsec:stage2_patch_test}), this unsymmetric assembly is completely consistent: because affine fields belong identically to $\mathcal{P}^1$, the local residual vanishes on every disk, and the condensed system $\mathbf{K}_{\mathrm{int}, \mathrm{int}} \mathbf{e}_{\mathrm{int}} = \mathbf{F}_{\mathrm{int}} - \mathbf{K}_{\mathrm{int}, \mathrm{bnd}} \mathbf{e}_{\mathrm{bnd}}$ reproduces the exact field to machine precision. However, for general non-polynomial fields with truncation error (as demonstrated under the MMS in Section~\ref{subsubsec:mms_assembly_comparison}) as well as conservative eigenvalue problems, testing with only one test function per disk fails to maintain discrete inf-sup coercivity and breaks operator self-adjointness, making the symmetric Subdomain Galerkin assembly developed next indispensable.
\end{remark}

\subsubsection{Symmetric Subdomain Galerkin Assembly for Eigenvalue and Inhomogeneous Problems}
\label{subsubsec:symmetric_assembly}

In conservative electromagnetic eigenvalue problems, such as resonant cavities governed by Maxwell's curl-curl operator without physical dissipation, the continuous differential operator is self-adjoint. If an unsymmetric discretization matrix is employed, the loss of discrete self-adjointness can introduce non-physical complex-conjugate eigenvalues ($\mathrm{Im}(\lambda) \ne 0$) and precludes the application of efficient Hermitian eigensolvers \cite{saad2011numerical, boffi2013mixed}. Furthermore, as established in Section~\ref{subsubsec:mms_assembly_comparison}, unsymmetric row-by-row assembly on non-polynomial fields leads to spectral indefiniteness and numerical divergence.

To resolve these challenges while retaining the cell-free nature of local meshless methods (MLPG and LPIM), we introduce the \textbf{Symmetric Subdomain Galerkin} assembly as a primary methodological contribution. Rather than testing each local subdomain against only the single nodal shape function of its evaluation node, we test the local weak form over each independent disk $\Omega_s^I$ against \textit{all six} vector trial shape functions $\vec{W}_a^I = \vec{N}_a^I$ ($a=1,\dots,6$). This generates the full $6 \times 6$ local matrices $\mathbf{K}_I^{\mathrm{local}}$ and $\mathbf{M}_I^{\mathrm{local}}$ defined analytically in closed form in \eqref{eq:k_curl_analytic}--\eqref{eq:m_local_analytic}.

Global assembly is then achieved by accumulating the local subdomain contributions into the global system via standard Boolean scatter operations:
\begin{equation}
    \mathbf{K} = \sum_{I=1}^N \mathbf{L}_I^T \mathbf{K}_I^{\mathrm{local}} \mathbf{L}_I, \qquad \mathbf{M} = \sum_{I=1}^N \mathbf{L}_I^T \mathbf{M}_I^{\mathrm{local}} \mathbf{L}_I, \qquad \mathbf{F} = \sum_{I=1}^N \mathbf{L}_I^T \mathbf{f}_I^{\mathrm{local}},
    \label{eq:global_assembly}
\end{equation}
where $\mathbf{L}_I \in \mathbb{R}^{6 \times N}$ is the Boolean gathering matrix mapping the 6 stencil nodes of $\Omega_u^I$ to global DoF indices, and $\mathbf{f}_I^{\mathrm{local}} = [\ell_{\Omega_s^I}(\vec{N}_1^I), \dots, \ell_{\Omega_s^I}(\vec{N}_6^I)]^T \in \mathbb{R}^6$ is the local load vector evaluated over the exact same active subdomain $\Omega_s^I \cap \Omega$.

Because each local block $\mathbf{K}_I^{\mathrm{local}}$ and $\mathbf{M}_I^{\mathrm{local}}$ is an outer-product or Gram matrix of real vectors, each is individually symmetric and positive-(semi)definite:
\begin{equation}
    \mathbf{K}_I^{\mathrm{local}} = (\mathbf{K}_I^{\mathrm{local}})^T \ge 0, \qquad \mathbf{M}_I^{\mathrm{local}} = (\mathbf{M}_I^{\mathrm{local}})^T > 0.
\end{equation}
Consequently, their congruence sum \eqref{eq:global_assembly} guarantees that the global matrices are \textbf{strictly symmetric and positive-(semi)definite}:
\begin{equation}
    \mathbf{K} = \mathbf{K}^T \ge 0, \qquad \mathbf{M} = \mathbf{M}^T > 0.
    \label{eq:symmetric_global_properties}
\end{equation}
Imposing essential PEC boundary conditions by eliminating boundary DoFs preserves this symmetry on the condensed generalized eigenvalue problem $\mathbf{K}_{\mathrm{red}} \mathbf{e}_{\mathrm{red}} = \lambda \mathbf{M}_{\mathrm{red}} \mathbf{e}_{\mathrm{red}}$, strictly guaranteeing real eigenvalues ($\lambda = k_c^2 \in \mathbb{R}$) and precluding complex spurious modes.

\paragraph{Effect of Overlap Multiplicity on Global Assembly:}
With support size $\alpha_s = 0.80$, the cumulative subdomain area exceeds the physical domain measure ($\sum_{I=1}^N \mathrm{Area}(\Omega_s^I \cap \Omega) \approx 2\pi \alpha_s^2 \, \mathrm{Area}(\Omega) \approx 4\,\mathrm{Area}(\Omega)$), and the pointwise overlap multiplicity $M(\mathbf{x})$ varies near truncated boundaries and under nodal jitter. As established theoretically in Section~\ref{subsec:local_weak_form} (Eqs.~\eqref{eq:subdomain_residual_zero}--\eqref{eq:global_residual_sum}), this redundant coverage does not introduce multi-counting errors or require partition-of-unity renormalization weights. Because the discrete operators and load vectors are assembled congruently from identically balanced local subdomains ($\mathbf{r}_I(\mathbf{e}^*) \equiv \mathbf{0}$), the global assembly \eqref{eq:global_assembly} directly inherits exact machine-precision consistency on affine fields and optimal $\mathcal{O}(h_{\mathrm{nom}}^2)$ asymptotic convergence on general smooth fields without partition-of-unity distortion.

Both assembly modes (the classical row-by-row testing of Section~\ref{subsubsec:row_mlpg_symmetry} and the symmetric Subdomain Galerkin assembly of Section~\ref{subsubsec:symmetric_assembly}) are supported in our implementation.

\subsection{Discrete Operator Rank Deficiency and Orthogonal Energy Stabilization}
\label{subsec:rank_deficiency_stabilization}

While the closed-form integration over circular disks $\Omega_s^I$ yields exact geometric moments without quadrature grids, a critical structural limitation of the linear polynomial trial space $\mathcal{P}^1$ arises when evaluating the local stiffness matrix. A rigorous algebraic analysis reveals that the standard cell-free stiffness operator exhibits an intrinsic rank deficiency analogous to single-point reduced integration in finite elements, which must be stabilized to achieve quadratic asymptotic convergence in spectral eigenvalue problems.

\subsubsection{Local Operator Rank Analysis under Linear Vector Basis}
\label{subsubsec:local_rank_analysis}

Within each support stencil $\Omega_u^I$, the local vector electric field is approximated via the complete linear polynomial expansion \eqref{eq:P1_shape_function_expansion_mlpg}--\eqref{eq:field_interpolation_mlpg}:
\begin{equation}
    \vec{E}(\mathbf{x}) = \begin{bmatrix} a_1 + a_3 (x - x_I) + a_4 (y - y_I) \\ a_2 + a_5 (x - x_I) + a_6 (y - y_I) \end{bmatrix}, \qquad \mathbf{a} = \boldsymbol{\beta}^I \mathbf{u}_I,
    \label{eq:p1_field_vector_expanded}
\end{equation}
where $\mathbf{u}_I \equiv \mathbf{e}_{\mathrm{stencil}} = \mathbf{L}_I \mathbf{e} \in \mathbb{R}^6$ represents the vector of directional nodal degrees of freedom in stencil $\Omega_u^I$, $\mathbf{a} = [a_1, \dots, a_6]^T \in \mathbb{R}^6$ collects the polynomial basis coefficients, and $\boldsymbol{\beta}^I = (\mathbf{A}_I)^{-1} \in \mathbb{R}^{6 \times 6}$ is the inverse collocation matrix \eqref{eq:beta_matrix_inversion}.

Because the basis functions in $\mathcal{P}^1$ are linear in spatial coordinates, their spatial derivatives are strictly constant on the disk $\Omega_s^I$ and match Eqs.~\eqref{eq:curl_constant}--\eqref{eq:div_constant}:
\begin{align}
    (\nabla \times \vec{E})_z & = a_5 - a_4 = (\boldsymbol{\beta}_{5, :}^I - \boldsymbol{\beta}_{4, :}^I) \mathbf{u}_I \equiv (\mathbf{v}_{\mathrm{curl}}^I)^T \mathbf{u}_I, \\
    \nabla \cdot \vec{E}      & = a_3 + a_6 = (\boldsymbol{\beta}_{3, :}^I + \boldsymbol{\beta}_{6, :}^I) \mathbf{u}_I \equiv (\mathbf{v}_{\mathrm{div}}^I)^T \mathbf{u}_I,
\end{align}
where $\mathbf{v}_{\mathrm{curl}}^I \equiv \mathbf{c}_{\mathrm{curl}}^I \in \mathbb{R}^6$ and $\mathbf{v}_{\mathrm{div}}^I \equiv \mathbf{c}_{\mathrm{div}}^I \in \mathbb{R}^6$ are column vectors.
Consequently, evaluating the weak-form stiffness bilinear forms \eqref{eq:k_curl_analytic}--\eqref{eq:k_div_analytic} over the active subdomain $\Omega_s^I$ of area $S_0^I = \pi (r_s^I)^2$ yields dyadic outer products of single discrete vectors:
\begin{equation}
    \mathbf{K}_{I, \mathrm{curl}}^{\mathrm{local}} = \frac{S_0^I}{\mu_r} \mathbf{v}_{\mathrm{curl}}^I (\mathbf{v}_{\mathrm{curl}}^I)^T, \qquad
    \mathbf{K}_{I, \mathrm{div}}^{\mathrm{local}} = \frac{s_{\mathrm{div}} S_0^I}{\mu_r} \mathbf{v}_{\mathrm{div}}^I (\mathbf{v}_{\mathrm{div}}^I)^T.
\end{equation}
Each constituent matrix is strictly of \textbf{rank 1}. Therefore, the unregularized local stiffness matrix has an algebraic rank bounded by:
\begin{equation}
    \mathrm{rank}(\mathbf{K}_I^{\mathrm{local}}) = \mathrm{rank}\left(\mathbf{K}_{I, \mathrm{curl}}^{\mathrm{local}} + \mathbf{K}_{I, \mathrm{div}}^{\mathrm{local}}\right) \le 1 + 1 = 2.
    \label{eq:stiffness_rank_deficiency}
\end{equation}
In contrast, the local mass matrix $\mathbf{M}_I^{\mathrm{local}}$ \eqref{eq:m_local_analytic} integrates zeroth-order polynomial products and has \textbf{full rank 6}.

This creates a fundamental structural imbalance: out of the 6 degrees of freedom present in each stencil, \textbf{4 degrees of freedom possess identically zero local stiffness}. In the cell-based VNMM-Gauss \cite{stable_vnmm2d_2026}, this rank deficiency is prevented because the $3 \times 3 = 9$ Gauss points within each background Cartesian cell sample distinct neighboring stencils $\Omega_u(P_g)$, coupling inter-stencil derivative jumps to yield a full-rank ($\mathrm{rank} = 6$) cell stiffness matrix. In cell-free MLPG with single-subdomain integration, however, each disk $\Omega_s^I$ evaluates a single fixed polynomial stencil $\Omega_u^I$; consequently, the unconstrained 4-dimensional local nullspace permits discrete modes to soften, which accounts for the sublinear eigenvalue error scaling ($p \approx 0.53$ on regular grids and $p \approx 0.84$ under $25\%$ jitter) in the unaugmented formulation, as demonstrated numerically in Sections~\ref{subsubsec:cavity_jitter_error} and \ref{subsubsec:resolution_stagnation_stabilization}.

\subsubsection{Canonical Orthogonal Decomposition of the Gradient Tensor}
\label{subsubsec:canonical_decomposition}

To eliminate this rank deficiency without corrupting the consistency of physical fields, we analyze the structure of the 4-dimensional local nullspace $\ker(\mathbf{K}_I^{\mathrm{local}})$ via the classical Cauchy--Stokes decomposition of a 2D linear vector gradient tensor into antisymmetric rotation, volumetric trace, and symmetric traceless (deviatoric/shear) components. In polynomial coefficient space $\mathbb{R}^6$ (with canonical unit vectors $\mathbf{e}_1, \dots, \mathbf{e}_6$), any vector field satisfying both $(\nabla \times \vec{E})_z = 0$ ($a_5 = a_4$) and $\nabla \cdot \vec{E} = 0$ ($a_6 = -a_3$) is parametrized by:
\begin{equation}
    \mathbf{a} = a_1 \mathbf{e}_1 + a_2 \mathbf{e}_2 + a_3 (\mathbf{e}_3 - \mathbf{e}_6) + a_4 (\mathbf{e}_4 + \mathbf{e}_5).
\end{equation}
This nullspace decomposes into two physically distinct categories:
\begin{enumerate}
    \item \textbf{Admissible Constant Translation Modes ($a_1, a_2$):} The components $\vec{E} = [a_1, a_2]^T$ represent uniform vector fields. In continuous electromagnetics, constant fields have identically zero curl and zero divergence. They constitute the physical zero-energy gauge subspace and must retain strictly zero stiffness to pass the Patch Test.
    \item \textbf{Harmonic Zero-Energy Linear Deformation Modes ($a_3, a_4$):} The components $\vec{E}_{\mathrm{hyp}} = [x - x_I, -(y - y_I)]^T = \frac{1}{2}\nabla((x-x_I)^2 - (y-y_I)^2)$ and $\vec{E}_{\mathrm{shear}} = [y - y_I, x - x_I]^T = \nabla((x-x_I)(y-y_I))$ are harmonic conjugate gradient fields ($\Delta\phi = 0$). Because their continuous curl and divergence vanish identically throughout $\mathbb{R}^2$, any single-stencil div-curl integral assigns them strictly zero energy regardless of quadrature order, even though their tensor gradient $\nabla\vec{E}$ is non-zero (pure deviatoric normal strain $\varepsilon_{xx} - \varepsilon_{yy} \ne 0$ and pure shear strain $\varepsilon_{xy} \ne 0$). In a bounded PEC cavity, non-trivial harmonic gradients cannot satisfy tangential Dirichlet boundary conditions on all four walls; however, on local interior disks with single-stencil evaluation, they evade discrete curl and divergence penalties entirely.
\end{enumerate}

To isolate these modes, we define the four canonical derivative selector vectors in polynomial coefficient space $\mathbb{R}^6$:
\begin{equation}
    \mathbf{d}_{\mathrm{curl}} = \mathbf{e}_5 - \mathbf{e}_4, \quad
    \mathbf{d}_{\mathrm{div}} = \mathbf{e}_3 + \mathbf{e}_6, \quad
    \mathbf{d}_{\mathrm{shear}} = \mathbf{e}_5 + \mathbf{e}_4, \quad
    \mathbf{d}_{\mathrm{hyp}} = \mathbf{e}_3 - \mathbf{e}_6,
    \label{eq:canonical_d_vectors}
\end{equation}
which map via $\mathbf{v}_\alpha^I = (\boldsymbol{\beta}^I)^T \mathbf{d}_\alpha$ to the complete basis of the 4-dimensional nodal derivative space:
\begin{align}
    \mathbf{v}_{\mathrm{curl}}^I  & \equiv \mathbf{c}_{\mathrm{curl}}^I = (\boldsymbol{\beta}_{5, :}^I - \boldsymbol{\beta}_{4, :}^I)^T \quad (\text{rotational / antisymmetric}), \label{eq:v_curl_def}                   \\
    \mathbf{v}_{\mathrm{div}}^I   & \equiv \mathbf{c}_{\mathrm{div}}^I = (\boldsymbol{\beta}_{3, :}^I + \boldsymbol{\beta}_{6, :}^I)^T \quad (\text{volumetric dilation / trace}), \label{eq:v_div_def}                     \\
    \mathbf{v}_{\mathrm{shear}}^I & = (\boldsymbol{\beta}_{5, :}^I + \boldsymbol{\beta}_{4, :}^I)^T \quad (\text{symmetric shear strain } 2\varepsilon_{xy}), \label{eq:v_shear_def}                   \\
    \mathbf{v}_{\mathrm{hyp}}^I   & = (\boldsymbol{\beta}_{3, :}^I - \boldsymbol{\beta}_{6, :}^I)^T \quad (\text{deviatoric normal strain } \varepsilon_{xx} - \varepsilon_{yy}). \label{eq:v_hyp_def}
\end{align}
\begin{lemma}[Strict Mutual Orthogonality of Gradient Modes]
    Let $\mathbf{D} = [\mathbf{d}_{\mathrm{curl}}, \mathbf{d}_{\mathrm{div}}, \mathbf{d}_{\mathrm{shear}}, \mathbf{d}_{\mathrm{hyp}}]^T \in \mathbb{R}^{4 \times 6}$ and $\mathbf{V}_I = [\mathbf{v}_{\mathrm{curl}}^I, \mathbf{v}_{\mathrm{div}}^I, \mathbf{v}_{\mathrm{shear}}^I, \mathbf{v}_{\mathrm{hyp}}^I]^T = \mathbf{D} \boldsymbol{\beta}^I \in \mathbb{R}^{4 \times 6}$. The canonical derivative selectors are strictly mutually orthogonal in $\mathbb{R}^6$, and the nodal derivative mode vectors $\mathbf{v}_\alpha^I$ are strictly mutually orthogonal with respect to the placement metric $\mathbf{G}_I = \mathbf{A}_I^T \mathbf{A}_I = (\boldsymbol{\beta}^I (\boldsymbol{\beta}^I)^T)^{-1}$:
    \begin{equation}
        \mathbf{V}_I (\mathbf{A}_I^T \mathbf{A}_I) \mathbf{V}_I^T = \mathbf{D} \mathbf{D}^T =
        \begin{bmatrix}
            \|\mathbf{d}_{\mathrm{curl}}\|^2 & 0                               & 0                                 & 0                               \\
            0                                & \|\mathbf{d}_{\mathrm{div}}\|^2 & 0                                 & 0                               \\
            0                                & 0                               & \|\mathbf{d}_{\mathrm{shear}}\|^2 & 0                               \\
            0                                & 0                               & 0                                 & \|\mathbf{d}_{\mathrm{hyp}}\|^2
        \end{bmatrix}
        = 2 \, \mathbf{I}_4.
        \label{eq:orthogonality_gram_matrix}
    \end{equation}
\end{lemma}
\begin{proof}
    Direct computation of the Euclidean inner products in polynomial coefficient space $\mathbb{R}^6$ yields $\mathbf{d}_{\mathrm{curl}} \cdot \mathbf{d}_{\mathrm{div}} = 0$, $\mathbf{d}_{\mathrm{curl}} \cdot \mathbf{d}_{\mathrm{shear}} = 1 - 1 = 0$, $\mathbf{d}_{\mathrm{curl}} \cdot \mathbf{d}_{\mathrm{hyp}} = 0$, $\mathbf{d}_{\mathrm{div}} \cdot \mathbf{d}_{\mathrm{shear}} = 0$, $\mathbf{d}_{\mathrm{div}} \cdot \mathbf{d}_{\mathrm{hyp}} = 1 - 1 = 0$, and $\mathbf{d}_{\mathrm{shear}} \cdot \mathbf{d}_{\mathrm{hyp}} = 0$, with squared norms $\|\mathbf{d}_\alpha\|^2 = 1^2 + (\pm 1)^2 = 2$. Substituting $\mathbf{V}_I = \mathbf{D}\boldsymbol{\beta}^I$ and $\boldsymbol{\beta}^I \mathbf{A}_I = \mathbf{I}_6$ gives $\mathbf{V}_I \mathbf{A}_I^T \mathbf{A}_I \mathbf{V}_I^T = \mathbf{D}(\boldsymbol{\beta}^I \mathbf{A}_I)(\boldsymbol{\beta}^I \mathbf{A}_I)^T \mathbf{D}^T = \mathbf{D}\mathbf{D}^T = 2\mathbf{I}_4$. Because $\mathbf{A}_I$ is non-singular on every qualified stencil, $\mathrm{rank}(\mathbf{V}_I) = 4$.
\end{proof}

\subsubsection{Derivation of the Orthogonal Energy Stabilization Operator}
\label{subsubsec:derivation_orthogonal_stabilization}

Exploiting the mutual orthogonality established in \eqref{eq:orthogonality_gram_matrix}, we construct an \textbf{Orthogonal Energy Stabilization Matrix} $\mathbf{K}_{I, \mathrm{stab}}^{\mathrm{local}}$ that penalizes exclusively the two zero-curl, zero-divergence deformation modes ($\mathbf{v}_{\mathrm{shear}}^I$ and $\mathbf{v}_{\mathrm{hyp}}^I$):
\begin{equation}
    \mathbf{K}_{I, \mathrm{stab}}^{\mathrm{local}} = \frac{\alpha_{\mathrm{stab}}}{\mu_r} S_0^I \left[ \mathbf{v}_{\mathrm{shear}}^I (\mathbf{v}_{\mathrm{shear}}^I)^T + \mathbf{v}_{\mathrm{hyp}}^I (\mathbf{v}_{\mathrm{hyp}}^I)^T \right],
    \label{eq:k_stab_analytic}
\end{equation}
where $\alpha_{\mathrm{stab}} \ge 0$ is a non-dimensional stabilization parameter and $S_0^I = A_{\mathrm{active}}^I = \eta_I \pi (r_s^I)^2$ is the active subdomain area (ensuring that the exact boundary truncation scaling $\mathbf{K}_I^{\mathrm{trunc}} = \eta_I \mathbf{K}_I^{\mathrm{full}}$ in \eqref{eq:k_truncated_scaled} holds identically for the stabilized operator). The augmented local stiffness matrix becomes:
\begin{equation}
    \mathbf{K}_I^{\mathrm{local, stab}} = \mathbf{K}_{I, \mathrm{curl}}^{\mathrm{local}} + \mathbf{K}_{I, \mathrm{div}}^{\mathrm{local}} + \mathbf{K}_{I, \mathrm{stab}}^{\mathrm{local}}.
    \label{eq:k_augmented_local}
\end{equation}

The operator \eqref{eq:k_stab_analytic} satisfies six fundamental properties:
\begin{enumerate}
    \item \textbf{Exact Nullspace for Translations, Pure Rotations, and Isotropic Dilations:} By \eqref{eq:orthogonality_gram_matrix}, $(\mathbf{v}_{\mathrm{shear}}^I)^T \mathbf{u}_I = a_5 + a_4$ and $(\mathbf{v}_{\mathrm{hyp}}^I)^T \mathbf{u}_I = a_3 - a_6$. Consequently, for any uniform field $\vec{E}_{\mathrm{const}} = [a_1, a_2]^T$, pure rigid-body rotation $\vec{E}_{\mathrm{rot}} = [-y, x]^T$ ($a_4 = -1, a_5 = 1 \implies a_5 + a_4 = 0$), or pure isotropic volumetric dilation $\vec{E}_{\mathrm{iso}} = [x, y]^T$ ($a_3 = a_6 = 1 \implies a_3 - a_6 = 0$), the projections onto both stabilization modes vanish identically:
          \begin{equation}
              \mathbf{K}_{I, \mathrm{stab}}^{\mathrm{local}} \mathbf{u}_{\mathrm{const}} = \mathbf{K}_{I, \mathrm{stab}}^{\mathrm{local}} \mathbf{u}_{\mathrm{rot}} = \mathbf{K}_{I, \mathrm{stab}}^{\mathrm{local}} \mathbf{u}_{\mathrm{iso}} = \mathbf{0}.
          \end{equation}
          The stabilization adds strictly zero energy to rigid translations, pure rotational curl states, and isotropic divergence states without requiring any auxiliary load terms.
    \item \textbf{Restoration of Full Local Rank:} Adding the two rank-1 dyadic products in \eqref{eq:k_stab_analytic} elevates the local rank from 2 to \textbf{rank 4}:
          \begin{equation}
              \mathrm{rank}\left(\mathbf{K}_I^{\mathrm{local, stab}}\right) = 1 + 1 + 1 + 1 = 4.
          \end{equation}
          Because 2 degrees of freedom correspond strictly to physical constant translations, rank 4 represents the theoretical maximum rank achievable for the derivative space of $\mathcal{P}^1$.
    \item \textbf{Coordinate-Frame Invariance (Objectivity / Isotropy):} Although the individual projections $\mathbf{v}_{\mathrm{shear}}^I \cdot \mathbf{u}_I = \partial_x E_y^h + \partial_y E_x^h$ and $\mathbf{v}_{\mathrm{hyp}}^I \cdot \mathbf{u}_I = \partial_x E_x^h - \partial_y E_y^h$ depend on the orientation of the Cartesian axes and rotate into one another by an angle $2\theta$ under an in-plane coordinate rotation $\mathbf{Q}(\theta) \in \mathrm{SO}(2)$, their equal-weight quadratic sum in \eqref{eq:k_stab_analytic} forms an exact rotational invariant:
          \begin{equation}
              \begin{aligned}
                  (\mathbf{v}_{\mathrm{shear}}^I \cdot \mathbf{u}_I)^2 + (\mathbf{v}_{\mathrm{hyp}}^I \cdot \mathbf{u}_I)^2 & = (\partial_x E_y^h + \partial_y E_x^h)^2 + (\partial_x E_x^h - \partial_y E_y^h)^2 \\
                                                                                                                            & = 2 \, \|\operatorname{dev}\boldsymbol{\varepsilon}(\vec{E}^h)\|_F^2,
              \end{aligned}
          \end{equation}
          where $\operatorname{dev}\boldsymbol{\varepsilon}(\vec{E}^h) = \boldsymbol{\varepsilon}(\vec{E}^h) - \frac{1}{2}(\nabla \cdot \vec{E}^h)\mathbf{I}_2$ is the symmetric traceless (deviatoric) gradient tensor and $\|\cdot\|_F$ is the Frobenius norm. Consequently, $\mathbf{K}_{I, \mathrm{stab}}^{\mathrm{local}}$ is strictly isotropic and objective with respect to arbitrary rotations of the global coordinate system.
    \item \textbf{Complementary Roles of $s_{\mathrm{div}}$ and $\alpha_{\mathrm{stab}}$:} It is important to emphasize that $\mathbf{K}_{\mathrm{stab}}$ does \emph{not} replace the divergence penalty parameter $s_{\mathrm{div}}$. The term $s_{\mathrm{div}} \mathbf{K}_{\mathrm{div}}$ acts on $\mathbf{v}_{\mathrm{div}}$ to expel non-physical gradient modes ($\vec{E} = \nabla\phi$ with $\nabla^2\phi \ne 0$) to high frequencies, whereas $\mathbf{K}_{\mathrm{stab}}$ acts on the orthogonal subspace $\{\mathbf{v}_{\mathrm{shear}}, \mathbf{v}_{\mathrm{hyp}}\}$ with zero divergence. Both operators are mutually orthogonal and operate in concert.
    \item \textbf{Variational Consistency and VMS Asymptotic Compatibility (Patch Test and MMS):} For affine states containing non-zero symmetric shear ($a_5 + a_4 \ne 0$) or deviatoric normal strain ($a_3 - a_6 \ne 0$), two regimes are distinguished: (a) when the local linear functional $\ell_{\Omega_s^I}(\vec{W}_I)$ includes the consistent stabilization boundary flux $a_{\Omega_s^I}^{\mathrm{stab}}(\vec{W}_I, \vec{E}^*)$, all seven canonical affine states (Modes~1--7) are reproduced to machine precision ($\le 1.3 \times 10^{-13}$) for arbitrary $\alpha_{\mathrm{stab}} > 0$ (Remark~\ref{rem:patch_test_stab}); (b) when $\mathbf{K}_{\mathrm{stab}}$ is applied purely to the left-hand bilinear form without auxiliary strain fluxes (as in homogeneous cavity eigenproblems and standard source-driven MMS), it acts as a weakly consistent VMS subgrid operator whose magnitude $\alpha_{\mathrm{stab}} = c_{\mathrm{stab}} (h_{\mathrm{nom}}/L)^2$ matches the $\mathcal{O}(h_{\mathrm{nom}}^2)$ $\mathcal{P}^1$ interpolation truncation error, strictly preserving the optimal rates of $\mathcal{O}(h_{\mathrm{nom}}^2)$ in $L^2(\vec{E})$ and $\mathcal{O}(h_{\mathrm{nom}})$ in the $H(\mathbf{curl})$ semi-norm $|\cdot|_{H(\mathbf{curl})}$ (Section~\ref{subsec:stage3_mms}).
    \item \textbf{Connection to Orthogonal Hourglass Control and Modified Nodal Integration (MSCNI/SNNI):} Conceptually, projecting out zero-energy deformation modes orthogonally to the active stiffness modes parallels classical orthogonal hourglass control in under-integrated finite elements \cite{flanagan1981uniform} and stabilized nodal integration in meshfree mechanics (Modified SCNI and SNNI \cite{chen2001stabilized, chen2013arbitrary, hillman2018performance}). However, whereas MSCNI/SNNI stabilize rank-deficient modes by subdividing each nodal smoothing cell into multiple sub-cells and evaluating finite strain differences across sub-cell centroids, the explicit polynomial basis $\mathcal{P}^1$ of VNMM-MLPG allows Lemma~1 to project the two harmonic zero-curl, zero-divergence modes ($\mathbf{v}_{\mathrm{shear}}^I$ and $\mathbf{v}_{\mathrm{hyp}}^I$) directly onto the directional degrees of freedom in closed form via rows of $\boldsymbol{\beta}^I$, restoring full local rank 4 without sub-cell partitioning or auxiliary quadrature points.
\end{enumerate}

\subsubsection{Dimensionless Variational Multiscale Scaling Law (\texorpdfstring{$\alpha_{\mathrm{stab}} \propto (h_{\mathrm{nom}}/L)^2$}{alpha\_stab proportional to (h\_nom/L)2})}
\label{subsubsec:vms_scaling_law}

In accordance with Variational Multiscale (VMS) subgrid theory \cite{hughes1998variational, codina2000stabilization}, the stabilization parameter must not remain a constant $\mathcal{O}(1)$ across grid refinement; an $\mathcal{O}(1)$ penalty over-stiffens physical cavity modes (whose continuous fields contain genuine spatial shear). Instead, the stabilization must scale with the unresolved subgrid truncation error. To ensure rigorous dimensional consistency and invariance across arbitrary physical scales and unit conventions, we normalize the nominal nodal spacing $h_{\mathrm{nom}}$ by a macroscopic domain reference dimension $L = \max(L_x, L_y)$ (with $L = \pi\text{ m}$ for the standard square cavity), formulating the dimensionless quadratic scaling law:
\begin{equation}
    \alpha_{\mathrm{stab}}(h_{\mathrm{nom}}) = c_{\mathrm{stab}} \left(\frac{h_{\mathrm{nom}}}{L}\right)^2,
    \label{eq:alpha_stab_vms_scaling}
\end{equation}
where $c_{\mathrm{stab}} \ge 0$ is a strictly dimensionless calibration parameter. On general unstructured or locally adaptive point clouds that lack a global Cartesian resolution parameter $N_x$, $h_{\mathrm{nom}}$ in \eqref{eq:alpha_stab_vms_scaling} is naturally replaced by either the local nodal spacing $h_I$ on each disk ($\alpha_{\mathrm{stab}}^I = c_{\mathrm{stab}} (h_I / L)^2$) or the cloud-averaged spacing $\overline{h} = \frac{1}{N}\sum_{I=1}^N h_I$. As demonstrated in Section~\ref{subsubsec:resolution_stagnation_stabilization}, values in the range $c_{\mathrm{stab}} \in [40, 100]$ systematically accelerate the empirical convergence rate from sublinear ($p = 0.53$) to full quadratic ($p = 2.39$), resolving the eigenvalue error stagnation on dense grids while offering an explicit control parameter for balancing coarse-grid stiffness and fine-grid asymptotic accuracy.

\subsection{Analysis and Treatment of Boundary Conditions}
\label{subsec:boundary_treatment}

A well-known challenge in meshless methods, particularly within the MLPG paradigm, is the enforcement of boundary conditions. In conventional scalar meshless methods based on Moving Least Squares (MLS) or radial basis functions (RBF), shape functions lack the Kronecker-delta property ($\phi_i(\mathbf{x}_j) \ne \delta_{ij}$), requiring complex external mechanisms such as Lagrange multipliers, penalty parameters, or boundary element coupling \cite{atluri2002meshless}.

In the VNMM-MLPG framework, the directional vector degree of freedom $e_k = \vec{E}(\mathbf{x}_k) \cdot \vec{t}_k$ combined with the Subdomain Galerkin / LPIM weak form \eqref{eq:local_weak_form} provides a direct, consistent treatment for standard electromagnetic boundary conditions. Consider a general partition of the physical boundary $\partial\Omega = \Gamma_D \cup \Gamma_N \cup \Gamma_R$, where $\Gamma_D$ denotes Dirichlet boundaries (electric walls / PEC), $\Gamma_N$ denotes Neumann boundaries (magnetic walls / PMC / surface current sheets), and $\Gamma_R$ denotes Robin boundaries (impedance / absorbing radiation).

For any node $\mathbf{x}_I$ whose local subdomain $\Omega_s^I$ intersects the physical boundary $\partial\Omega$, the active integration domain is the truncated region $\Omega_s^I \cap \Omega$, and the physical boundary intersection segment $\Gamma_s^I = \Omega_s^I \cap \partial\Omega$ partitions naturally according to the boundary types:
\begin{equation}
    \Gamma_s^I = (\Omega_s^I \cap \Gamma_D) \cup (\Omega_s^I \cap \Gamma_N) \cup (\Omega_s^I \cap \Gamma_R).
\end{equation}
In the local weak form \eqref{eq:local_weak_form}, physical boundary conditions enter exclusively through the natural boundary flux line integral along $\Gamma_s^I$:
\begin{equation}
    I_{\Gamma_s^I} = \int_{\Gamma_s^I} \vec{W}_I \cdot \vec{\tau}_{\mathrm{ext}} \, ds,
    \label{eq:boundary_flux_partition}
\end{equation}
where $\vec{\tau}_{\mathrm{ext}}$ denotes the prescribed physical surface boundary flux corresponding to the natural electromagnetic flux vector along outward unit normal $\vec{n}$:
\begin{equation}
    \vec{\tau} = \frac{1}{\mu_r} \left( \vec{n} \times (\nabla \times \vec{E}) \right) + \frac{s_{\mathrm{div}}}{\mu_r} (\nabla \cdot \vec{E}) \vec{n}.
    \label{eq:traction_vector}
\end{equation}
We now analyze the enforcement of each physical boundary condition along $\Gamma_s^I$ in detail.

\subsubsection{Inhomogeneous and Homogeneous Dirichlet Boundary Conditions}
\label{subsubsec:dirichlet_bc}

On a Dirichlet boundary $\Gamma_D$, the tangential trace of the electric field is prescribed:
\begin{equation}
    \vec{n} \times \vec{E} = \vec{g}_D \quad \iff \quad E_t(\mathbf{x}) = g_D(\mathbf{x}) \quad \forall \mathbf{x} \in \Gamma_D.
    \label{eq:dirichlet_prescribed}
\end{equation}
In the VNMM, the tangential director $\vec{t}_k$ of each boundary node $\mathbf{x}_k \in \Gamma_D$ is aligned strictly parallel to the boundary tangent ($\vec{t}_k \parallel \Gamma_D$). Because the VNMM trial space satisfies the exact directional nodal projection property at the node itself:
\begin{equation}
    e_k = \vec{E}(\mathbf{x}_k) \cdot \vec{t}_k = g_D(\mathbf{x}_k).
    \label{eq:dirichlet_nodal_value}
\end{equation}
Thus, essential Dirichlet conditions are satisfied directly at the discrete nodal level without any Lagrange multipliers or penalty parameters. Partitioning the global DoF vector into active interior DoFs $\mathbf{e}_{\mathrm{int}}$ and prescribed boundary DoFs $\mathbf{e}_{\mathrm{bnd}}$:
\begin{equation}
    \begin{bmatrix}
        \mathbf{A}_{\mathrm{int}, \mathrm{int}} & \mathbf{A}_{\mathrm{int}, \mathrm{bnd}} \\
        \mathbf{A}_{\mathrm{bnd}, \mathrm{int}} & \mathbf{A}_{\mathrm{bnd}, \mathrm{bnd}}
    \end{bmatrix}
    \begin{bmatrix}
        \mathbf{e}_{\mathrm{int}} \\
        \mathbf{e}_{\mathrm{bnd}}
    \end{bmatrix}
    =
    \begin{bmatrix}
        \mathbf{F}_{\mathrm{int}} \\
        \mathbf{F}_{\mathrm{bnd}}
    \end{bmatrix},
\end{equation}
the exact boundary values $\mathbf{e}_{\mathrm{bnd}} = [g_D(\mathbf{x}_k)]$ are substituted directly into the system, yielding the condensed, strictly symmetric positive-definite linear system matching \cite{stable_vnmm2d_2026}:
\begin{equation}
    \mathbf{A}_{\mathrm{int}, \mathrm{int}} \mathbf{e}_{\mathrm{int}} = \mathbf{F}_{\mathrm{int}} - \mathbf{A}_{\mathrm{int}, \mathrm{bnd}} \mathbf{e}_{\mathrm{bnd}}.
    \label{eq:dirichlet_condensed_system}
\end{equation}
For homogeneous Dirichlet boundaries (PEC, with $g_D \equiv 0$), $\mathbf{e}_{\mathrm{bnd}} = \mathbf{0}$, reducing the source problem to $\mathbf{A}_{\mathrm{int}, \mathrm{int}} \mathbf{e}_{\mathrm{int}} = \mathbf{F}_{\mathrm{int}}$. Similarly, in resonant cavity eigenvalue problems, imposing $\mathbf{e}_{\mathrm{bnd}} = \mathbf{0}$ yields the condensed generalized eigenvalue problem:
\begin{equation}
    \mathbf{K}_{\mathrm{red}} \mathbf{e}_{\mathrm{red}} = \lambda \mathbf{M}_{\mathrm{red}} \mathbf{e}_{\mathrm{red}}, \qquad \mathbf{K}_{\mathrm{red}} = \mathbf{K}_{\mathrm{int}, \mathrm{int}}, \quad \mathbf{M}_{\mathrm{red}} = \mathbf{M}_{\mathrm{int}, \mathrm{int}},
    \label{eq:reduced_eigenproblem}
\end{equation}
with $\mathbf{e}_{\mathrm{red}} \equiv \mathbf{e}_{\mathrm{int}}$ matching the formulation in \cite{stable_vnmm2d_2026}.
Because boundary DoFs are eliminated directly, evaluating boundary integrals along $\Gamma_D$ is avoided entirely.

\subsubsection{Homogeneous Neumann Boundary Conditions (Natural / PMC / Symmetry)}
\label{subsubsec:neumann_homogeneous}

On a homogeneous Neumann boundary $\Gamma_N$ (such as a perfect magnetic conductor, PMC, or a plane of electromagnetic geometric symmetry), the natural boundary condition requires the tangential magnetic field (surface curl flux) to vanish:
\begin{equation}
    \vec{n} \times (\nabla \times \vec{E}) = \mathbf{0} \quad \text{on } \Gamma_N.
    \label{eq:neumann_homogeneous_pde}
\end{equation}
Moreover, for physical solenoidal fields satisfying the divergence-free condition in source-free regions, $\nabla \cdot \vec{E} = 0$. Substituting these into the boundary flux vector \eqref{eq:traction_vector} shows that:
\begin{equation}
    \vec{\tau} \equiv \mathbf{0} \quad \text{on } \Gamma_N.
\end{equation}
Consequently, the physical boundary flux integral over $\Gamma_s^I \subset \Gamma_N$ vanishes identically:
\begin{equation}
    \int_{\Omega_s^I \cap \Gamma_N} \vec{\tau} \cdot \vec{W}_I \, ds \equiv 0.
\end{equation}
This establishes a key theoretical property of the VNMM-MLPG: \textbf{homogeneous Neumann boundary conditions are strictly natural}. No line integrals, no boundary segment quadratures, and no boundary corrections are needed. Nodes located on $\Gamma_N$ are simply treated as active, unconstrained degrees of freedom belonging to $\mathbf{e}_{\mathrm{int}}$. Their equations are assembled normally from their local subdomains $\Omega_s^I$, and their tangential fields are resolved automatically as part of the global system solution.

\subsubsection{Inhomogeneous Neumann Boundary Conditions (Surface Currents)}
\label{subsubsec:neumann_inhomogeneous}

When an inhomogeneous Neumann condition is prescribed on $\Gamma_N$, representing an impressed surface magnetic current density $\vec{M}_s$ or a non-zero natural tangential magnetic flux:
\begin{equation}
    \frac{1}{\mu_r} \left( \vec{n} \times (\nabla \times \vec{E}) \right) = \vec{g}_N(\mathbf{x}) \quad \text{on } \Gamma_N,
    \label{eq:neumann_inhomogeneous_pde}
\end{equation}
the boundary flux vector becomes $\vec{\tau} = \vec{g}_N(\mathbf{x})$ (under $\nabla \cdot \vec{E} = 0$). In the local weak form for any node whose subdomain intersects $\Gamma_N$, this boundary flux introduces a non-zero line integral over the intersection segment $\Gamma_s^I = \Omega_s^I \cap \Gamma_N$:
\begin{equation}
    F_a^{\mathrm{Neumann}, I} = \int_{\Gamma_s^I} \vec{g}_N(\mathbf{x}) \cdot \vec{W}_a^I(\mathbf{x}) \, ds = \int_{\Gamma_s^I} \vec{g}_N(\mathbf{x}) \cdot \vec{N}_a^I(\mathbf{x}) \, ds.
    \label{eq:neumann_load_integral}
\end{equation}
In the Local Galerkin formulation (MLPG6), the test basis $\vec{W}_a^I = \vec{N}_a^I(\mathbf{x})$ is a complete linear vector polynomial. For planar boundaries (such as straight cavity walls or polygonal waveguides), the intersection $\Gamma_s^I$ is a simple straight line segment $[-L_0, L_0]$. Because the test functions are linear polynomials in Cartesian coordinates, the line integral \eqref{eq:neumann_load_integral} can be evaluated either in closed analytical form or via a standard 2-point 1D Gauss--Legendre quadrature. The resulting vector $\mathbf{F}^{\mathrm{Neumann}, I}$ is added directly to the local load vector, fully maintaining the cell-free character of the method.

\subsubsection{Mixed Boundary Conditions and Robin / Impedance Boundaries}
\label{subsubsec:mixed_robin_bc}

In general computational domains, mixed boundary conditions are enforced by partitioning the boundary into pairwise disjoint sub-domains $\partial\Omega = \Gamma_D \cup \Gamma_N \cup \Gamma_R$. At transition corners where a Dirichlet wall meets a Neumann or Robin boundary, Dirichlet conditions take precedence, fixing the corner DoF, while adjacent interior nodes on the natural segment remain free.

Furthermore, first-order absorbing (ABC) or surface impedance conditions on $\Gamma_R$ take the Robin form:
\begin{equation}
    \frac{1}{\mu_r} \left( \vec{n} \times (\nabla \times \vec{E}) \right) + \gamma \vec{E}_t = \vec{g}_R(\mathbf{x}) \quad \text{on } \Gamma_R,
    \label{eq:robin_pde}
\end{equation}
where $\gamma = j k_0 \sqrt{\epsilon_r / \mu_r}$ represents the surface admittance and $\vec{E}_t = (\vec{E} \cdot \vec{t})\vec{t}$ is the tangential electric field vector along the unit tangent $\vec{t}$. Substituting \eqref{eq:robin_pde} into the local boundary flux \eqref{eq:boundary_flux_partition} shifts the dependent field term $\gamma \vec{E}_t$ to the bilinear form, yielding a local boundary damping matrix:
\begin{equation}
    C_{ab}^{\mathrm{Robin}, I} = \gamma \int_{\Omega_s^I \cap \Gamma_R} \left( \vec{N}_a^I(\mathbf{x}) \cdot \vec{t} \right) \left( \vec{N}_b^I(\mathbf{x}) \cdot \vec{t} \right) \, ds,
    \label{eq:robin_matrix}
\end{equation}
while the independent term $\vec{g}_R$ contributes to the right-hand load vector. Because $\Omega_s^I \cap \Gamma_R$ is a 1D boundary segment (e.g., a straight segment along planar walls), \eqref{eq:robin_matrix} can be evaluated either in exact closed analytical form---since the shape functions $\vec{N}_a^I$ are linear polynomials in Cartesian coordinates, yielding a quadratic integrand of degree at most 2---or via a standard 2-point 1D Gauss--Legendre quadrature, entirely avoiding background domain meshes.

\subsection{Local Conformity vs. Global Non-Conformity}
\label{subsec:conformity_discussion}

A central conceptual question in vector meshless methods concerns conformity:
\begin{itemize}
    \item \textbf{Local Conformity:} Within each local subdomain $\Omega_s^I$, the trial function $\vec{E}^h(\mathbf{x}) = \mathbf{P}(\mathbf{x} - \mathbf{x}_I) \boldsymbol{\beta}^I \mathbf{e}_{\mathrm{stencil}}$ is a smooth, continuous vector polynomial $\in C^\infty(\Omega_s^I) \subset H(\mathbf{curl}, \Omega_s^I)$. Inside $\Omega_s^I$, there are no internal edges, no jump discontinuities, and no singular distributions.
    \item \textbf{Global Non-Conformity:} Because each node $P_I$ independently selects its own trial stencil $\Omega_u^I$, evaluating $\vec{E}^h(\mathbf{x})$ across overlapping support domains produces inter-stencil trace jumps across shifting support boundaries. Consequently, the global discrete space does not strictly conform to $H(\mathbf{curl}, \Omega)$.
\end{itemize}

In earlier cell-based VNMM-Gauss, non-conformity caused instability: when Gauss points sampled stencil transitions, derivative jumps contaminated the discrete curl, requiring dense over-integration to suppress rank deficiencies. In contrast, VNMM-MLPG resolves this issue by localizing integration:
\begin{enumerate}
    \item Integration is performed strictly over autonomous circular disks $\Omega_s^I$.
    \item Inside each disk $\Omega_s^I$, the trial stencil is fixed and invariant (the 6 nodes of $\Omega_u^I$).
    \item Therefore, the integrand is smooth ($C^\infty$) across the entire integration domain $\Omega_s^I$, eliminating intra-domain jump discontinuities.
    \item As the nominal nodal spacing $h_{\mathrm{nom}} \to 0$, inter-stencil discrepancies between adjacent overlapping disks decay asymptotically at $\mathcal{O}(h_{\mathrm{nom}}^2)$, ensuring asymptotic consistency and convergence.
\end{enumerate}

%% =========================================================================
%% SECTION 3: NUMERICAL VALIDATION AND RESULTS
%% =========================================================================
\section{Numerical Validation and Results}
\label{sec:results}

To establish the mathematical consistency, asymptotic convergence, physical fidelity, and computational efficiency of the VNMM-MLPG, we construct a comprehensive six-stage verification roadmap on the canonical square domain $\Omega = [0, \pi]^2\text{ m}^2$ (adopted throughout to enable a direct, apple-to-apple comparison against the six-mesh benchmark suite of cell-based VNMM-Gauss \cite{stable_vnmm2d_2026} and exact closed-form analytical solutions):

\begin{enumerate}
    \item \textbf{Stage 1: Element-Level Analytical Verification.} Verify that closed-form analytical formulas for stiffness $\mathbf{K}$ and mass $\mathbf{M}$ match high-precision numerical quadrature to machine tolerance across both full circular disks and truncated subdomains (flat walls, corner vertices, and chord cuts), and demonstrate that 2-point Gauss quadrature evaluates 1D boundary terms (Neumann load vectors and Robin damping matrices) to machine precision.
    \item \textbf{Stage 2: Galerkin Patch Test for Linear Consistency.} Verify that under both the classical unsymmetric row-by-row MLPG assembly \eqref{eq:k_row_analytic} and the symmetric Subdomain Galerkin assembly \eqref{eq:global_assembly}, on both regular grids and irregular point clouds with $25\%$ stochastic position jitter, the discrete system reproduces all seven canonical affine vector states to machine precision:
          \begin{equation}
              \frac{\|\mathbf{e}_{\mathrm{int}} - \mathbf{e}^*\|_{\infty}}{\|\mathbf{e}^*\|_{\infty}} \le 10^{-11}.
          \end{equation}
    \item \textbf{Stage 3: Method of Manufactured Solutions (MMS).} Verify asymptotic convergence orders under spatial refinement ($h_{\mathrm{nom}} \to 0$) for an inhomogeneous vector Helmholtz problem with exact manufactured vector field $\vec{E}^*$ under both:
          \begin{enumerate}[(a)]
              \item Pure essential Dirichlet boundary conditions on all four physical walls ($\partial\Omega$).
              \item Mixed Dirichlet--Neumann boundary conditions, where inhomogeneous Dirichlet traces are imposed on $x \in \{0, \pi\}$ and natural homogeneous Neumann conditions ($\vec{n} \times (\nabla \times \vec{E}) = \mathbf{0}$) are prescribed on $y \in \{0, \pi\}$ with boundary nodes treated as active free unknowns.
          \end{enumerate}
          The theoretical polynomial convergence rates are recovered in both configurations:
          \begin{equation}
              \begin{aligned}
                  \|\vec{E}^h - \vec{E}^*\|_{L^2}                                                                       & \sim \mathcal{O}(h_{\mathrm{nom}}^2), \\
                  |\vec{E}^h - \vec{E}^*|_{H(\mathbf{curl})} \equiv \|(\nabla \times \vec{E}^h)_z - (\nabla \times \vec{E}^*)_z\|_{L^2} & \sim \mathcal{O}(h_{\mathrm{nom}}).
              \end{aligned}
          \end{equation}
    \item \textbf{Stage 4: Subdomain Scaling Bounds and Parametric Sensitivity of $\alpha_s$.} Verify the classical MLPG/LPIM domain coverage lower bound ($\alpha_s \ge \sqrt{2}/2 \approx 0.7071$ \cite{atluri2002meshless, liu2009mesh}) alongside the 6-node directional stencil containment upper bound ($\alpha_s \le 1.15$) across a wide sweep $\alpha_s \in [0.30, 2.00]$ on regular and perturbed clouds, confirming $\alpha_s = 0.80$ as a robust operating standard.
    \item \textbf{Stage 5: Physical Spectral Validation on a 2D PEC Cavity.} Employing the universal standard $\alpha_s = 0.80$ and symmetric Subdomain Galerkin assembly \eqref{eq:global_assembly}, solve the generalized eigenvalue problem $\mathbf{K}_{\mathrm{red}} \mathbf{e}_{\mathrm{red}} = k_c^2 \mathbf{M}_{\mathrm{red}} \mathbf{e}_{\mathrm{red}}$ on a square resonant cavity $\Omega = [0, \pi]^2\text{ m}^2$. Compare cutoff wavenumbers $k_c$ for the first 10 physical transverse electric (TE) modes against exact analytical solutions \eqref{eq:analytical_cavity_kc}, assess spurious mode shifting via $s_{\mathrm{div}}$, and formulate the superlinear penalty scaling $s_{\mathrm{div}}(h_{\mathrm{nom}}) \propto h_{\mathrm{nom}}^{-1.5}$ coupled with orthogonal energy stabilization.
    \item \textbf{Stage 6: Computational Benchmark.} Directly benchmark CPU execution time for global matrix assembly, eigensolver solve time, total runtime, and matrix sparsity across the sequence of six grid resolutions from \cite{stable_vnmm2d_2026} ($15 \times 15, 22 \times 22, 29 \times 29, 36 \times 36, 43 \times 43, 49 \times 49$), alongside a direct comparison of cutoff wavenumber errors and convergence rates.
\end{enumerate}

\subsection{Stage 1: Element-Level Closed-Form Verification}
\label{subsec:stage1_results}

To establish the mathematical fidelity and cell-free precision of the local operators before undertaking global system assembly, the element-level stiffness matrix $\mathbf{K}_I^{\mathrm{local}}$, mass matrix $\mathbf{M}_I^{\mathrm{local}}$, and boundary terms were comprehensively verified by comparing closed-form analytical evaluations against high-order numerical quadratures. We investigate three distinct geometric regimes:
\begin{enumerate}
    \item Full interior circular disks ($\eta = 1.0$),
    \item Truncated subdomains along or near physical boundaries ($\eta \in (0, 1)$),
    \item 1D natural boundary integrals entering Neumann and Robin boundary conditions along $\Gamma_s^I = \Omega_s^I \cap \partial\Omega$.
\end{enumerate}

\begin{table}[!t]
    \centering
    \small
    \caption{Element-level verification of stiffness $\mathbf{K}$, mass $\mathbf{M}$, and boundary terms across full and truncated integration domains, comparing analytical closed-form formulas against high-precision numerical quadrature.}
    \label{tab:stage1_element_verification}

    \vspace{0.4em}
    \textbf{(a) 2D Subdomain Integrals: Stiffness $\mathbf{K}$ and Mass $\mathbf{M}$ on Full and Truncated Domains}
    \vspace{0.2em}

    \resizebox{\textwidth}{!}{
        \begin{tabular}{llcccc}
            \toprule
            \textbf{Domain}    & \textbf{Geometric}        & \textbf{Fraction}     & \textbf{Stiffness Error}         & \textbf{Mass Error}                               & \textbf{Scaling Error}                                            \\
            \textbf{Type}      & \textbf{Configuration}    & $\boldsymbol{\eta_I}$ & $\boldsymbol{\Delta \mathbf{K}}$ & $\boldsymbol{\Delta \mathbf{M}}$ \textbf{(Exact)} & $\boldsymbol{\eta \mathbf{M}^{\mathrm{full}}}$ \textbf{vs.~Exact} \\
            \midrule
            Full Circular Disk & Interior node             & $1.0$                 & $2.14 \times 10^{-14}$           & $1.87 \times 10^{-14}$                            & $0.00$                                                            \\
            Flat Boundary Wall & Half-disk ($\eta=1/2$)    & $1/2$                 & $6.91 \times 10^{-17}$           & $2.93 \times 10^{-15}$                            & $1.56 \times 10^{0}$                                              \\
            Corner Vertex      & Quarter-disk ($\eta=1/4$) & $1/4$                 & $1.85 \times 10^{-16}$           & $3.70 \times 10^{-16}$                            & $4.22 \times 10^{-1}$                                             \\
            Near-Wall Cut      & Chord ($d = 0.5 r_s$)     & $0.8045$              & $2.09 \times 10^{-7}$            & $1.69 \times 10^{-7}$                             & $1.55 \times 10^{-1}$                                             \\
            \bottomrule
        \end{tabular}
    }

    \vspace{0.8em}
    \textbf{(b) 1D Natural Boundary Line Integrals along $\Gamma_s^I = \Omega_s^I \cap \partial\Omega$}
    \vspace{0.2em}

    \resizebox{\textwidth}{!}{
        \begin{tabular}{llccc}
            \toprule
            \textbf{Boundary}                                 & \textbf{Intersection}         & \textbf{1-Point Gauss} & \textbf{2-Point Gauss} & \textbf{3-Point Gauss} \\
            \textbf{Term}                                     & \textbf{Segment $\Gamma_s^I$} & \textbf{Error}         & \textbf{Error}         & \textbf{Error}         \\
            \midrule
            Inhomogeneous Neumann $F_a^{\mathrm{Neumann}, I}$ & Straight wall $[-r_s, r_s]$   & $1.01 \times 10^{-1}$  & $1.92 \times 10^{-16}$ & $4.81 \times 10^{-17}$ \\
            Robin Damping Matrix $C_{ab}^{\mathrm{Robin}, I}$ & Straight wall $[-r_s, r_s]$   & $1.76 \times 10^{-1}$  & $1.09 \times 10^{-16}$ & $2.18 \times 10^{-16}$ \\
            \bottomrule
        \end{tabular}
    }
\end{table}

\subsubsection{Full Circular Disks (Interior Subdomains)}
For an uncut circular disk $\Omega_s^I$ ($r_s = 0.35\text{ m}$, $s_{\mathrm{div}} = 4.0$), polar Gauss--Legendre quadrature was executed with $n_r = 6$ radial points and $n_\theta = 16$ circumferential points. Because the trial shape functions are linear polynomials in Cartesian coordinates ($\mathcal{P}^1$), the stiffness integrand is strictly constant, while the mass integrand is a complete quadratic polynomial. As tabulated in Table~\ref{tab:stage1_element_verification}, the relative difference between the closed-form analytical formulas \eqref{eq:k_curl_analytic}--\eqref{eq:m_local_analytic} and numerical quadrature satisfies:
\begin{align}
    \frac{\|\mathbf{K}_{\mathrm{analytic}} - \mathbf{K}_{\mathrm{numeric}}\|_{\max}}{\|\mathbf{K}_{\mathrm{analytic}}\|_{\max}} & \le 2.14 \times 10^{-14}, \\
    \frac{\|\mathbf{M}_{\mathrm{analytic}} - \mathbf{M}_{\mathrm{numeric}}\|_{\max}}{\|\mathbf{M}_{\mathrm{analytic}}\|_{\max}} & \le 1.87 \times 10^{-14},
\end{align}
confirming exact agreement down to double-precision machine epsilon.

\subsubsection{Truncated Subdomains Near Boundaries}
For the truncated boundary subdomains formulated in Section~\ref{subsubsec:truncated_subdomains_analytic} and \ref{app:truncated_geometry} (flat walls, corner vertices, and interior chord cuts), Table~\ref{tab:stage1_element_verification}(a) summarizes the numerical verification against high-order quadrature.

For the stiffness matrix $\mathbf{K}_I^{\mathrm{trunc}}$, the analytical area scaling \eqref{eq:k_truncated_scaled} reproduces numerical integration to machine precision: $\Delta \mathbf{K} \le 6.91 \times 10^{-17}$ on flat walls and $\Delta \mathbf{K} \le 1.85 \times 10^{-16}$ on corners.

For the mass matrix $\mathbf{M}_I^{\mathrm{trunc}}$, Table~\ref{tab:stage1_element_verification}(a) confirms the necessity of the moment-based formulation \eqref{eq:m_truncated_moments}: while simplified area scaling $\eta_I \mathbf{M}_I^{\mathrm{full}}$ incurs severe errors ($156\%$ on half-disks, $42\%$ on quarter-disks, and $15.5\%$ on chord cuts), evaluating the complete moment expansion recovers the quadrature results down to machine epsilon ($\Delta \mathbf{M} \le 2.93 \times 10^{-15}$ on flat walls and $\Delta \mathbf{M} \le 3.70 \times 10^{-16}$ on corners).

For interior nodes with near-wall chord cuts ($d = 0.5 r_s$), a small discrepancy on the order of $10^{-7}$ is observed for both operators ($\Delta \mathbf{K} = 2.09 \times 10^{-7}$ and $\Delta \mathbf{M} = 1.69 \times 10^{-7}$). It is important to emphasize that this discrepancy reflects the $\mathcal{O}(h_{\mathrm{quad}})$ boundary-clipping discretization error of the numerical polar quadrature across the straight chord (where the planar wall $r(\theta) = d/\cos\theta$ is not aligned with polar coordinate lines), rather than any limitation of the analytical expressions. The exact closed-form cap area and moment formulas derived in \ref{app:truncated_geometry} completely bypass this numerical quadrature error.

\subsubsection{Boundary Condition Integrals (Neumann and Robin)}
Finally, we verify the natural boundary line integrals along straight boundary segments $\Gamma_s^I = \Omega_s^I \cap \partial\Omega$ of length $2L$ ($L = r_s$):
\begin{enumerate}[(a)]
    \item \textit{Inhomogeneous Neumann load vector} $F_a^{\mathrm{Neumann}, I} = \int_{\Gamma_s^I} \vec{g}_N(s) \cdot \vec{N}_a^I(s) \, ds$ under an affine surface current $\vec{g}_N(s)$.
    \item \textit{Robin boundary damping matrix} $C_{ab}^{\mathrm{Robin}, I} = \gamma \int_{\Gamma_s^I} (\vec{N}_a^I(s) \cdot \vec{t})(\vec{N}_b^I(s) \cdot \vec{t}) \, ds$.
\end{enumerate}
Because the shape functions $\vec{N}_a^I$ are linear in $s$, the Robin integrand is quadratic (degree 2), while the Neumann integrand with linear $\vec{g}_N$ is also of degree at most 2. Standard 1D Gauss--Legendre quadrature with $n$ points integrates polynomials up to degree $2n - 1$ exactly. As shown in Table~\ref{tab:stage1_element_verification}, a 1-point Gauss rule ($2n-1 = 1$) exhibits a truncation error of $10.1\%$ on $F^{\mathrm{Neumann}}$ and $17.6\%$ on $C^{\mathrm{Robin}}$. Conversely, a 2-point Gauss rule ($2n-1 = 3 \ge 2$) matches the closed-form analytical integrals to machine precision ($\Delta F^{\mathrm{Neumann}} \le 1.92 \times 10^{-16}$ and $\Delta C^{\mathrm{Robin}} \le 1.09 \times 10^{-16}$). Both analytical line formulas and 2-point Gauss quadrature are thus exact to floating-point precision, ensuring cell-free boundary imposition.

\subsection{Stage 2: Galerkin Patch Test for Linear Consistency}
\label{subsec:stage2_patch_test}

The Galerkin patch test was evaluated across the seven canonical affine vector states:
\begin{enumerate}
    \item Rigid Translation $X$: $\vec{E}^* = [1, 0]^T$ ($\nabla \times \vec{E}^* = 0, \nabla \cdot \vec{E}^* = 0$).
    \item Rigid Translation $Y$: $\vec{E}^* = [0, 1]^T$ ($\nabla \times \vec{E}^* = 0, \nabla \cdot \vec{E}^* = 0$).
    \item Normal Dilatation (Harmonic): $\vec{E}^* = [x, -y]^T$ ($\nabla \times \vec{E}^* = 0, \nabla \cdot \vec{E}^* = 0$).
    \item Pure Shear (Harmonic): $\vec{E}^* = [y, x]^T$ ($\nabla \times \vec{E}^* = 0, \nabla \cdot \vec{E}^* = 0$).
    \item Rigid-Body Rotation (Pure Solenoidal Curl): $\vec{E}^* = [-y, x]^T$ ($(\nabla \times \vec{E}^*)_z = 2, \nabla \cdot \vec{E}^* = 0$).
    \item Isotropic Volumetric Expansion (Pure Irrotational Divergence): $\vec{E}^* = [x, y]^T$ ($(\nabla \times \vec{E}^*)_z = 0, \nabla \cdot \vec{E}^* = 2$).
    \item General Affine Field: $\vec{E}^* = [1.5 + 0.8x - 0.3y, -0.7 + 0.4x + 1.2y]^T$ ($(\nabla \times \vec{E}^*)_z = 0.7, \nabla \cdot \vec{E}^* = 2.0$).
\end{enumerate}

In each mode, the boundary-value problem is formulated on $\Omega$ with exact essential Dirichlet conditions prescribed along the entire physical boundary: $\vec{E} \cdot \vec{t} = \vec{E}^* \cdot \vec{t}$ on $\partial\Omega$. Physically, these states represent canonical electromagnetic configurations spanning the complete 6-dimensional linear basis $\mathcal{P}^1$ (the two uniform translations plus the four mutually orthogonal derivative modes of Lemma~1):
\begin{itemize}
    \item \textit{Modes 1--2 (Uniform Fields):} Constant vector electric fields ($\vec{E}^* = [1, 0]^T$ and $[0, 1]^T$), yielding piecewise-constant tangential traces $\vec{E} \cdot \vec{t} \in \{-1, 0, 1\}$ along the boundary walls (as in a transverse uniform field or TEM mode).
    \item \textit{Modes 3--4 (Harmonic Quadrupoles):} Curl-free and divergence-free linear fields ($\nabla \times \vec{E}^* = 0, \nabla \cdot \vec{E}^* = 0$) with linear tangential traces $\vec{E} \cdot \vec{t}$ along each boundary edge (as in electrostatic quadrupole field profiles).
    \item \textit{Mode 5 (Faraday Induction):} A pure rotational solenoidal electric field ($\vec{E}^* = [-y, x]^T$, with $(\nabla \times \vec{E}^*)_z = 2, \nabla \cdot \vec{E}^* = 0$, and zero symmetric strain), where the non-zero tangential circulation along $\partial\Omega$ balances a time-varying axial magnetic flux ($\dot{B}_z = -\mathrm{const}$).
    \item \textit{Mode 6 (Space Charge):} A pure isotropic irrotational field ($\vec{E}^* = [x, y]^T$, with $\nabla \cdot \vec{E}^* = 2, \nabla \times \vec{E}^* = 0$, and zero deviatoric strain), associated with a uniform space-charge density ($\rho_0 = 2\epsilon_0$) via Gauss's law.
    \item \textit{Mode 7 (General Affine):} A complete linear superposition of uniform, quadrupole, inductive, and space-charge field components.
\end{itemize}
Because both curl and divergence of any affine field are spatial constants, the continuous volumetric source in \eqref{eq:helmholtz_pde} vanishes identically ($\vec{f}_{\mathrm{vol}} \equiv \mathbf{0}$), while the weak-form boundary fluxes produce exact, consistent discrete nodal forces $F_I$.

Simulations were performed comparing both the classical unsymmetric row-by-row MLPG assembly \eqref{eq:k_row_analytic} and the symmetric Subdomain Galerkin assembly \eqref{eq:global_assembly} on an $11 \times 11$ discretization ($N = 121$ nodes, 81 free interior DoFs) on: (a) a uniform Cartesian grid with alternating interior directors ($45^\circ / 135^\circ$), and (b) an irregular cloud with $25\%$ random nodal position perturbation ($x_k \gets x_k + \delta_x, y_k \gets y_k + \delta_y$ with $|\delta| \le 0.25 h_{\mathrm{nom}}$) and isotropic random directors. Exact Dirichlet values $e_k = \vec{E}^*(\mathbf{x}_k) \cdot \vec{t}_k$ were imposed on all boundary nodes, and the partitioned linear system $\mathbf{K}_{\mathrm{int}, \mathrm{int}} \mathbf{e}_{\mathrm{int}} = \mathbf{F}_{\mathrm{int}} - \mathbf{K}_{\mathrm{int}, \mathrm{bnd}} \mathbf{e}_{\mathrm{bnd}}$ was solved for the interior DoFs using standard sparse factorization.

\begin{table}[!t]
    \centering
    \small
    \caption{Relative infinity-norm error on free interior degrees of freedom ($\|\mathbf{e}_{\mathrm{int}} - \mathbf{e}^*\|_{\infty} / \|\mathbf{e}^*\|_{\infty}$) for the Galerkin Patch Test across 7 affine states, comparing unsymmetric row-by-row and symmetric Subdomain Galerkin assemblies.}
    \label{tab:patch_test_results}
    \resizebox{\textwidth}{!}{
        \begin{tabular}{lcccccc}
            \toprule
                                       &                                          &                                     & \multicolumn{2}{c}{\textbf{Regular Grid ($11 \times 11$)}} & \multicolumn{2}{c}{\textbf{Irregular Cloud (25\% Jitter)}}                                                   \\
            \cmidrule(lr){4-5} \cmidrule(lr){6-7}
            \textbf{Affine Field Mode} & $\boldsymbol{(\nabla\times\vec{E}^*)_z}$ & $\boldsymbol{\nabla\cdot\vec{E}^*}$ & \textbf{Unsymmetric}                                       & \textbf{Symmetric}                                         & \textbf{Unsymmetric}   & \textbf{Symmetric}     \\
            \midrule
            1. Translation $X$         & 0.0                                      & 0.0                                 & $1.41 \times 10^{-14}$                                     & $2.45 \times 10^{-14}$                                     & $1.04 \times 10^{-14}$ & $9.67 \times 10^{-14}$ \\
            2. Translation $Y$         & 0.0                                      & 0.0                                 & $1.77 \times 10^{-14}$                                     & $1.60 \times 10^{-14}$                                     & $3.80 \times 10^{-14}$ & $7.48 \times 10^{-14}$ \\
            3. Normal Dilatation       & 0.0                                      & 0.0                                 & $4.47 \times 10^{-15}$                                     & $8.47 \times 10^{-15}$                                     & $9.62 \times 10^{-15}$ & $7.20 \times 10^{-14}$ \\
            4. Pure Shear              & 0.0                                      & 0.0                                 & $1.66 \times 10^{-14}$                                     & $1.23 \times 10^{-14}$                                     & $2.40 \times 10^{-14}$ & $6.03 \times 10^{-14}$ \\
            5. Pure Rotation (Curl)    & 2.0                                      & 0.0                                 & $8.70 \times 10^{-15}$                                     & $1.26 \times 10^{-14}$                                     & $1.71 \times 10^{-14}$ & $1.91 \times 10^{-14}$ \\
            6. Volumetric Expansion    & 0.0                                      & 2.0                                 & $1.11 \times 10^{-14}$                                     & $1.21 \times 10^{-14}$                                     & $2.40 \times 10^{-14}$ & $1.29 \times 10^{-13}$ \\
            7. General Affine Field    & 0.7                                      & 2.0                                 & $2.03 \times 10^{-14}$                                     & $1.26 \times 10^{-14}$                                     & $2.31 \times 10^{-14}$ & $1.21 \times 10^{-13}$ \\
            \bottomrule
        \end{tabular}
    }
\end{table}

As reported in Table~\ref{tab:patch_test_results}, the relative error across all seven affine modes remains on the order of $10^{-14}\text{--}10^{-15}$ on regular grids and $\le 1.29 \times 10^{-13}$ on perturbed point clouds for both formulations. This demonstrates that VNMM-MLPG satisfies the Galerkin patch test to machine precision under both unsymmetric row-by-row and symmetric Subdomain Galerkin assemblies, verifying the fundamental linear consistency of the local polynomial reproduction space.

\begin{remark}[Patch Test Consistency under Orthogonal Energy Stabilization]
\label{rem:patch_test_stab}
    Activating the orthogonal energy stabilization $\mathbf{K}_{\mathrm{stab}}$ \eqref{eq:k_stab_analytic} with parameter $\alpha_{\mathrm{stab}} \in [0.1, 5.0]$ preserves consistency in two distinct ways. First, without any modification to the right-hand load vector, $\mathbf{K}_{I, \mathrm{stab}}^{\mathrm{local}}$ annihilates all uniform translation states (Modes~1 and 2, $\vec{E}^* = [a_1, a_2]^T$), pure rigid-body rotation states ($\vec{E}^* = [-y, x]^T$), and pure isotropic dilation states ($\vec{E}^* = [x, y]^T$) identically ($\mathbf{K}_{I, \mathrm{stab}}^{\mathrm{local}} \mathbf{u}^* \equiv \mathbf{0}$), keeping their relative errors below $7.0 \times 10^{-14}$ across all meshes. Second, for affine modes possessing non-zero deviatoric or symmetric shear strain (Modes~3, 4, and 7), including the consistent stabilization boundary flux $a_{\Omega_s^I}^{\mathrm{stab}}(\vec{W}_I, \vec{E}^*)$ in the local linear functional $\ell_{\Omega_s^I}$ reproduces all seven patch test states to machine precision ($\le 1.3 \times 10^{-13}$); whereas when applied without auxiliary strain boundary fluxes under the VMS scaling law $\alpha_{\mathrm{stab}}(h_{\mathrm{nom}}) = c_{\mathrm{stab}}(h_{\mathrm{nom}}/L)^2$, the subgrid stabilization residual decays quadratically as $\mathcal{O}(h_{\mathrm{nom}}^2)$, preserving full asymptotic consistency.
\end{remark}

\subsection{Stage 3: Asymptotic Refinement Convergence via MMS (Dirichlet and Neumann)}
\label{subsec:stage3_mms}

To assess asymptotic convergence rates under spatial refinement ($h_{\mathrm{nom}} \to 0$), the Method of Manufactured Solutions (MMS) was solved on a square domain $\Omega = [0, \pi]^2\text{ m}^2$ (with side length $L = \pi\text{ m}$) with exact analytical vector field:
\begin{equation}
    \vec{E}^*(x, y) = \begin{bmatrix} \sin(x) \cos(y) \\ -\cos(x) \sin(y) \end{bmatrix}, \qquad (\nabla \times \vec{E}^*)_z = 2\sin(x)\sin(y), \quad \nabla \cdot \vec{E}^* = 0.
\end{equation}
To isolate spatial discretization errors from internal cavity resonances and ensure a strictly coercive, positive-definite elliptic resolvent operator ($\mathbf{A} = \mathbf{K} + \kappa \mathbf{M} > 0$), we solve the positive-definite shifted curl-curl boundary-value problem $\nabla \times (\mu_r^{-1} \nabla \times \vec{E}) - s_{\mathrm{div}} \mu_r^{-1} \nabla(\nabla \cdot \vec{E}) + \kappa \epsilon_r \vec{E} = \vec{J}_s$ (corresponding to replacing $-k_0^2 \mapsto +\kappa$ in \eqref{eq:helmholtz_pde}) with divergence factor $s_{\mathrm{div}} = 6.0$, positive shift parameter $\kappa = +1.0\text{ m}^{-2}$ ($\mu_r = \epsilon_r = 1$), and volumetric current source $\vec{J}_s = (2 + \kappa) \vec{E}^*$. In accordance with the Overlapping Subdomain Ritz--Galerkin formulation \eqref{eq:volume_energy_form}, the local right-hand load vector on each active subdomain $\Omega_s^I \cap \Omega$ is evaluated in consistent weak form via Green's identity \eqref{eq:local_weak_form} as $f_{I, a}^{\mathrm{local}} = \ell_{\Omega_s^I}(\vec{N}_a^I) = a_{\Omega_s^I}(\vec{N}_a^I, \vec{E}^*) + \kappa m_{\Omega_s^I}(\vec{N}_a^I, \vec{E}^*)$, which incorporates both the volumetric source $\vec{J}_s$ and the consistent subdomain boundary flux along $\partial(\Omega_s^I \cap \Omega)$. Refinement was evaluated across the standard sequence of six uniform point clouds matching the benchmark in \cite{stable_vnmm2d_2026} with $N \in \{15 \times 15, 22 \times 22, 29 \times 29, 36 \times 36, 43 \times 43, 49 \times 49\}$, corresponding to nominal grid spacings $h_{\mathrm{nom}} = \pi/(N_x-1) \in [0.2244, 0.0654]\text{ m}$. In all simulations, isotropic random vector directors ($\text{seed}=42$) were employed to ensure that spatial convergence is strictly unbiased by Cartesian lattice alignments.

To systematically test the boundary formulation developed in Section~\ref{subsec:boundary_treatment}, two distinct boundary-value problem configurations were investigated:
\begin{enumerate}
    \item \textbf{Case A (Pure Inhomogeneous Dirichlet):} On all four boundary walls $\partial\Omega$, the tangential trace is prescribed by the exact manufactured field $e_k = \vec{E}^*(\mathbf{x}_k) \cdot \vec{t}_k$, and boundary nodes are eliminated via condensed system \eqref{eq:dirichlet_condensed_system}.
    \item \textbf{Case B (Mixed Dirichlet--Neumann / PMC):} Inhomogeneous Dirichlet conditions are imposed on the vertical walls $x \in \{0, \pi\}$, while the horizontal walls $y \in \{0, \pi\}$ are treated as homogeneous natural Neumann (PMC) boundaries. Because $(\nabla \times \vec{E}^*)_z = 2\sin(x)\sin(y)$ vanishes identically along $y=0$ and $y=\pi$, the exact electromagnetic boundary flux satisfies $\vec{\tau} = \frac{1}{\mu_r} (\vec{n} \times (\nabla \times \vec{E}^*)) \equiv \mathbf{0}$. Therefore, the boundary nodes along $y \in \{0, \pi\}$ are retained as active, unconstrained degrees of freedom and determined directly by the global system.
\end{enumerate}

\begin{table}[!t]
    \centering
    \small
    \caption{MMS refinement convergence comparison between Pure Dirichlet and Mixed Dirichlet--Neumann (PMC) boundary conditions for VNMM-MLPG 2D ($s_{\mathrm{div}} = 6.0, \alpha_s = 0.80$, isotropic random directors) across the 6-mesh sequence from \cite{stable_vnmm2d_2026}.}
    \label{tab:mms_convergence}
    \resizebox{\textwidth}{!}{
        \begin{tabular}{cc|cccc|cccc}
            \toprule
                                                                 &                                              & \multicolumn{4}{c|}{\textbf{Case A: Pure Dirichlet}} & \multicolumn{4}{c}{\textbf{Case B: Mixed Dirichlet--Neumann (PMC)}}                                                                                                                                                                                                                                                                   \\
            \textbf{Mesh}                                        & $\boldsymbol{h_{\mathrm{nom}}}$ \textbf{[m]} & \textbf{Active DoFs}                                 & $\boldsymbol{L^2(\vec{E})}$                                         & $\boldsymbol{|H(\mathbf{curl})|}$                   & \textbf{DoF Error} & \textbf{Active DoFs}                                & $\boldsymbol{L^2(\vec{E})}$                         & $\boldsymbol{|H(\mathbf{curl})|}$                   & \textbf{DoF Error} \\
            \midrule
            $15 \times 15$                                       & 0.2244                                       & 169                                                  & $6.529\%$                                                           & $21.817\%$                                          & $8.383\%$          & 197                                                 & $7.330\%$                                           & $20.865\%$                                          & $9.599\%$          \\
            $22 \times 22$                                       & 0.1496                                       & 400                                                  & $3.398\%$                                                           & $13.991\%$                                          & $5.555\%$          & 442                                                 & $3.943\%$                                           & $13.661\%$                                          & $13.563\%$         \\
            $29 \times 29$                                       & 0.1122                                       & 729                                                  & $1.922\%$                                                           & $10.556\%$                                          & $2.407\%$          & 785                                                 & $2.187\%$                                           & $10.282\%$                                          & $3.545\%$          \\
            $36 \times 36$                                       & 0.0898                                       & 1156                                                 & $1.184\%$                                                           & $8.336\%$                                           & $1.721\%$          & 1226                                                & $1.352\%$                                           & $8.250\%$                                           & $3.141\%$          \\
            $43 \times 43$                                       & 0.0748                                       & 1681                                                 & $0.804\%$                                                           & $6.835\%$                                           & $1.129\%$          & 1765                                                & $1.022\%$                                           & $6.741\%$                                           & $3.242\%$          \\
            $49 \times 49$                                       & 0.0654                                       & 2209                                                 & $0.637\%$                                                           & $5.941\%$                                           & $1.010\%$          & 2305                                                & $0.780\%$                                           & $5.892\%$                                           & $1.678\%$          \\
            \midrule
            \multicolumn{2}{c|}{\textbf{Fitted Asymptotic Rate}} & ---                                          & $\boldsymbol{\mathcal{O}(h_{\mathrm{nom}}^{1.93})}$  & $\boldsymbol{\mathcal{O}(h_{\mathrm{nom}}^{1.05})}$                 & $\boldsymbol{\mathcal{O}(h_{\mathrm{nom}}^{1.84})}$ & ---                & $\boldsymbol{\mathcal{O}(h_{\mathrm{nom}}^{1.85})}$ & $\boldsymbol{\mathcal{O}(h_{\mathrm{nom}}^{1.02})}$ & $\boldsymbol{\mathcal{O}(h_{\mathrm{nom}}^{1.49})}$                      \\
            \multicolumn{2}{c|}{\textbf{Theoretical Rate}}       & ---                                          & $\boldsymbol{\mathcal{O}(h_{\mathrm{nom}}^{2.00})}$  & $\boldsymbol{\mathcal{O}(h_{\mathrm{nom}}^{1.00})}$                 & ---                                                 & ---                & $\boldsymbol{\mathcal{O}(h_{\mathrm{nom}}^{2.00})}$ & $\boldsymbol{\mathcal{O}(h_{\mathrm{nom}}^{1.00})}$ & ---                                                                      \\
            \bottomrule
        \end{tabular}
    }
\end{table}

\begin{figure}[!t]
    \centering
    \includegraphics[width=0.78\textwidth]{fig_mlpg_mms_convergence.png}
    \caption{Asymptotic spatial convergence for the VNMM-MLPG under MMS with Pure Dirichlet and Mixed Dirichlet--Neumann (PMC) boundary conditions across the six meshes in \cite{stable_vnmm2d_2026} with isotropic random directors. Both boundary formulations achieve the theoretical polynomial convergence rates: strictly monotonic quadratic $\mathcal{O}(h_{\mathrm{nom}}^2)$ in the $L^2(\vec{E})$ field norm and linear $\mathcal{O}(h_{\mathrm{nom}}^1)$ in the $H(\mathbf{curl})$ rotational semi-norm $|\cdot|_{H(\mathbf{curl})}$, demonstrating the numerical stability and accuracy of natural Neumann boundary treatment.}
    \label{fig:mms_convergence}
\end{figure}

The numerical results presented in Table~\ref{tab:mms_convergence} and Fig.~\ref{fig:mms_convergence} confirm several key findings:
\begin{enumerate}
    \item \textbf{Field Convergence Rates:} In both Case A (Dirichlet) and Case B (Mixed Neumann), the observed convergence rate in the continuous $L^2(\vec{E})$ field norm is strictly monotonic across all six refinement levels, achieving asymptotic rates of $\mathcal{O}(h_{\mathrm{nom}}^{1.93})$ and $\mathcal{O}(h_{\mathrm{nom}}^{1.85})$, respectively, in close agreement with the theoretical quadratic rate $\mathcal{O}(h_{\mathrm{nom}}^2)$ for linear vector polynomials ($p=1$). Note that in Table~\ref{tab:mms_convergence}, the auxiliary ``DoF Error'' column reports the relative discrete $\ell_2$ error $\|\mathbf{e} - \mathbf{e}^*\|_2 / \|\mathbf{e}^*\|_2$ evaluated across the active directional scalar unknowns $e_k = \vec{E}^h(\mathbf{x}_k) \cdot \vec{t}_k$; while Case~B exhibits a localized fluctuation in this pointwise directional metric at $22 \times 22$ ($13.563\%$) due to random interior director orientations ($\vec{t}_k$) near the unconstrained Neumann boundary, the physically objective continuous functional norms $L^2(\vec{E})$ ($3.943\%$) and $|H(\mathbf{curl})|$ ($13.661\%$) decrease strictly monotonically across every refinement step.
    \item \textbf{Rotational Convergence Rates:} In the $H(\mathbf{curl})$ semi-norm $|\vec{E}^h - \vec{E}^*|_{H(\mathbf{curl})}$, both formulations achieve monotonic rates of $\mathcal{O}(h_{\mathrm{nom}}^{1.05})$ and $\mathcal{O}(h_{\mathrm{nom}}^{1.02})$, in exact agreement with the theoretical linear rate $\mathcal{O}(h_{\mathrm{nom}})$.
    \item \textbf{Boundary Error Equivalence:} Releasing boundary nodes on the Neumann walls $y \in \{0, \pi\}$ introduces no numerical degradation or boundary layer oscillations. In fact, across all grid levels, the $H(\mathbf{curl})$ semi-norm error with active Neumann nodes is marginally lower than with fully clamped Dirichlet nodes (e.g., $20.865\%$ vs. $21.817\%$ on $15 \times 15$ and $5.892\%$ vs. $5.941\%$ on $49 \times 49$), demonstrating that the natural variational weak form relaxes artificial Dirichlet over-constraining.
    \item \textbf{Solvability without Penalty Parameters:} The mixed Dirichlet--Neumann problem was solved without adding any boundary penalty terms or boundary element integrations, validating the theoretical analysis in Section~\ref{subsec:boundary_treatment}.
    \item \textbf{Asymptotic Convergence Invariance under Orthogonal Energy Stabilization:} Solving the MMS with the orthogonal energy stabilization operator $\mathbf{K}_{\mathrm{stab}}$ \eqref{eq:k_stab_analytic} parameterized by the dimensionless subgrid law $\alpha_{\mathrm{stab}}(h_{\mathrm{nom}}) = c_{\mathrm{stab}} (h_{\mathrm{nom}}/L)^2$ ($c_{\mathrm{stab}} \in [10, 100]$) preserves the asymptotic convergence rate without loss of accuracy. Because $\alpha_{\mathrm{stab}}(h_{\mathrm{nom}})$ decays quadratically with $h_{\mathrm{nom}}$, the stabilization energy scales at the same rate as the $\mathcal{P}^1$ interpolation truncation error. As a result, the asymptotic rates of $\mathcal{O}(h_{\mathrm{nom}}^{1.92})$ in $L^2(\vec{E})$ and $\mathcal{O}(h_{\mathrm{nom}}^{1.02})$ in the $H(\mathbf{curl})$ semi-norm remain strictly unaltered, confirming that stabilizing the rank-deficient shear/deviatoric directions does not induce artificial locking in source-driven boundary-value problems.
\end{enumerate}

\subsubsection{Comparison of Symmetric Subdomain Galerkin and Unsymmetric Row Assemblies}
\label{subsubsec:mms_assembly_comparison}

To empirically validate the matrix symmetry and stability analysis developed in Sections~\ref{subsubsec:row_mlpg_symmetry} and \ref{subsubsec:symmetric_assembly}, the manufactured solution benchmark was solved under both assembly strategies:
\begin{enumerate}
    \item \textbf{Symmetric Subdomain Galerkin Assembly \eqref{eq:global_assembly}:} Within each disk $\Omega_s^I$, testing is carried out against all six local shape functions $\vec{N}_a^I(\mathbf{x})$, yielding symmetric positive-semidefinite $6 \times 6$ local blocks whose sum produces a strictly symmetric, positive-definite global operator $\mathbf{A}_{\mathrm{sym}} = \mathbf{A}_{\mathrm{sym}}^T > 0$.
    \item \textbf{Classical Unsymmetric Row-by-Row Assembly \eqref{eq:k_row_analytic}:} Each disk $\Omega_s^I$ is tested with only the single vector shape function assigned to its evaluation node ($\vec{W}_I(\mathbf{x}) = \vec{N}_{i_0}^I(\mathbf{x})$), generating an intrinsically unsymmetric matrix $\mathbf{A}_{\mathrm{unsym}} \ne \mathbf{A}_{\mathrm{unsym}}^T$.
\end{enumerate}

Table~\ref{tab:mms_assembly_comparison} reports the comparative performance of both formulations across three coarse-to-medium refinement levels ($9 \times 9$, $13 \times 13$, and $17 \times 17$; coarser meshes are shown because the unsymmetric row-by-row system exhibits severe numerical divergence by $17 \times 17$, exceeding $3300\%$ error in $L^2(\vec{E})$ and $7800\%$ in the $H(\mathbf{curl})$ semi-norm).

\begin{table}[!t]
    \centering
    \small
    \caption{MMS performance comparison between the Symmetric Subdomain Galerkin assembly and the Unsymmetric Row-by-Row Petrov-Galerkin assembly ($s_{\mathrm{div}} = 6.0, \alpha_s = 0.80$, Pure Dirichlet conditions).}
    \label{tab:mms_assembly_comparison}
    \resizebox{\textwidth}{!}{
        \begin{tabular}{cccccc}
            \toprule
            \textbf{Mesh}  & \textbf{Formulation}           & $\boldsymbol{L^2(\vec{E})}$ \textbf{Error} & $\boldsymbol{|H(\mathbf{curl})|}$ \textbf{Error} & $\mathrm{cond}(\mathbf{A}_{\mathrm{int}, \mathrm{int}})$ & $\min \mathrm{Re}(\lambda)$ \\
            \midrule
            $9 \times 9$   & Symmetric (Subdomain Galerkin) & $\mathbf{12.31\%}$                         & $\mathbf{47.86\%}$                             & $1.85 \times 10^2$                                       & $+1.64$                     \\
            $9 \times 9$   & Unsymmetric (Row Petrov-Gal.)  & $75.09\%$                                  & $104.31\%$                                     & $1.82 \times 10^2$                                       & $-14.19$ (10 neg.)          \\
            \midrule
            $13 \times 13$ & Symmetric (Subdomain Galerkin) & $\mathbf{7.97\%}$                          & $\mathbf{29.13\%}$                             & $8.25 \times 10^2$                                       & $+0.36$                     \\
            $13 \times 13$ & Unsymmetric (Row Petrov-Gal.)  & $895.71\%$                                 & $2063.16\%$                                    & $2.80 \times 10^4$                                       & $-13.25$ (26 neg.)          \\
            \midrule
            $17 \times 17$ & Symmetric (Subdomain Galerkin) & $\mathbf{3.10\%}$                          & $\mathbf{19.17\%}$                             & $1.26 \times 10^3$                                       & $+0.24$                     \\
            $17 \times 17$ & Unsymmetric (Row Petrov-Gal.)  & $3310.09\%$                                & $7870.43\%$                                    & $2.18 \times 10^5$                                       & $-12.73$ (48 neg.)          \\
            \bottomrule
        \end{tabular}
    }
\end{table}

The comparative results highlight a fundamental theoretical and practical distinction:
\begin{itemize}
    \item \textbf{Patch Test vs. MMS General PDE Solutions:} While both formulations pass the linear patch test to machine precision (Table~\ref{tab:patch_test_results}) because affine fields have zero local residual in $\mathcal{P}^1$, their behavior diverges sharply on general non-polynomial fields.
    \item \textbf{Loss of Discrete Coercivity in Row Assembly:} Under classical row-by-row assembly, testing with only one vector function per disk fails to constrain all 6 degrees of freedom within the local trial stencil $\Omega_u^I$. As a consequence, the reduced interior matrix $\mathbf{A}_{\mathrm{int}, \mathrm{int}}$ becomes indefinite: for the $9 \times 9$ mesh, it possesses 10 negative real eigenvalues and 26 complex eigenvalues. Upon grid refinement, this inf-sup instability causes the condition number to escalate rapidly to $2.18 \times 10^5$, causing severe numerical divergence.
    \item \textbf{Discrete Coercivity of Subdomain Galerkin Assembly:} In contrast, the symmetric Subdomain Galerkin formulation tests against the complete local polynomial space on each disk $\Omega_s^I$. This guarantees that the local blocks are symmetric positive-semidefinite and their global sum is strictly positive-definite ($\min \lambda_i > 0$). Consequently, the linear system remains well-conditioned, discrete coercivity is preserved, and the method achieves stable, monotonic $\mathcal{O}(h_{\mathrm{nom}}^2)$ convergence.
\end{itemize}

\subsection{Stage 4: Subdomain Scaling Bounds and Parametric Sensitivity of $\alpha_s$}
\label{subsec:stage4_alpha_s}

The local integration subdomain radius is defined in \eqref{eq:subdomain_disk} as $r_s^I = \alpha_s h_I$. In the classical scalar MLPG and LPIM literature \cite{atluri2002meshless, liu2009mesh}, $\alpha_s$ is typically selected in the interval $[0.75, 0.85]$ (commonly $\alpha_s = 0.80$) to balance domain coverage against subdomain overlap. To verify that this standard carries over to 6-node directional vector stencils ($\mathcal{P}^1$) on both regular and $25\%$ perturbed clouds, we examine the two geometric limits governing $\alpha_s$:
\begin{enumerate}
    \item \textbf{Lower Bound (Global Domain Coverage \cite{atluri2002meshless, liu2009mesh}):} Spatial integration in MLPG relies on the collection of local subdomains forming a complete cover of the physical domain, $\bigcup_{I=1}^N \Omega_s^I \supseteq \Omega$. In a quasi-uniform nodal distribution with nominal grid spacing $h_{\mathrm{nom}}$, the maximum distance from any point in a Cartesian cell to its nearest node (the cell circumradius) is:
          \begin{equation}
              r_{\mathrm{cov}} = \frac{\sqrt{2}}{2} h_{\mathrm{nom}} \approx 0.7071 h_{\mathrm{nom}}.
              \label{eq:alpha_s_lower_bound}
          \end{equation}
          Because the circular integration disk radius is defined locally as $r_s^I = \alpha_s h_{\mathrm{local}}(\mathbf{x}_I)$ (where $h_{\mathrm{local}}(\mathbf{x}_I) \equiv h_I = \min_{J \ne I} \|\mathbf{x}_J - \mathbf{x}_I\|$ equals $h_{\mathrm{nom}}$ on regular lattices), global domain coverage requires $r_s \ge r_{\mathrm{cov}}$, which yields the classical threshold $\alpha_s \ge \sqrt{2}/2 \approx 0.7071$ \cite{atluri2002meshless, liu2009mesh}. If $\alpha_s < \sqrt{2}/2$, unintegrated voids or interstitial gaps emerge between adjacent circular subdomains. As demonstrated in Fig.~\ref{fig:alpha_s_sensitivity}(a), on a regular lattice ($0\%$ jitter), full coverage ($100\%$) is reached at $\alpha_s \ge 0.75$, whereas below $0.7071$ coverage drops precipitously to $95.7\%$ at $\alpha_s = 0.60$, $78.9\%$ at $\alpha_s = 0.50$, and $29.4\%$ at $\alpha_s = 0.30$. Under nodal perturbations ($25\%$ jitter), localized pairwise node clustering shrinks the minimum nearest-neighbor distance $h_I$ for closely paired nodes while adjacent regions dilate: at $\alpha_s = 0.80$, coverage is $94.4\%$ for $25\%$ jitter, reaching asymptotic saturation ($\ge 95\%$) for $\alpha_s \ge 0.85\text{--}1.00$ (which can be further smoothed toward $100\%$ if $h_I$ is alternatively defined as the stencil-averaged neighbor distance).
    \item \textbf{Upper Bound (6-Node Vector Stencil Containment):} Each node $\mathbf{x}_I$ constructs its VNMM trial vector field $\vec{E}^h$ over a 6-node directional stencil $\Omega_u^I$ whose spatial support radius is $r_u^I = \max_{k \in \Omega_u^I} \|\mathbf{x}_k - \mathbf{x}_I\|$. For a regular or perturbed Cartesian grid, a 6-node stencil includes nearest neighbors at distance $h_{\mathrm{local}} \approx h_{\mathrm{nom}}$ and diagonal neighbors at distance $\sqrt{2}h_{\mathrm{nom}}$, giving an average support radius $r_u \approx 1.15\text{--}1.20 h_{\mathrm{nom}}$. If $\alpha_s$ is enlarged beyond this limit ($r_s^I > r_u^I$), the integration disk $\Omega_s^I$ spills outside the convex hull of its generating stencil. In this extrapolation regime, polynomial shape functions $\vec{W}_k^I(\mathbf{x})$ grow rapidly outside their interpolating cloud, causing numerical degradation. Geometric analysis confirms that support domain violations ($r_s^I > r_u^I$) remain strictly zero for all $\alpha_s \le 1.10$ across both regular and perturbed clouds ($0\%$ and $25\%$ jitter), initiating only when $\alpha_s \approx 1.15\text{--}1.20$ and growing systematically for larger subdomains.
\end{enumerate}

\begin{figure}[!t]
    \centering
    \includegraphics[width=\textwidth]{fig_mlpg_alpha_s_sensitivity.png}
    \caption{Parametric sensitivity of the subdomain radius factor $\alpha_s$ ($21\times 21$ mesh, random directors): (a) Domain integration coverage percentage $\bigcup \Omega_s^I \supseteq \Omega$; (b) Relative $L^2(\vec{E})$ error in the 2D MMS problem across regular and perturbed clouds.}
    \label{fig:alpha_s_sensitivity}
\end{figure}

To quantify these geometric bounds, Fig.~\ref{fig:alpha_s_sensitivity}(b) illustrates the boundary-value convergence under the Method of Manufactured Solutions on the shared $21 \times 21$ baseline grid with randomized vector directors across regular and perturbed clouds:
\begin{itemize}
    \item \textbf{MMS Convergence Accuracy:} For the MMS benchmark ($21\times 21$ grid, random directors), the relative $L^2(\vec{E})$ error exhibits a smooth, stable response across the operational interval $\alpha_s \in [0.70, 1.20]$: at $\alpha_s = 0.80$, the error is $2.94\%$ for $0\%$ jitter and $4.26\%$ for $25\%$ jitter. Beyond $\alpha_s > 1.20$, extrapolation errors cause accuracy to degrade (reaching $4.71\%$ and $5.46\%$ at $\alpha_s = 2.00$, respectively). Notably, the apparent flattening of error below $\alpha_s < 0.70$ does not indicate enhanced physical accuracy, but rather an unphysical plateau caused by extensive unintegrated voids (where up to $70\%\text{--}85\%$ of the domain is unrepresented by any local test subdomain).
\end{itemize}

Based on this comprehensive study across both regular and disordered clouds:
\begin{itemize}
    \item $\alpha_s < 0.70$ suffers from severe under-integration voids, failing to capture the global field dynamics across the physical domain.
    \item $\alpha_s > 1.20$ triggers polynomial extrapolation beyond stencil support radii, leading to systematic error amplification.
    \item $\alpha_s = 0.80$ provides a well-balanced trade-off: on regular lattices, it ensures complete ($100\%$) domain coverage with a $13\%$ safety buffer above $\sqrt{2}/2$, maintains high coverage on perturbed clouds ($94.4\%$ at $25\%$ jitter), and stays comfortably below the extrapolation boundary ($r_s^I \approx 0.72 r_u^I < r_u^I$) with exactly $0\%$ containment violations.
\end{itemize}
In addition, choosing $\alpha_s = 0.80$ mirrors the scale-invariant stencil qualification threshold $\widehat{\mathrm{Tol}} = 0.80$ established in Section~\ref{subsec:stencil_selection}. Adopting $\alpha_s = 0.80$ as the single, universal operating standard throughout all simulations in this paper establishes parameter economy, eliminating problem-dependent empirical tuning while guaranteeing domain coverage, strict support containment, and accuracy across regular and perturbed nodal distributions.

\begin{remark}[Geometric Invariance of $\alpha_s$ under Energy Stabilization]
    The bounding limits $\alpha_s \ge \sqrt{2}/2 \approx 0.7071$ (domain integration coverage) and $\alpha_s \le 1.15$ (support domain containment) are purely geometric invariants of the local partition and trial stencils, completely independent of algebraic operator modifications. Repeating the parametric MMS sweep of Fig.~\ref{fig:alpha_s_sensitivity}(b) with the orthogonal energy stabilization operator $\mathbf{K}_{\mathrm{stab}}$ activated ($\alpha_{\mathrm{stab}} \propto (h_{\mathrm{nom}}/L)^2$, Section~\ref{subsec:rank_deficiency_stabilization}) reveals an identical flat error profile across $\alpha_s \in [0.70, 1.20]$ (with variations below $0.3\%$). Because algebraic stabilization acts strictly within existing test disks and cannot integrate unrepresented physical voids, the selection $\alpha_s = 0.80$ remains a robust, canonical geometric foundation whether energy stabilization is active or inactive.
\end{remark}

\subsection{Stage 5: Physical Spectral Validation and Divergence Penalty Scaling}
\label{subsec:stage5_cavity}

Having established $\alpha_s = 0.80$ as the universal subdomain standard in Section~\ref{subsec:stage4_alpha_s}, the physical spectral fidelity of VNMM-MLPG is evaluated on the 2D PEC resonant cavity $\Omega = [0, \pi]^2\text{ m}^2$ ($L_x = L_y = \pi\text{ m}$) with $\mu_r = 1, \epsilon_r = 1$. The analytical cutoff wavenumbers for transverse electric modes $TE_{nm}$ evaluate to:
\begin{equation}
    k_{c, nm} = \sqrt{\lambda_{nm}} = \sqrt{n^2 + m^2}\text{ m}^{-1}, \qquad n, m \in \mathbb{N}_0, \quad (n, m) \ne (0, 0).
    \label{eq:analytical_cavity_kc}
\end{equation}
Global stiffness and mass operators $\mathbf{K}_{\mathrm{red}}$ and $\mathbf{M}_{\mathrm{red}}$ are assembled via the symmetric Subdomain Galerkin scheme \eqref{eq:global_assembly}, guaranteeing strictly real eigenvalues ($\lambda = k_c^2 \in \mathbb{R}$). Physical solenoidal modes are distinguished from spurious gradient modes via the divergence energy ratio matching \cite{stable_vnmm2d_2026}:
\begin{equation}
    R_{\mathrm{div}}(\mathbf{e}) = \frac{\mathbf{e}^T \mathbf{K}_{\mathrm{div}} \mathbf{e}}{\mathbf{e}^T \mathbf{K}_{\mathrm{curl}} \mathbf{e}} \times 100\%.
    \label{eq:R_div_definition}
\end{equation}
Physical solenoidal modes satisfy $R_{\mathrm{div}} \le 50\%$, whereas unregularized gradient modes exhibit $R_{\mathrm{div}} \gg 100\%$. In all simulations below, isotropic random vector directors ($\text{seed}=42$) and $\alpha_s = 0.80$ are employed.

\subsubsection{Parametric Analysis of $s_{\mathrm{div}}$ on the Baseline $21\times 21$ Grid}
\label{subsubsec:s_div_sweep}

We first investigate the sensitivity of the divergence regularization penalty $s_{\mathrm{div}}$ on the baseline $21 \times 21$ grid ($N = 441$ nodes, 361 active free interior DoFs, nominal spacing $h_{\mathrm{nom}} = \pi/20 \approx 0.1571\text{ m}$). In previous cell-based formulations \cite{stable_vnmm2d_2026}, a penalty of $s_{\mathrm{div}} = 6.0$ was adopted as standard. To systematically assess its efficacy in the cell-free MLPG setting, a parametric sweep across $s_{\mathrm{div}} \in [0.0, 50.0]$ was performed, as documented in Table~\ref{tab:s_div_sweep}.

In computational electromagnetics, non-physical spurious gradient modes originating from discrete nullspace leakage ($\nabla_h \times \nabla_h \phi_h \neq \mathbf{0}$) scale linearly with the divergence penalty ($d\lambda_{\mathrm{spur}} / ds_{\mathrm{div}} \gg 0$) and are shifted rapidly toward high frequencies. Conversely, physical cavity resonances stabilize asymptotically near their analytical cutoff values ($d\lambda / ds_{\mathrm{div}} \to 0$).

As reported in Table~\ref{tab:s_div_sweep}:
\begin{itemize}
    \item At $s_{\mathrm{div}} = 0.0$ (unregularized), the stiffness matrix $\mathbf{K}_{\mathrm{curl}}$ is numerically singular; all initial 10 numerical modes are spurious gradient fields with dominant divergence energy ($\max R_{\mathrm{div}} > 10^4\%$), giving an apparent error of $83.23\%$.
    \item For $s_{\mathrm{div}} \ge 4.0$, all 10 computed states correspond strictly to physical cavity modes, with the lowest spurious mode shifted completely outside the operational band ($\lambda_{\mathrm{spur}} > 78$).
    \item Adopting the previous standard $s_{\mathrm{div}} = 6.0$ from \cite{stable_vnmm2d_2026} yields a mean relative error of $8.83\%$.
    \item Increasing the penalty to $s_{\mathrm{div}} = 15.0$ substantially reduces the divergence residue of physical modes (down to $\max R_{\mathrm{div}} = 3.4\%$), yielding a minimum mean error of $\mathbf{2.46\%}$ across the first 10 physical modes, with max error below $4.95\%$ and stable conditioning ($\mathrm{cond} = 5.14 \times 10^4$). Note that the slight non-monotonicity in the ``Lowest Spur $\lambda$'' column between $s_{\mathrm{div}} = 10.0$ ($\lambda = 252.60$) and $s_{\mathrm{div}} = 15.0$ ($\lambda = 237.24$) occurs far above the operational band ($\lambda \le 9$, near the ~80th discrete eigenmode approaching the grid Nyquist limit), where ascending gradient modes undergo spectral mode veering (hybridization) with high-order dispersive cavity modes across the fixed classification threshold $R_{\mathrm{div}} = 50\%$.
    \item For $s_{\mathrm{div}} \ge 30.0$, over-penalization begins to slightly elevate the mean error ($4.83\%$ at $s_{\mathrm{div}} = 30.0$ and $8.65\%$ at $s_{\mathrm{div}} = 50.0$).
\end{itemize}
Consequently, for this $21 \times 21$ discretization, the most effective penalty calibration is $s_{\mathrm{div}} = 15.0$, significantly outperforming the value $s_{\mathrm{div}} = 6.0$ adopted in \cite{stable_vnmm2d_2026}.

\begin{table}[!t]
    \centering
    \small
    \caption{Parametric sensitivity of the div-curl penalty parameter $s_{\mathrm{div}}$ on spurious mode shifting and cutoff accuracy ($21 \times 21$ grid, $\alpha_s = 0.80$, isotropic random directors, criterion $R_{\mathrm{div}} > 50\%$). The recommended calibration $s_{\mathrm{div}} = 15.0$ and the previous cell-based standard $s_{\mathrm{div}} = 6.0$ are highlighted.}
    \label{tab:s_div_sweep}
    \resizebox{\textwidth}{!}{
        \begin{tabular}{ccccccc}
            \toprule
            $\boldsymbol{s_{\mathrm{div}}}$ & \textbf{Spur in Band} & \textbf{Lowest Spur} $\boldsymbol{\lambda}$ & \textbf{Mean Error} $\boldsymbol{k_c}$ \textbf{(\%)} & \textbf{Max Error} $\boldsymbol{k_c}$ \textbf{(\%)} & $\boldsymbol{\mathrm{cond}(\mathbf{K}_{\mathrm{red}})}$ & \textbf{Max} $\boldsymbol{R_{\mathrm{div}}}$ \textbf{(\%)} \\
            \midrule
            0.0                             & 10                    & 0.05                                        & $83.23\%$                                            & $87.39\%$                                           & Singular                                                & $> 10^4\%$                                                 \\
            1.0                             & 3                     & 1.23                                        & $27.82\%$                                            & $40.83\%$                                           & $9.02 \times 10^3$                                      & $466.0\%$                                                  \\
            2.0                             & 3                     & 2.17                                        & $19.00\%$                                            & $28.74\%$                                           & $9.67 \times 10^3$                                      & $180.9\%$                                                  \\
            4.0                             & 0                     & 78.39                                       & $12.19\%$                                            & $25.31\%$                                           & $1.58 \times 10^4$                                      & $41.5\%$                                                   \\
            6.0 (\cite{stable_vnmm2d_2026}) & 0                     & 173.66                                      & $8.83\%$                                             & $21.40\%$                                           & $2.25 \times 10^4$                                      & $30.2\%$                                                   \\
            8.0                             & 0                     & 207.15                                      & $6.56\%$                                             & $13.70\%$                                           & $2.91 \times 10^4$                                      & $21.9\%$                                                   \\
            10.0                            & 0                     & 252.60                                      & $4.82\%$                                             & $7.51\%$                                            & $3.56 \times 10^4$                                      & $15.1\%$                                                   \\
            \textbf{15.0 (Recommended)}     & \textbf{0}            & \textbf{237.24}                             & \textbf{2.46\%}                                      & \textbf{4.95\%}                                     & $\mathbf{5.14 \times 10^4}$                             & \textbf{3.4\%}                                             \\
            20.0                            & 0                     & 471.77                                      & $2.85\%$                                             & $5.81\%$                                            & $6.69 \times 10^4$                                      & $2.2\%$                                                    \\
            30.0                            & 0                     & 693.46                                      & $4.83\%$                                             & $10.82\%$                                           & $9.68 \times 10^4$                                      & $1.0\%$                                                    \\
            50.0                            & 0                     & 1035.15                                     & $8.65\%$                                             & $18.03\%$                                           & $1.54 \times 10^5$                                      & $0.8\%$                                                    \\
            \bottomrule
        \end{tabular}
    }
\end{table}

\subsubsection{Cavity Spectrum and Mode-by-Mode Accuracy at $s_{\mathrm{div}} = 15.0$}
\label{subsubsec:cavity_modes_15}

Using the recommended penalty $s_{\mathrm{div}} = 15.0$, Table~\ref{tab:cavity_eigenvalues} presents the spectrum for the first 10 numerical transverse electric modes on the $21 \times 21$ mesh, compared against the exact analytical benchmark from Eq.~\eqref{eq:analytical_cavity_kc}.

\begin{table}[!t]
    \centering
    \small
    \caption{Comparison of the first 10 numerical cavity eigenvalues ($\lambda$ in $\mathrm{m}^{-2}$) and cutoff wavenumbers ($k_c$ in $\mathrm{m}^{-1}$) computed by VNMM-MLPG ($21 \times 21, s_{\mathrm{div}} = 15.0, \alpha_s = 0.80$, isotropic random directors) against analytical reference solutions from Eq.~\eqref{eq:analytical_cavity_kc} on a cavity with dimensions $\pi\text{ m} \times \pi\text{ m}$.}
    \label{tab:cavity_eigenvalues}
    \resizebox{\textwidth}{!}{
        \begin{tabular}{ccccccc}
            \toprule
            \textbf{Mode}                                                                    & $\boldsymbol{\lambda_{\mathrm{num}}}$ \textbf{[m$^{-2}$]} & $\boldsymbol{\lambda_{\mathrm{ref}}}$ \textbf{[m$^{-2}$]} & $\boldsymbol{k_{c, \mathrm{num}}}$ \textbf{[m$^{-1}$]} & $\boldsymbol{k_{c, \mathrm{ref}}}$ \textbf{[m$^{-1}$]} & \textbf{Error} $\boldsymbol{k_c}$ \textbf{(\%)} & $\boldsymbol{R_{\mathrm{div}}}$ \textbf{(\%)} \\
            \midrule
            $TE_{10}$                                                                        & 0.9035                                                    & 1.0000                                                    & 0.9505                                                 & 1.0000                                                 & \textbf{4.95\%}                                 & $0.63\%$                                      \\
            $TE_{01}$                                                                        & 0.9213                                                    & 1.0000                                                    & 0.9598                                                 & 1.0000                                                 & \textbf{4.02\%}                                 & $0.53\%$                                      \\
            $TE_{11}$                                                                        & 1.9366                                                    & 2.0000                                                    & 1.3916                                                 & 1.4142                                                 & \textbf{1.60\%}                                 & $0.67\%$                                      \\
            $TE_{20}$                                                                        & 3.7689                                                    & 4.0000                                                    & 1.9414                                                 & 2.0000                                                 & \textbf{2.93\%}                                 & $1.19\%$                                      \\
            $TE_{02}$                                                                        & 3.9924                                                    & 4.0000                                                    & 1.9981                                                 & 2.0000                                                 & \textbf{0.10\%}                                 & $0.91\%$                                      \\
            $TE_{21}$                                                                        & 4.6468                                                    & 5.0000                                                    & 2.1556                                                 & 2.2361                                                 & \textbf{3.60\%}                                 & $1.05\%$                                      \\
            $TE_{12}$                                                                        & 5.3363                                                    & 5.0000                                                    & 2.3101                                                 & 2.2361                                                 & \textbf{3.31\%}                                 & $1.19\%$                                      \\
            $TE_{22}$                                                                        & 8.2883                                                    & 8.0000                                                    & 2.8789                                                 & 2.8284                                                 & \textbf{1.79\%}                                 & $1.82\%$                                      \\
            $TE_{30}$                                                                        & 8.6164                                                    & 9.0000                                                    & 2.9354                                                 & 3.0000                                                 & \textbf{2.15\%}                                 & $3.43\%$                                      \\
            $TE_{03}$                                                                        & 9.0222                                                    & 9.0000                                                    & 3.0037                                                 & 3.0000                                                 & \textbf{0.12\%}                                 & $3.20\%$                                      \\
            \midrule
            \multicolumn{5}{l}{\textbf{Mean Relative Error across First 10 Physical Modes:}} & \textbf{2.46\%}                                           & ---                                                                                                                                                                                                                                                                           \\
            \multicolumn{5}{l}{\textbf{Spurious Modes in Operational Band ($\le 10$):}}      & \textbf{0}                                                & ---                                                                                                                                                                                                                                                                           \\
            \bottomrule
        \end{tabular}
    }
\end{table}

The spectral results in Table~\ref{tab:cavity_eigenvalues} demonstrate close agreement with analytical theory:
\begin{itemize}
    \item The fundamental degenerate pair ($TE_{10}$ and $TE_{01}$, with exact $k_c = 1.0\text{ m}^{-1}$) achieves cutoff wavenumber errors of $4.95\%$ and $4.02\%$, with low divergence ratio ($R_{\mathrm{div}} \approx 0.5\%\text{--}0.6\%$). Note that on the square cavity $[0, \pi]^2$, analytical doublets ($TE_{nm}$ and $TE_{mn}$ for $n \ne m$) span a 2D degenerate eigenspace ($\lambda_{nm} = \lambda_{mn}$); the slight numerical splitting between paired modes in Table~\ref{tab:cavity_eigenvalues} arises because the randomized interior director orientations ($\vec{t}_I$, $\text{seed}=42$) intentionally break the $90^\circ$ rotational symmetry ($C_{4v}$) of the $21\times 21$ lattice, a splitting that vanishes asymptotically as $h_{\mathrm{nom}} \to 0$.
    \item Highly symmetric modes such as $TE_{02}$ ($k_c = 1.9981\text{ m}^{-1}$) and $TE_{03}$ ($k_c = 3.0037\text{ m}^{-1}$) achieve sub-$0.15\%$ precision ($0.10\%$ and $0.12\%$).
    \item Mode $TE_{22}$, which exhibited an error of $21.40\%$ and $R_{\mathrm{div}} = 30.18\%$ under $s_{\mathrm{div}} = 6.0$, is now resolved with an error of only $\mathbf{1.79\%}$ and $R_{\mathrm{div}} = 1.82\%$.
    \item The mean relative error across all 10 physical modes is $\mathbf{2.46\%}$ (a $>3.5\times$ improvement over $s_{\mathrm{div}} = 6.0$), without spurious modes in the operational band.
\end{itemize}

\subsubsection{Spectral Stagnation under Fixed Divergence Penalty}
\label{subsubsec:stagnation_fixed_sdiv}

Having achieved high accuracy on the $21 \times 21$ mesh with $s_{\mathrm{div}} = 15.0$, we examine the error behavior under systematic grid refinement across the six benchmark resolutions from \cite{stable_vnmm2d_2026}: $15 \times 15, 22 \times 22, 29 \times 29, 36 \times 36, 43 \times 43,$ and $49 \times 49$.

When holding the penalty fixed at $s_{\mathrm{div}} = 15.0$ on regular nodal distributions (isotropic random directors, $\text{seed}=42$), the mean relative cutoff wavenumber error behaves as follows:
\begin{equation}
    \begin{aligned}
        \overline{\varepsilon}_{k_c}(s_{\mathrm{div}}=15.0): \quad & 6.46\% \; (15^2) \longrightarrow 3.84\% \; (22^2) \longrightarrow 5.27\% \; (29^2)                  \\
                                                                   & \longrightarrow 5.16\% \; (36^2) \longrightarrow 6.44\% \; (43^2) \longrightarrow 7.32\% \; (49^2).
    \end{aligned}
\end{equation}
Notably, after an initial decrease from $15 \times 15$ to $22 \times 22$, the error does not converge toward zero under further refinement. Instead, it exhibits an artificial stagnation, plateauing between $5.1\%$ and $7.3\%$ for dense meshes ($N_x \ge 29$). A similar stagnation occurred with $s_{\mathrm{div}} = 6.0$ (which plateaued at $\approx 11\%\text{--}13\%$).

\subsubsection{Constraint-to-DoF Ratios and Theoretical Need for $h_{\mathrm{nom}}$-Scaling}
\label{subsubsec:constraint_ratio_scaling}

To understand the origin of this spectral stagnation, we analyze the discrete balance between the curl and divergence operators under mesh refinement:
\begin{enumerate}
    \item \textbf{Expansion of the Discrete Irrotational Nullspace:} In the continuous domain, physical solenoidal cavity fields satisfy $\nabla \cdot \vec{E} = 0$, rendering the exact cutoff spectrum independent of $s_{\mathrm{div}}$. In the discrete setting with linear polynomial trial functions ($\mathcal{P}^1$), however, the discrete divergence residue $\nabla \cdot \vec{E}_h$ is piecewise constant and non-zero. As the nominal mesh spacing $h_{\mathrm{nom}} \to 0$, the dimension of the discrete irrotational vector space expands as $\mathcal{O}(h_{\mathrm{nom}}^{-2})$.
    \item \textbf{Constraint-to-DoF Ratio in Gauss vs. MLPG:} A fundamental structural distinction separates cell-based Gauss integration from cell-free MLPG. In the cell-based VNMM-Gauss \cite{stable_vnmm2d_2026}, resolution-matched $3 \times 3$ Gauss points per quadrilateral cell evaluate the weak form at $9(N_x-1)^2$ discrete points. This provides an intrinsic over-constraint ratio of $\approx 9:1$ relative to the active interior degrees of freedom, which inherently stiffens $\mathbf{K}_{\mathrm{curl}}$ and maintains monotonic convergence even with a fixed penalty $s_{\mathrm{div}} = 6.0$. In contrast, cell-free VNMM-MLPG employs exactly one integration subdomain $\Omega_s^I$ per node. Because $\nabla \times \vec{N}$ is spatially constant on each disk under $\mathcal{P}^1$, each local stiffness matrix $\mathbf{K}_{\mathrm{curl}}^I$ contributes rank 1. The effective constraint-to-DoF ratio is therefore only $\approx 1.1:1$.
    \item \textbf{Compliance and Discrete Relaxation on Dense Meshes:} Holding $s_{\mathrm{div}}$ fixed (whether at $6.0$ or $15.0$) while refining the cloud causes the divergence constraint to become comparatively compliant relative to the expanding nullspace. The discrete eigensolver relaxes into quasi-solenoidal directions, producing the observed error plateau. Evaluating the Rayleigh quotient on the exact continuous eigenmode $\vec{E}^*$ yields $\mathbf{e}^{*T} \mathbf{K}_{\mathrm{curl}} \mathbf{e}^* / (\mathbf{e}^{*T} \mathbf{M} \mathbf{e}^*) = 0.9987$ on the $43 \times 43$ grid (an error of only $0.13\%$), demonstrating that the MLPG operators themselves represent physical fields with high fidelity.
    \item \textbf{Deduction and Justification of Superlinear Penalty Scaling ($h_{\mathrm{nom}}^{-3/2}$):} To maintain uniform spectral rigidity between discrete curl stiffness and divergence constraint across mesh refinement, the penalty parameter must scale with mesh resolution. While standard Nitsche-type trace formulations in continuous PDEs suggest a linear scaling $s_{\mathrm{div}} \propto h_{\mathrm{nom}}^{-1}$ \cite{nitsche1971uber, juntunen2009nitsche}, the discrete setting in 2D cell-free MLPG presents two competing asymptotic mechanisms:
          \begin{itemize}
              \item \emph{Continuous Scaling ($\mathcal{O}(h_{\mathrm{nom}}^{-1})$):} The divergence penalty term scales dimensionally with inverse length $h_{\mathrm{nom}}^{-1}$ relative to the kinetic/mass bilinear form.
              \item \emph{Discrete Nullspace Expansion ($\mathcal{O}(h_{\mathrm{nom}}^{-2})$):} The dimension of spurious gradient-like discrete modes grows as $\mathcal{O}(h_{\mathrm{nom}}^{-2})$, proportional to the total node count. A linear penalty $s_{\mathrm{div}} \propto h_{\mathrm{nom}}^{-1}$ grows too slowly to restrain this expanding subspace on fine clouds ($N_x \ge 36$), leading to error stagnation near $\sim 3\%$. Conversely, quadratic scaling $s_{\mathrm{div}} \propto h_{\mathrm{nom}}^{-2}$ over-stiffens coarse grids, inducing discrete locking.
          \end{itemize}
          A balanced compromise is given by the geometric mean between the continuous energy scaling and the 2D discrete nullspace expansion:
          \begin{equation}
              \sqrt{\mathcal{O}(h_{\mathrm{nom}}^{-1}) \times \mathcal{O}(h_{\mathrm{nom}}^{-2})} = \mathcal{O}(h_{\mathrm{nom}}^{-1.5}) = \mathcal{O}(h_{\mathrm{nom}}^{-3/2}).
          \end{equation}
          Consequently, we adopt the superlinear power-law scaling:
          \begin{equation}
              s_{\mathrm{div}}(h_{\mathrm{nom}}) = s_{\mathrm{div}}^{(0)} \left( \frac{h_0}{h_{\mathrm{nom}}} \right)^{1.5} \propto h_{\mathrm{nom}}^{-3/2},
              \label{eq:s_div_scaling}
          \end{equation}
          anchored to the baseline $21\times 21$ mesh ($s_{\mathrm{div}}^{(0)} = 15.0$ at $h_0 = \pi/20\text{ m}$). On general unstructured or locally adaptive point clouds without a structured grid index $N_x$, $h_{\mathrm{nom}}$ in \eqref{eq:s_div_scaling} is naturally evaluated via the cloud-averaged spacing $\overline{h} = \frac{1}{N}\sum_{I=1}^N h_I$ (or locally on each disk via $h_I$). This superlinear law grants sufficient flexibility on coarse discretizations to avoid locking, while providing the progressive stiffness needed to completely eliminate discrete nullspace leakage as $h_{\mathrm{nom}} \to 0$.
\end{enumerate}

\subsubsection{Spectral Error Reduction and Stagnation across Regular and Disordered Clouds}
\label{subsubsec:cavity_jitter_error}

To evaluate the baseline cell-free VNMM-MLPG formulation, the cavity eigenvalue problem was first solved across the six benchmark grid resolutions using the superlinear scaling law \eqref{eq:s_div_scaling}, $s_{\mathrm{div}}(h_{\mathrm{nom}}) = s_0 (h_0 / h_{\mathrm{nom}})^{1.5}$ ($s_0 = 15.0, h_0 = \pi/20\text{ m}$, yielding $s_{\mathrm{div}} \in [8.78, 55.77]$) on regular nodal distributions with isotropic random directors ($\alpha_s = 0.80$, $\text{seed}=42$). Table~\ref{tab:cavity_scaling_errors} contrasts the resulting cutoff wavenumber errors against the fixed penalty $s_{\mathrm{div}} = 15.0$.

\begin{table}[!t]
    \centering
    \small
    \caption{Cavity eigenvalue cutoff wavenumber error comparison on the 2D PEC resonant cavity across the six grid resolutions from \cite{stable_vnmm2d_2026}: fixed penalty ($s_{\mathrm{div}} = 15.0$) vs. superlinear $h_{\mathrm{nom}}$-scaled penalty ($s_{\mathrm{div}}(h_{\mathrm{nom}}) = s_0(h_0/h_{\mathrm{nom}})^{1.5}$, $s_{\mathrm{div}} \in [8.78, 55.77]$) with isotropic random directors ($\alpha_s = 0.80$, $\text{seed}=42$).}
    \label{tab:cavity_scaling_errors}
    \resizebox{\textwidth}{!}{
        \begin{tabular}{ccccccc}
            \toprule
            \textbf{Grid}             & \textbf{Nominal Spacing}                     & \textbf{Active} & \multicolumn{2}{c}{\textbf{Fixed Penalty ($s_{\mathrm{div}} = 15.0$)}} & \multicolumn{2}{c}{\textbf{Superlinear Scaled Penalty ($s_{\mathrm{div}} \propto h_{\mathrm{nom}}^{-1.5}$)}}                                                                                             \\
            \cmidrule(lr){4-5} \cmidrule(lr){6-7}
            $\boldsymbol{N_x \times N_y}$ & $\boldsymbol{h_{\mathrm{nom}}}$ \textbf{[m]} & \textbf{DoFs}   & $\boldsymbol{s_{\mathrm{div}}}$                                        & \textbf{Mean Error} $\boldsymbol{k_c}$ \textbf{(\%)}                                          & $\boldsymbol{s_{\mathrm{div}}(h_{\mathrm{nom}})}$ & \textbf{Mean Error} $\boldsymbol{k_c}$ \textbf{(\%)} \\
            \midrule
            $15 \times 15$            & 0.2244                                       & 169             & 15.0                                                                   & $6.46\%$                                                                                      & 8.78                               & $\mathbf{3.87\%}$                                    \\
            $22 \times 22$            & 0.1496                                       & 400             & 15.0                                                                   & $3.84\%$                                                                                      & 16.14                              & $\mathbf{3.55\%}$                                    \\
            $29 \times 29$            & 0.1122                                       & 729             & 15.0                                                                   & $5.27\%$                                                                                      & 24.85                              & $\mathbf{2.58\%}$                                    \\
            $36 \times 36$            & 0.0898                                       & 1156            & 15.0                                                                   & $5.16\%$                                                                                      & 34.73                              & $\mathbf{2.28\%}$                                    \\
            $43 \times 43$            & 0.0748                                       & 1681            & 15.0                                                                   & $6.44\%$                                                                                      & 45.65                              & $\mathbf{2.03\%}$                                    \\
            $49 \times 49$            & 0.0654                                       & 2209            & 15.0                                                                   & $7.32\%$                                                                                      & 55.77                              & $\mathbf{2.32\%}$                                    \\
            \bottomrule
        \end{tabular}
    }
\end{table}

As demonstrated in Table~\ref{tab:cavity_scaling_errors}, adopting the superlinear $h_{\mathrm{nom}}$-scaled penalty substantially reduces the spectral error compared to a fixed penalty:
\begin{itemize}
    \item On coarse grids ($15 \times 15$), the moderate penalty $s_{\mathrm{div}} = 8.78$ avoids numerical over-stiffening, reducing the error from $6.46\%$ to $3.87\%$.
    \item Across intermediate and fine grids ($22 \times 22$ to $49 \times 49$), the superlinear progression $s_{\mathrm{div}} \in [16.14, 55.77]$ counteracts discrete irrotational nullspace expansion, lowering mean cutoff errors from $5.1\%\text{--}7.3\%$ with fixed penalty down to $2.03\%\text{--}3.55\%$.
    \item However, error stagnation continues to exist, albeit at this lower level: the error plateaus around $\approx 2.0\%\text{--}2.3\%$ on dense grids ($2.03\%$ on $43 \times 43$ and $2.32\%$ on $49 \times 49$), yielding a sublinear empirical rate of only $p \approx 0.53$.
\end{itemize}

To evaluate the robustness of this unaugmented formulation under geometric irregularity, stochastic perturbations were next applied to all interior nodal coordinates:
\begin{equation}
    \mathbf{x}_i \longleftarrow \mathbf{x}_i + \delta\mathbf{x}_i, \quad \delta x_i, \delta y_i \sim \mathcal{U}(-j_{\mathrm{frac}} h_{\mathrm{nom}}, +j_{\mathrm{frac}} h_{\mathrm{nom}}),
\end{equation}
where $j_{\mathrm{frac}} \in \{0.0, 0.25\}$ denotes the jitter fraction relative to the nominal spacing $h_{\mathrm{nom}}$. Boundary nodes $\mathbf{x} \in \Gamma = \partial\Omega$ remain strictly unperturbed, guaranteeing exact geometric enforcement of the PEC boundary conditions ($E_t = 0$).

Fig.~\ref{fig:mlpg_cavity_jitter} depicts the resulting eigenvalue cutoff wavenumber error behavior across the six benchmark grid resolutions ($15 \times 15$ to $49 \times 49$) for $j_{\mathrm{frac}} = 0.0$ (regular) and $0.25$ ($25\%$), under the superlinear scaled penalty $s_{\mathrm{div}}(h_{\mathrm{nom}}) \propto h_{\mathrm{nom}}^{-1.5}$ and universal standard $\alpha_s = 0.80$ ($\text{seed}=42$).

\begin{figure}[!t]
    \centering
    \includegraphics[width=0.68\textwidth]{fig_mlpg_cavity_jitter.png}
    \caption{VNMM-MLPG cavity eigenvalue cutoff wavenumber error across the six benchmark grid resolutions ($15 \times 15$ to $49 \times 49$) under interior nodal position perturbations: log-log scaling of the mean cutoff error $\overline{\varepsilon}_{k_c}$ vs. nominal nodal spacing $h_{\mathrm{nom}}$ for $0\%$ (regular) and $25\%$ nodal jitter, showing empirical scaling exponents $p = 0.53$ and $p = 0.84$.}
    \label{fig:mlpg_cavity_jitter}
\end{figure}

The error trends in Fig.~\ref{fig:mlpg_cavity_jitter} reveal distinct physical regimes:
\begin{itemize}
    \item \textbf{Regular Distribution ($0\%$ Jitter, $p = 0.53$):} With regular point spacing, the superlinear penalty suppresses discrete nullspace leakage, confining the mean error to the range $2.03\%\text{--}3.87\%$ ($3.87\%$ on $15\times 15$ and $2.32\%$ on $49\times 49$), but exhibiting sublinear scaling due to local mode softening.
    \item \textbf{Moderate Perturbation ($25\%$ Jitter, $p = 0.84$):} When nodes undergo $25\%$ random dislocation, circular subdomains $\Omega_s^I$ with $\alpha_s = 0.80$ retain sufficient pairwise overlap across adjacent stencils. The eigenspectrum error systematically decreases with grid refinement: starting at $12.67\%$ on the coarse $15 \times 15$ grid, the error decreases to $5.70\%$ ($22\times 22$), $5.61\%$ ($29\times 29$), $5.15\%$ ($36\times 36$), $4.27\%$ ($43\times 43$), and reaches $\mathbf{3.95\%}$ on $49 \times 49$. Importantly, all 10 physical modes are recovered with zero spurious modes in the operational band.
\end{itemize}

Together, these regular and disordered benchmarks demonstrate that the unaugmented MLPG formulation suffers from a structural limitation: an unconstrained local nullspace that allows discrete algebraic modes to soften, producing sublinear error scaling ($p \approx 0.53$ on regular grids and $p \approx 0.84$ under $25\%$ jitter). The theoretical origin of this discrete operator rank deficiency and its mathematical resolution via Orthogonal Energy Stabilization were derived in Section~\ref{subsec:rank_deficiency_stabilization}; its numerical activation and spectral recovery are demonstrated next.

\subsubsection{Restoration of Monotonic Quadratic Convergence and Spectral Accuracy via Orthogonal Energy Stabilization}
\label{subsubsec:resolution_stagnation_stabilization}

Although the superlinear penalty law $s_{\mathrm{div}}(h_{\mathrm{nom}}) \propto h_{\mathrm{nom}}^{-1.5}$ successfully arrests the divergence leakage that plagued fixed-penalty formulations, a deeper inspection of Table~\ref{tab:cavity_scaling_errors} and Fig.~\ref{fig:mlpg_cavity_jitter} reveals that asymptotic error scaling remains sublinear ($p = 0.53$) on regular grids, while the fundamental mode ($TE_{10}$, exact $k_c = 1.0000\text{ m}^{-1}$) stagnates at $k_c \approx 0.9551\text{ m}^{-1}$ (an error of $4.49\%$ persisting across all fine meshes).

As established analytically in Section~\ref{subsec:rank_deficiency_stabilization}, this slow convergence does not stem from an inability to interpolate physical fields, but rather from the \textbf{local rank-2 deficiency of the unaugmented linear stiffness operator $\mathbf{K}_I^{\mathrm{local}}$}. Evaluating constant gradients on a single disk leaves local shear and deviatoric modes with zero stiffness, permitting discrete algebraic modes to soften under unconstrained Rayleigh quotient minimization.

To restore full rank and eliminate this stagnation, we activate the Orthogonal Energy Stabilization operator $\mathbf{K}_{I, \mathrm{stab}}^{\mathrm{local}}$ derived in \eqref{eq:k_stab_analytic}, parameterized by the dimensionless VMS subgrid scaling law $\alpha_{\mathrm{stab}}(h_{\mathrm{nom}}) = c_{\mathrm{stab}} (h_{\mathrm{nom}}/L)^2$ from \eqref{eq:alpha_stab_vms_scaling} (with reference length $L = \pi\text{ m}$). Table~\ref{tab:cavity_convergence_stabilization} presents the resulting convergence across the benchmark grid resolutions for representative calibration parameters $c_{\mathrm{stab}} \in \{0, 20, 40, 70, 100\}$.

\begin{table}[!t]
    \centering
    \small
    \caption{Restoration of monotonic quadratic cavity eigenvalue convergence and high spectral accuracy via Dimensionless Orthogonal Energy Stabilization ($\alpha_{\mathrm{stab}} = c_{\mathrm{stab}} (h_{\mathrm{nom}}/L)^2$, $L = \pi\text{ m}$) across the benchmark grid resolutions ($0\%$ jitter, isotropic random directors, superlinear $s_{\mathrm{div}} \propto h_{\mathrm{nom}}^{-1.5}$). The unstabilized formulation ($c_{\mathrm{stab}} = 0$, $p = 0.53$) is compared against stabilized formulations ($c_{\mathrm{stab}} = 20, 40, 70, 100$), showing the recovery of asymptotic quadratic convergence ($p = 2.39$).}
    \label{tab:cavity_convergence_stabilization}
    \resizebox{\textwidth}{!}{
        \begin{tabular}{ccccccc}
            \toprule
            \textbf{Grid}                                                                                     & \textbf{Nominal Spacing}                     & \textbf{Unstabilized}                  & \multicolumn{4}{c}{\textbf{Dimensionless Orthogonally Stabilized} ($\boldsymbol{\alpha_{\mathrm{stab}} = c_{\mathrm{stab}} (h_{\mathrm{nom}}/L)^2}$)}                                                                                    \\
            \cmidrule(lr){4-7}
            $\boldsymbol{N_x \times N_y}$                                                                     & $\boldsymbol{h_{\mathrm{nom}}}$ \textbf{[m]} & $\boldsymbol{c_{\mathrm{stab}} = 0}$   & $\boldsymbol{c_{\mathrm{stab}} = 20}$ & $\boldsymbol{c_{\mathrm{stab}} = 40}$ & $\boldsymbol{c_{\mathrm{stab}} = 70}$ & $\boldsymbol{c_{\mathrm{stab}} = 100}$ \\
            \midrule
            $15 \times 15$                                                                                    & 0.2244                                       & $3.87\%$                               & $6.69\%$                              & $11.14\%$                             & $17.67\%$                             & $23.81\%$                              \\
            $22 \times 22$                                                                                    & 0.1496                                       & $3.55\%$                               & $3.31\%$                              & $4.41\%$                              & $7.11\%$                              & $9.73\%$                               \\
            $29 \times 29$                                                                                    & 0.1122                                       & $2.58\%$                               & $1.84\%$                              & $2.06\%$                              & $3.34\%$                              & $5.38\%$                               \\
            $36 \times 36$                                                                                    & 0.0898                                       & $2.28\%$                               & $2.11\%$                              & $2.24\%$                              & $2.87\%$                              & $3.94\%$                               \\
            $43 \times 43$                                                                                    & 0.0748                                       & $2.03\%$                               & $1.59\%$                              & $\mathbf{1.34\%}$                     & $1.39\%$                              & $1.67\%$                               \\
            $49 \times 49$                                                                                    & 0.0654                                       & $2.32\%$                               & $1.92\%$                              & $1.59\%$                              & $1.29\%$                              & $\mathbf{1.17\%}$                      \\
            \midrule
            \multicolumn{2}{c}{\textbf{Convergence Rate} $\boldsymbol{p}$ ($\boldsymbol{\mathcal{O}(h_{\mathrm{nom}}^p)}$):} & $\mathbf{0.53}$                              & $\mathbf{1.06}$                       & $\mathbf{1.64}$                       & $\mathbf{2.15}$                       & $\mathbf{2.39}$                        \\
            \bottomrule
        \end{tabular}
    }
\end{table}

\begin{figure}[!t]
    \centering
    \includegraphics[width=\textwidth]{fig_mlpg_cavity_stabilization.png}
    \caption{Cavity eigenvalue cutoff wavenumber convergence under Dimensionless Orthogonal Energy Stabilization ($\alpha_{\mathrm{stab}} = c_{\mathrm{stab}} (h_{\mathrm{nom}}/L)^2$, $L = \pi\text{ m}$) across parameters $c_{\mathrm{stab}} \in \{0, 20, 40, 70, 100\}$: (a) Regular grids ($0\%$ jitter); (b) Perturbed point clouds ($25\%$ jitter).}
    \label{fig:mlpg_cavity_stabilization}
\end{figure}

The numerical results in Table~\ref{tab:cavity_convergence_stabilization} and Fig.~\ref{fig:mlpg_cavity_stabilization} demonstrate the significant impact of the orthogonal stabilization operator:
\begin{enumerate}
    \item \textbf{Restoration of Monotonic Quadratic Convergence:} As illustrated in Fig.~\ref{fig:mlpg_cavity_stabilization}(a), increasing the dimensionless stabilization parameter systematically accelerates the empirical convergence rate $p$ from sublinear ($p = 0.53$ at $c_{\mathrm{stab}} = 0$) through linear ($p = 1.06$ at $c_{\mathrm{stab}} = 20$) and superlinear ($p = 1.64$ at $c_{\mathrm{stab}} = 40$, $p = 2.15$ at $c_{\mathrm{stab}} = 70$) up to quadratic asymptotic convergence ($p = 2.39$ at $c_{\mathrm{stab}} = 100$). This restores the asymptotic rate achieved by the cell-based VNMM-Gauss ($p = 2.25$ in Table~\ref{tab:jitter_benchmark_comparison}).
    \item \textbf{Attainment of High Spectral Accuracy on Fine Meshes:} On the dense $49 \times 49$ grid, where the unstabilized formulation was stalled at $2.32\%$, orthogonal stabilization drives the mean error down to $1.29\%$ ($c_{\mathrm{stab}} = 70$) and $\mathbf{1.17\%}$ ($c_{\mathrm{stab}} = 100$). This surpasses the finest-grid accuracy of VNMM-Gauss ($1.53\%$).
    \item \textbf{Coarse-Grid Stiffness vs. Fine-Grid Asymptotics:} The parametric variation reveals the trade-off inherent in quadratic subgrid scaling: on very coarse grids ($15 \times 15$), where $h_{\mathrm{nom}}/L$ is relatively large ($h_{\mathrm{nom}}/L \approx 0.071$), a high value such as $c_{\mathrm{stab}} = 100$ injects non-negligible artificial stiffness ($\alpha_{\mathrm{stab}} \approx 0.51$), raising the coarse-grid error to $23.81\%$. Conversely, moderate constants such as $c_{\mathrm{stab}} \in [20, 40]$ keep the coarse-grid error well-contained ($6.69\%$ and $11.14\%$) while still significantly lowering errors on intermediate grids. For $c_{\mathrm{stab}} = 100$, as soon as the grid reaches $29 \times 29$, the quadratic $\mathcal{O}((h_{\mathrm{nom}}/L)^2)$ decay rapidly reduces the stabilization magnitude ($\alpha_{\mathrm{stab}} \le 0.13$), allowing physical modes to dominate and yielding the highest asymptotic accuracy.
    \item \textbf{Robust Quadratic Convergence under Perturbed Clouds ($25\%$ Jitter):} As shown in Fig.~\ref{fig:mlpg_cavity_stabilization}(b), orthogonal energy stabilization is equally effective on disordered point clouds. Under $25\%$ jitter, the unaugmented formulation ($c_{\mathrm{stab}} = 0$) exhibits a large coarse-grid error ($12.67\%$ on $15\times 15$) and stalls at $3.95\%$ on $49 \times 49$ ($p = 0.84$). Moderate stabilization ($c_{\mathrm{stab}} = 20$ and $40$) reduces the error across every single mesh compared to $c_{\mathrm{stab}} = 0$: in particular, it lowers the coarse $15\times 15$ error by nearly $3\times$ (from $12.67\%$ down to $4.44\%$ and $5.02\%$, respectively, reaching $3.37\%$ and $2.80\%$ on $49\times 49$), which explains why the overall log-log regression slope across the 6 meshes appears flatter ($p = 0.18$ and $0.54$). When stronger stabilization is activated ($c_{\mathrm{stab}} = 100$), the asymptotic $(h_{\mathrm{nom}}/L)^2$ decay dominates intermediate and fine meshes, recovering full quadratic convergence ($p = 2.30$) and driving the mean cutoff error down to $\mathbf{1.45\%}$ on $49 \times 49$. For perturbed meshes, $c_{\mathrm{stab}} = 70$ provides the most balanced overall performance across all resolutions, achieving $p = 1.62$ and an error of $1.97\%$ on $49 \times 49$ while keeping the coarse $15 \times 15$ error down to $11.88\%$ (vs.\ $18.61\%$ for $c_{\mathrm{stab}} = 100$). This confirms that restoring full local stiffness rank eliminates discrete mode softening in both regular and disordered meshless configurations.
\end{enumerate}

Consequently, based on these comparative findings across regular and perturbed point distributions, $c_{\mathrm{stab}} = 70$ is adopted as the standard operational stabilization parameter for the computational benchmarks in Stage~6: while $c_{\mathrm{stab}} = 100$ achieves the steepest asymptotic slope ($p = 2.30\text{--}2.39$) and lowest error on the finest $49\times 49$ grid ($1.17\%$ regular and $1.45\%$ perturbed), $c_{\mathrm{stab}} = 70$ retains near-quadratic convergence ($p = 2.15$ regular, $1.29\%$ error on $49\times 49$) while substantially reducing artificial stiffness on coarse $15\times 15$ and $22\times 22$ point clouds.

\subsection{Stage 6: Computational Benchmark: VNMM-MLPG vs. VNMM-Gauss}
\label{subsec:stage6_benchmark}

A direct performance comparison was conducted between the proposed cell-free VNMM-MLPG (with superlinear scaled penalty $s_{\mathrm{div}}(h_{\mathrm{nom}}) \propto h_{\mathrm{nom}}^{-1.5}$, $s_{\mathrm{div}} \in [8.78, 55.77]$, and dimensionless orthogonal energy stabilization $c_{\mathrm{stab}} = 70$) and cell-based VNMM-Gauss ($3 \times 3$ Gauss points per cell, $s_{\mathrm{div}} = 6.0$ \cite{stable_vnmm2d_2026}) across the standard six grid resolutions from \cite{stable_vnmm2d_2026}: $15 \times 15, 22 \times 22, 29 \times 29, 36 \times 36, 43 \times 43,$ and $49 \times 49$. Note that each method is evaluated under its respective optimal variational calibration: in VNMM-Gauss, the $9:1$ quadrature over-constraint ratio of $3\times 3$ cell Gauss points already confers full rank to cell stiffness matrices but causes severe volumetric locking if $s_{\mathrm{div}}$ is scaled to large values, making the fixed penalty $s_{\mathrm{div}} = 6.0$ optimal \cite{stable_vnmm2d_2026}; whereas in single-disk VNMM-MLPG ($1.1:1$ constraint ratio), the superlinear law $s_{\mathrm{div}}(h_{\mathrm{nom}}) \propto h_{\mathrm{nom}}^{-1.5}$ coupled with $c_{\mathrm{stab}} = 70$ represents the corresponding canonical parameterization (moreover, even with $c_{\mathrm{stab}} = 0$, Table~\ref{tab:cavity_scaling_errors} shows that VNMM-MLPG already achieves $3.87\%\text{--}6.46\%$ error on $15\times 15$ compared to $25.25\%$ for VNMM-Gauss). Both solvers were executed on identical hardware (AMD Ryzen / Linux) under single-thread execution with isotropic random vector directors ($\text{seed}=42$). To evaluate both computational efficiency and resilience to geometric disorder, evaluations were conducted on regular lattices ($0\%$ jitter) as well as perturbed point clouds with $25\%$ stochastic nodal coordinate jitter ($j_{\mathrm{frac}} = 0.25$).

Table~\ref{tab:benchmark_comparison} compares the global matrix assembly times and stiffness matrix sparsity (non-zero entries, NNZ). Table~\ref{tab:jitter_benchmark_comparison} and Fig.~\ref{fig:benchmark} present the resulting cutoff wavenumber errors and empirical asymptotic convergence rates $p$ ($\overline{\varepsilon}_{k_c} \propto h_{\mathrm{nom}}^p$) under both $0\%$ and $25\%$ nodal jitter.

\begin{table}[!t]
    \centering
    \small
    \caption{Computational benchmark on the 2D PEC resonant cavity across the six grid resolutions matching \cite{stable_vnmm2d_2026}: VNMM-MLPG (cell-free with superlinear scaled $s_{\mathrm{div}} \in [8.78, 55.77]$ and stabilization $c_{\mathrm{stab}} = 70$) vs. VNMM-Gauss (cell-based with $3 \times 3$ Gauss quadrature and $s_{\mathrm{div}} = 6.0$). Assembly runtime and global stiffness matrix non-zero entries (NNZ) are compared.}
    \label{tab:benchmark_comparison}
    \resizebox{0.88\textwidth}{!}{
        \begin{tabular}{ccccccc}
            \toprule
            \textbf{Grid}                 & \textbf{Active} & \multicolumn{2}{c}{\textbf{Assembly Time [s]}} & \textbf{Assembly} & \multicolumn{2}{c}{\textbf{Stiffness NNZ}}                                                      \\
            \cmidrule(lr){3-4} \cmidrule(lr){6-7}
            $\boldsymbol{N_x \times N_y}$ & \textbf{DoFs}   & \textbf{Gauss}                                 & \textbf{MLPG}     & \textbf{Speedup}                           & \textbf{Gauss} & \textbf{MLPG}            \\
            \midrule
            $15 \times 15$            & 169             & 0.1239 s                                       & \textbf{0.0624 s} & \textbf{1.99}$\times$                      & 3,091          & \textbf{2,467} ($-20\%$) \\
            $22 \times 22$            & 400             & 0.2670 s                                       & \textbf{0.1060 s} & \textbf{2.52}$\times$                      & 7,802          & \textbf{6,140} ($-21\%$) \\
            $29 \times 29$            & 729             & 0.4529 s                                       & \textbf{0.1593 s} & \textbf{2.84}$\times$                      & 14,677         & \textbf{11,509} ($-22\%$) \\
            $36 \times 36$            & 1156            & 0.6813 s                                       & \textbf{0.2263 s} & \textbf{3.01}$\times$                      & 23,634         & \textbf{18,460} ($-22\%$) \\
            $43 \times 43$            & 1681            & 0.9754 s                                       & \textbf{0.2855 s} & \textbf{3.42}$\times$                      & 35,409         & \textbf{26,813} ($-24\%$) \\
            $49 \times 49$            & 2209            & 1.2750 s                                       & \textbf{0.3459 s} & \textbf{3.69}$\times$                      & 44,939         & \textbf{35,385} ($-21\%$) \\
            \bottomrule
        \end{tabular}
    }
\end{table}

\begin{figure}[!t]
    \centering
    \includegraphics[width=0.68\textwidth]{fig_mlpg_benchmark_comparison.png}
    \caption{Computational benchmark comparing stabilized VNMM-MLPG ($c_{\mathrm{stab}} = 70$) against VNMM-Gauss across the six grid resolutions matching \cite{stable_vnmm2d_2026}: Cavity eigenvalue cutoff wavenumber convergence between VNMM-Gauss (dashed lines) and VNMM-MLPG (solid lines) under $0\%$ (blue) and $25\%$ (green) nodal position jitter.}
    \label{fig:benchmark}
\end{figure}

\begin{table}[!t]
    \centering
    \small
    \caption{Cavity eigenvalue cutoff wavenumber convergence under nodal perturbations ($0\%$ and $25\%$ jitter) across the six benchmark grid resolutions: VNMM-Gauss (cell-based with $3 \times 3$ Gauss quadrature, $s_{\mathrm{div}} = 6.0$) vs. VNMM-MLPG (cell-free, superlinear $s_{\mathrm{div}} \propto h_{\mathrm{nom}}^{-1.5}$, stabilized with $c_{\mathrm{stab}} = 70$). Empirical asymptotic convergence rates $p$ are computed via linear regression of $\log(\overline{\varepsilon}_{k_c})$ vs. $\log(h_{\mathrm{nom}})$.}
    \label{tab:jitter_benchmark_comparison}
    \resizebox{0.85\textwidth}{!}{
        \begin{tabular}{cccccc}
            \toprule
            \textbf{Grid}                                                                                     & \textbf{Nominal}                          & \multicolumn{2}{c}{\textbf{0\% Jitter (Regular)}} & \multicolumn{2}{c}{\textbf{25\% Jitter}}                                           \\
            \cmidrule(lr){3-4} \cmidrule(lr){5-6}
            $\boldsymbol{N_x \times N_y}$                                                                     & $\boldsymbol{h_{\mathrm{nom}}}$ \textbf{[m]} & \textbf{Gauss}                                    & \textbf{MLPG}                            & \textbf{Gauss}    & \textbf{MLPG}      \\
            \midrule
            $15 \times 15$                                                                                    & 0.2244                                    & $25.25\%$                                         & $\mathbf{17.67\%}$                       & $24.67\%$         & $\mathbf{11.88\%}$ \\
            $22 \times 22$                                                                                    & 0.1496                                    & $13.46\%$                                         & $\mathbf{7.11\%}$                        & $11.49\%$         & $\mathbf{6.82\%}$  \\
            $29 \times 29$                                                                                    & 0.1122                                    & $6.18\%$                                          & $\mathbf{3.34\%}$                        & $5.94\%$          & $\mathbf{3.77\%}$  \\
            $36 \times 36$                                                                                    & 0.0898                                    & $4.22\%$                                          & $\mathbf{2.87\%}$                        & $3.10\%$          & $\mathbf{2.36\%}$  \\
            $43 \times 43$                                                                                    & 0.0748                                    & $2.46\%$                                          & $\mathbf{1.39\%}$                        & $1.96\%$          & $\mathbf{1.77\%}$  \\
            $49 \times 49$                                                                                    & 0.0654                                    & $1.53\%$                                          & $\mathbf{1.29\%}$                        & $\mathbf{1.63\%}$ & $1.97\%$           \\
            \midrule
            \multicolumn{2}{c}{\textbf{Convergence Rate} $\boldsymbol{p}$ ($\boldsymbol{\mathcal{O}(h_{\mathrm{nom}}^p)}$):} & $\mathbf{2.25}$                           & $\mathbf{2.15}$                                   & $\mathbf{2.29}$                          & $\mathbf{1.62}$                        \\
            \bottomrule
        \end{tabular}
    }
\end{table}

The comparative results in Tables~\ref{tab:benchmark_comparison} and \ref{tab:jitter_benchmark_comparison} and Fig.~\ref{fig:benchmark} demonstrate clear advantages of VNMM-MLPG over cell-based integration:
\begin{enumerate}
    \item \textbf{Assembly Speedup ($2.0\times\text{--}3.7\times$):} By evaluating closed-form analytical formulas \eqref{eq:k_curl_analytic}--\eqref{eq:m_local_analytic} over $N$ disk subdomains instead of looping over thousands of numerical Gauss points with point-in-cell tests and $k$-d tree searches, VNMM-MLPG accelerates global matrix assembly by a factor of $1.99\times$ to $3.69\times$ (scaling from $0.062\text{ s}$ on $15 \times 15$ to $0.346\text{ s}$ on $49 \times 49$, compared to $0.124\text{ s}$ and $1.275\text{ s}$ for Gauss integration).
    \item \textbf{$20\%\text{--}24\%$ Sparser Global Operators:} In cell-based integration, a single background cell interacts with all nodes belonging to the stencils of its 9 Gauss points, inflating the sparsity pattern of $\mathbf{K}$. In VNMM-MLPG, each node $I$ interacts strictly with the 6 nodes in its own stencil $\Omega_u^I$. Consequently, global matrices contain $20\%\text{--}24\%$ fewer non-zero entries (e.g., $35,385$ vs. $44,939$ NNZ on $49 \times 49$), reducing memory consumption and improving cache locality.
    \item \textbf{Global Spectral Accuracy and Convergence ($0\%$ Jitter):} Across every single regular grid resolution in Table~\ref{tab:jitter_benchmark_comparison}, stabilized VNMM-MLPG consistently outperforms cell-based VNMM-Gauss. On the coarsest mesh ($15 \times 15$, $h_{\mathrm{nom}} \approx 0.2244\text{ m}$), where cell-based Gauss integration suffers from severe quadrature softening and boundary truncation ($25.25\%$ error), stabilized VNMM-MLPG limits the error to $17.67\%$ ($1.43\times$ lower). As the grid is refined, stabilized VNMM-MLPG maintains its accuracy lead across all intermediate resolutions ($7.11\%$ vs. $13.46\%$ on $22\times 22$, $3.34\%$ vs. $6.18\%$ on $29\times 29$, $2.87\%$ vs. $4.22\%$ on $36\times 36$, and $1.39\%$ vs. $2.46\%$ on $43\times 43$), achieving $1.29\%$ on $49\times 49$ (surpassing the $1.53\%$ of VNMM-Gauss). Furthermore, orthogonal energy stabilization recovers near-quadratic asymptotic convergence ($p = 2.15$, matching the $p = 2.25$ of VNMM-Gauss as well as the theoretical $\mathcal{O}(h^2)$ eigenvalue convergence order of lowest-order conforming N\'{e}d\'{e}lec edge finite elements $\mathcal{ND}_1$ \cite{nedelec1980mixed, monk2003finite}), proving that discrete mode softening is eliminated while retaining all cell-free advantages.
    \item \textbf{Resilience to Point Cloud Disorder and Coarse-Grid Robustness ($25\%$ Jitter):} Under $25\%$ stochastic nodal jitter, stabilized VNMM-MLPG outperforms cell-based VNMM-Gauss across five of the six benchmark resolutions ($15\times 15$ through $43\times 43$), exhibiting smooth monotonic convergence ($p = 1.62$) and driving the cutoff frequency error below $2\%$ ($1.77\%$ on $43\times 43$ and $1.97\%$ on $49\times 49$). 
    
    On the coarsest perturbed point cloud ($15 \times 15$, $h_{\mathrm{nom}} \approx 0.2244\text{ m}$), stabilized VNMM-MLPG yields an error of $11.88\%$, compared to $24.67\%$ for VNMM-Gauss, representing a $2.08\times$ error reduction. This difference reflects a distinct geometric mechanism: in cell-based methods, irregular nodal displacements distort the geometric alignment between background Cartesian integration cells and shifting local support stencils. At coarse resolutions, individual background cells frequently span disjoint stencils with disparate local gradient orientations, resulting in intra-cell quadrature under-integration, artificial softening, and cell-boundary clipping. In contrast, VNMM-MLPG integrates the weak form over autonomous circular subdomains $\Omega_s^I$ centered directly at each physical node $\mathbf{x}_I$. Because each local disk is integrated in closed form using the unique, non-switching polynomial stencil $\Omega_u^I$ associated with node $I$, the discrete energy evaluation is entirely immune to cell-mesh skewness, element aspect-ratio degradation, and intra-cell stencil transitions. This confirms that VNMM-MLPG is inherently more robust against point-cloud disorder on coarse configurations than conventional cell-based quadrature schemes. While cell-based Gauss exhibits a marginally steeper asymptotic slope ($p = 2.29$) due to background grid rigidity on very fine meshes ($49\times 49$), it incurs a $2.0\times\text{--}3.7\times$ assembly time penalty and generates $20\%\text{--}24\%$ denser global operators. Overall, VNMM-MLPG delivers consistently superior accuracy across coarse-to-intermediate resolutions with substantial computational speedups and complete freedom from background integration grids.
\end{enumerate}

%% =========================================================================
%% SECTION 4: DISCUSSION
%% =========================================================================
\section{Discussion: Advantages, Deficiencies, and Open Questions}
\label{sec:discussion}

\subsection{Synthesis of Advantages}
\label{subsec:advantages_synthesis}

Table~\ref{tab:method_comparison} compares the features of VNMM-MLPG against the previous cell-based VNMM-Gauss and standard edge finite elements. The advantages achieved by VNMM-MLPG include:
\begin{itemize}
    \item \textbf{Freedom from Mesh Scaffolding:} The method is truly cell-free. Decoupling spatial integration from background grids eliminates the risk of geometric meshing failures.
    \item \textbf{Closed-Form Integration:} Numerical quadrature under-integration traps and artificial softening across stencil transitions are eliminated.
    \item \textbf{Symmetric Operators:} Unlike standard Petrov--Galerkin and LPIM formulations that produce unsymmetric matrices via row-by-row assembly, the proposed Symmetric Subdomain Galerkin scheme yields strictly symmetric positive-semidefinite operators, precluding complex spurious modes and restoring discrete self-adjointness.
    \item \textbf{Orthogonal Energy Stabilization:} Augmenting local linear stiffness with an analytical subgrid projection operator $\mathbf{K}_{\mathrm{stab}}^I$ scaled via VMS principles as $\alpha_{\mathrm{stab}} = c_{\mathrm{stab}}(h_{\mathrm{nom}}/L)^2$ overcomes the localized rank-2 deficiency of linear fields on circular disks (restoring full rank 4 over the linear deformation subspace), eliminating discrete mode softening and restoring superlinear-to-quadratic spectral convergence ($p = 2.15\text{--}2.39$ on regular grids and $p = 1.62$ under $25\%$ jitter).
    \item \textbf{Enhanced Spectral Accuracy:} Stabilized VNMM-MLPG ($c_{\mathrm{stab}} = 70$) outperforms cell-based VNMM-Gauss across all six regular benchmark resolutions (from $15\times 15$ to $49\times 49$, achieving $1.29\%$ error on $49\times 49$) and in five of the six perturbed configurations (achieving a $2.08\times$ error reduction on coarse $15\times 15$ clouds).
\end{itemize}

\begin{table}[!t]
    \centering
    \small
    \caption{Comparative synthesis of computational features among CEM methods.}
    \label{tab:method_comparison}
    \resizebox{\textwidth}{!}{
        \begin{tabular}{lccc}
            \toprule
            \textbf{Methodological Feature}      & \textbf{Edge FEM (Nédélec)} & \textbf{VNMM-Gauss (Previous)} & \textbf{VNMM-MLPG (Present)}                 \\
            \midrule
            Mesh Dependency                      & Conforming Triangulation    & Cartesian Background Cells     & \textbf{Truly Cell-Free}                     \\
            Integration Scheme                   & Numerical Simplex Gauss     & $3 \times 3$ Cell Quadrature   & \textbf{Closed-Form Analytic}                \\
            Intra-Cell Stencil Switching         & None                        & Present (requires dense Gauss) & \textbf{None (Autonomous Disks)}             \\
            Matrix Symmetry                      & Symmetric                   & Symmetric                      & \textbf{Symmetric}                           \\
            Matrix Sparsity (Stiffness NNZ)      & Baseline                    & Baseline                       & \textbf{20\%--24\% Sparser}                  \\
            Assembly Time Speedup                & Baseline                    & Baseline                       & \textbf{2.0$\times$--3.7$\times$ Faster}     \\
            Patch Test Satisfaction              & Machine Precision           & Machine Precision              & \textbf{Machine Precision ($\sim 10^{-14}\text{--}10^{-15}$)} \\
            Spurious Mode Treatment              & Discrete de Rham complex    & Div-Curl Regularization        & \textbf{Superlinear Div-Curl Scaling}        \\
            Local Operator Rank Stability        & Conforming exact sequence   & Dense $3 \times 3$ cell quadrature & \textbf{Orthogonal Energy Stabilization} \\
            \bottomrule
        \end{tabular}
    }
\end{table}

\subsection{Current Deficiencies and Open Challenges}
\label{subsec:deficiencies_challenges}

Despite its proven efficacy, several deficiencies, operational restrictions, and open questions remain for the VNMM-MLPG:
\begin{enumerate}
    \item \textbf{Higher-Order Cavity Mode Dispersion:} As demonstrated in the spectral cavity benchmarks (Section~\ref{subsec:stage5_cavity}), rapidly oscillating higher-order cavity modes exhibit higher sensitivity to the divergence penalty and larger discrete divergence ratios (e.g., $\mathrm{TE}_{22}$ reaching $21.40\%$ error with $R_{\mathrm{div}} = 30.18\%$ at $s_{\mathrm{div}} = 6.0$ before dropping to $1.79\%$ at $s_{\mathrm{div}} = 15.0$, and $\mathrm{TE}_{30}/\mathrm{TE}_{03}$ exhibiting $R_{\mathrm{div}} \approx 3.2\%\text{--}3.4\%$ vs.\ $0.5\%\text{--}0.6\%$ for $\mathrm{TE}_{10}/\mathrm{TE}_{01}$ in Table~\ref{tab:cavity_eigenvalues}). This numerical dispersion stems from the linear polynomial capacity of the $\mathcal{P}^1$ space: resolving modes with multiple spatial oscillations requires denser nodal distributions ($h_{\mathrm{local}} \ll \lambda_{\mathrm{wavelength}} / 10$).
    \item \textbf{Global Non-Conformity, Stabilization Trade-Offs, and Higher-Order PIM ($\mathcal{P}^2$):} Unlike conforming edge finite elements that satisfy the discrete de Rham sequence without penalty parameters, VNMM-MLPG operates on globally non-conforming point clouds and relies on the coupled action of superlinear divergence penalization ($s_{\mathrm{div}} \propto h_{\mathrm{nom}}^{-1.5}$) and orthogonal subgrid stabilization ($\alpha_{\mathrm{stab}} \propto h_{\mathrm{nom}}^2$). As shown in Sections~\ref{subsubsec:constraint_ratio_scaling} and~\ref{subsubsec:resolution_stagnation_stabilization}, excessive $s_{\mathrm{div}}$ on coarse grids can induce mild over-penalization of high-order physical modes, while choosing $c_{\mathrm{stab}}$ involves a trade-off between coarse-grid stiffness ($c_{\mathrm{stab}} = 70$) and fine-grid asymptotic accuracy ($1.29\%$ on $49\times 49$); furthermore, the superlinear penalty $s_{\mathrm{div}} \propto h_{\mathrm{nom}}^{-1.5}$ raises the asymptotic spectral condition number scaling of the stiffness matrix from $\mathcal{O}(h_{\mathrm{nom}}^{-2})$ to $\mathcal{O}(h_{\mathrm{nom}}^{-3.5})$ on extremely dense grids. A natural progression that simultaneously addresses both items (1) and (2) is extending the Vector PIM formulation to quadratic vector polynomial spaces ($\mathcal{P}^2$, requiring 12 nodes per stencil in 2D), where linear variations of $\nabla \times \vec{E}$ and $\nabla \cdot \vec{E}$ over each disk confer higher local operator rank naturally while substantially reducing high-frequency dispersion.
    \item \textbf{Curved Boundaries, Re-Entrant Corners, and Material Interfaces:} The closed-form analytical truncation formulas derived in Section~\ref{subsubsec:truncated_subdomains_analytic} and \ref{app:truncated_geometry} assume piecewise planar domain walls and orthogonal convex corners. Extending closed-form integration to arbitrary curved CAD boundaries (via local polygonal Green's theorem approximations), handling re-entrant metallic corners where electric fields exhibit non-polynomial singularities ($\vec{E} \sim r^{\pi/\omega - 1}$), and treating multi-dielectric interfaces with normal field jumps ($[\![ E_n ]\!] \ne 0$, $[\![ \epsilon_r \vec{E} \cdot \vec{n} ]\!] = 0$) through interface-aligned directors, split-disk integration, or interface Lagrange multipliers (as investigated for scalar wave scattering in the Method of Finite Spheres \cite{nicomedes2021acoustic}) represent key open challenges.
    \item \textbf{Extension to Three Dimensions:} The mathematical structure derived herein extends naturally to 3D by replacing circular disks with spheres of radius $r_s$. In 3D, the complete linear vector basis $\mathcal{P}_1^3$ contains 12 terms, and spherical moments evaluate analytically to $V_0 = \frac{4}{3}\pi r_s^3$ and $I_2 = \frac{4}{15}\pi r_s^5$. Implementing this closed-form 3D extension, along with analytical spherical-cap truncation at planar boundaries, constitutes the immediate next phase of this research.
\end{enumerate}

%% =========================================================================
%% SECTION 5: CONCLUSIONS AND FUTURE WORK
%% =========================================================================
\section{Conclusions and Future Work}
\label{sec:conclusions}

This work presented the formulation, implementation, and systematic validation of a truly cell-free Vector Nodal Meshless Method via Meshless Local Petrov--Galerkin (VNMM-MLPG) for two-dimensional computational electromagnetics. By combining the directional $\mathcal{P}^1$ Vector Nodal Meshless Method \cite{ortiz2021nodal, stable_vnmm2d_2026} with weak-form integration over autonomous circular subdomains $\Omega_s^I$, the method eliminates background quadrature grids, intra-cell stencil switching, and geometric meshing bottlenecks. Because fixed linear polynomial stencils yield spatially constant curl and divergence on each subdomain, Stokes' and Gauss' theorems rule out test functions vanishing on the subdomain boundary (MLPG1), motivating the adoption of the Local Galerkin formulation (MLPG6).

Key achievements of the proposed VNMM-MLPG include:
\begin{enumerate}
    \item \textbf{Closed-Form Analytical Operator Formulation:} Derived closed-form analytical expressions for local curl, divergence, and mass matrices on interior disks and boundary-truncated circular segments in terms of the directional inverse collocation matrix $\boldsymbol{\beta}^I$ and exact geometric moments, eliminating numerical quadrature error and lens-intersection cubatures.
    \item \textbf{Overlapping Symmetric Subdomain Galerkin Assembly:} Overcame the structural asymmetry and loss of discrete coercivity of classical row-by-row MLPG/LPIM by formulating an additive overlapping-subdomain Ritz--Galerkin energy assembly, guaranteeing symmetric positive-(semi)definite global operators.
    \item \textbf{Machine-Precision Linear Consistency:} Verified linear reproduction via the Galerkin patch test across seven affine vector states on regular and 25\% perturbed point clouds, achieving interior errors on the order of $\sim 10^{-14}\text{--}10^{-15}$ ($\le 1.29 \times 10^{-13}$).
    \item \textbf{MMS Convergence and Subdomain Containment Verification:} Confirmed monotonic convergence with rates of $\mathcal{O}(h_{\mathrm{nom}}^{1.93})$ (pure Dirichlet) and $\mathcal{O}(h_{\mathrm{nom}}^{1.85})$ (mixed Dirichlet--Neumann) in $L^2(\vec{E})$, and $\mathcal{O}(h_{\mathrm{nom}}^{1.05})$ and $\mathcal{O}(h_{\mathrm{nom}}^{1.02})$ in the $H(\mathbf{curl})$ semi-norm, while verifying that the classical MLPG/LPIM subdomain scaling $\alpha_s = 0.80$ \cite{atluri2002meshless, liu2009mesh} satisfies both global coverage ($\alpha_s \ge \sqrt{2}/2$) and 6-node directional stencil containment ($\alpha_s \le 1.15$).
    \item \textbf{Directional Modal Projection and VMS Orthogonal Stabilization:} Showed that single-disk integration leaves two harmonic zero-energy deformation modes ($\mathbf{v}_{\mathrm{shear}}^I, \mathbf{v}_{\mathrm{hyp}}^I$) in the local div-curl kernel, projected them analytically onto the directional degrees of freedom, and formulated a consistent VMS Orthogonal Energy Stabilization ($\alpha_{\mathrm{stab}} \propto h_{\mathrm{nom}}^2$) coupled with superlinear divergence penalty scaling ($s_{\mathrm{div}} \propto h_{\mathrm{nom}}^{-1.5}$). This restores full rank 4 over the linear deformation subspace, eliminates mode softening, recovers near-quadratic spectral convergence ($p = 2.15$ on regular grids, driving fine-grid error down to $1.29\%$ on $49\times 49$), and shifts all non-physical gradient modes outside the operational spectrum ($s_{\mathrm{div}} \ge 4$).
    \item \textbf{Benchmark Accuracy and Computational Efficiency:} Setting $c_{\mathrm{stab}} = 70$, stabilized VNMM-MLPG outperforms cell-based VNMM-Gauss across all six regular benchmark resolutions ($15\times 15$ to $49\times 49$) and five of the six perturbed resolutions ($11.88\%$ vs.\ $24.67\%$ on coarse $15\times 15$ clouds), while delivering $2.0\times\text{--}3.7\times$ faster matrix assembly and $20\%\text{--}24\%$ sparser global operators.
\end{enumerate}

Future work will focus on: (1) extending the closed-form spherical formulation to 3D electromagnetic cavities; (2) incorporating quadratic vector PIM spaces ($\mathcal{P}^2$) to reduce numerical dispersion in high-frequency regimes and provide inherently higher-rank local stiffness; (3) developing split-disk integration, interface-conforming director treatments, and singular corner enrichments for curved boundaries, re-entrant corners, and multi-dielectric media; (4) investigating locally adaptive subdomain scaling factors $\alpha_s(\mathbf{x}_I)$ for multiscale non-uniform point clouds; and (5) implementing open boundary conditions via cell-free absorbing boundary conditions (ABCs) and Perfectly Matched Layers (PML).

%% =========================================================================
%% ACKNOWLEDGMENTS
%% =========================================================================
\section*{Acknowledgments}
This work was supported in part by the Brazilian National Council for Scientific and Technological Development (CNPq), the Minas Gerais State Research Foundation (FAPEMIG), and the Coordination for the Improvement of Higher Education Personnel (CAPES).

%% =========================================================================
%% APPENDIX
%% =========================================================================
\appendix
\section{Closed-Form Geometry, Active Areas, and Moments for Truncated Subdomains}
\label{app:truncated_geometry}

This appendix provides the complete closed-form analytical formulas for evaluating the active physical area $S_0^I \equiv A_{\mathrm{active}}^I = \mathrm{Area}(\Omega_s^I \cap \Omega)$, the active area fraction $\eta_I = A_{\mathrm{active}}^I / (\pi (r_s^I)^2)$, and the first- and second-order geometric moments ($S_{1x}^I, S_{1y}^I, S_{2xx}^I, S_{2yy}^I, S_{2xy}^I$) required by the local stiffness and mass matrices in Section~\ref{subsubsec:truncated_subdomains_analytic} across piecewise planar boundaries.

\subsection{Active Area and Area Fraction Formulas}
Let $A_{\mathrm{cap}}(d, r_s)$ denote the area of a circular segment cut off by a chord at normal distance $d \in [0, r_s]$ from the center of a circular disk of radius $r_s$:
\begin{equation}
    A_{\mathrm{cap}}(d, r_s) = r_s^2 \left[ \arccos\left(\frac{d}{r_s}\right) - \frac{d}{r_s} \sqrt{1 - \left(\frac{d}{r_s}\right)^2} \right].
    \label{eq:cap_area_formula}
\end{equation}
The active physical area $A_{\mathrm{active}}^I$ and area fraction $\eta_I$ are then evaluated in closed form for all geometric configurations:
\begin{enumerate}
    \item \textbf{Case 3 on a Flat Boundary Wall Far from Corners (Case~3 in Fig.~\ref{fig:schematic_mlpg}):} For a node $\mathbf{x}_B \in \partial\Omega$ lying on a straight boundary segment at distance greater than $r_s^B$ from any adjacent corner, $\Omega_s^B \cap \Omega$ is an exact half-disk:
          \begin{equation}
              \eta_B = \frac{1}{2}, \qquad A_{\mathrm{active}}^B = \frac{1}{2} \pi (r_s^B)^2.
              \label{eq:eta_half_disk}
          \end{equation}
    \item \textbf{Case 3 at an Orthogonal Convex Corner Vertex:} For a node situated directly at a $90^\circ$ convex corner vertex, $\Omega_s^B \cap \Omega$ is an exact quarter-disk:
          \begin{equation}
              \eta_B = \frac{1}{4}, \qquad A_{\mathrm{active}}^B = \frac{1}{4} \pi (r_s^B)^2.
              \label{eq:eta_quarter_disk}
          \end{equation}
    \item \textbf{Case 3B for a Boundary Node on a Wall Near a Corner:} When a boundary node $\mathbf{x}_B$ lies on a wall at distance $d_{\parallel} < r_s^B$ from an adjacent perpendicular wall, the subdomain is truncated not only by the host wall (which halves the disk) but also by the transverse wall. Because the center lies on the boundary line, the chord at distance $d_{\parallel}$ cuts the full disk symmetrically, removing an exterior circular cap $A_{\mathrm{cap}}(d_{\parallel}, r_s^B)$, exactly half of which lies inside the half-plane of the domain. Consequently, the active area suffers an additional reduction, falling strictly below a half-disk:
          \begin{equation}
              \eta_B(d_{\parallel}) = \frac{1}{2} \left( 1 - \frac{A_{\mathrm{cap}}(d_{\parallel}, r_s^B)}{\pi (r_s^B)^2} \right) \in \left[ \frac{1}{4}, \frac{1}{2} \right),
              \label{eq:eta_boundary_near_corner}
          \end{equation}
          with $A_{\mathrm{active}}^B = \eta_B(d_{\parallel}) \pi (r_s^B)^2$. Equation~\eqref{eq:eta_boundary_near_corner} provides a continuous, exact transition between the interior-wall half-disk ($\eta \to 1/2$ as $d_{\parallel} \to r_s^B$) and the corner vertex quarter-disk ($\eta \to 1/4$ as $d_{\parallel} \to 0$).
    \item \textbf{Case 2 for an Interior Node Crossing a Single Wall (Case~2 in Fig.~\ref{fig:schematic_mlpg}):} For an interior node $\mathbf{x}_{I_2}$ at distance $d_I = \mathrm{dist}(\mathbf{x}_{I_2}, \partial\Omega) < r_s^{I_2}$ from a single boundary segment, subtracting the cut-off cap yields:
          \begin{equation}
              \eta_{I_2} = 1 - \frac{A_{\mathrm{cap}}(d_I, r_s^{I_2})}{\pi (r_s^{I_2})^2} \in \left( \frac{1}{2}, 1 \right).
              \label{eq:eta_circular_cap}
          \end{equation}
          As $d_I \to r_s^{I_2}$, $\eta_{I_2} \to 1$ smoothly recovering the uncut disk (Case~1 in Fig.~\ref{fig:schematic_mlpg}), whereas as $d_I \to 0$, $\eta_{I_2} \to 1/2$ smoothly matching the flat boundary limit (Case~3).
    \item \textbf{Case 2B for an Interior Node Near a Corner Crossing Two Walls:} When an interior node $\mathbf{x}_I$ is located close to a domain corner such that its distance to both perpendicular walls satisfies $d_x < r_s^I$ and $d_y < r_s^I$, the disk intersects both boundaries simultaneously. Two geometric regimes arise:
          \begin{itemize}
              \item If the corner vertex lies outside the disk ($d_x^2 + d_y^2 \ge (r_s^I)^2$), the two cut-off circular segments in the exterior half-planes are mutually disjoint, yielding:
                    \begin{equation}
                        A_{\mathrm{active}}^I = \pi (r_s^I)^2 - A_{\mathrm{cap}}(d_x, r_s^I) - A_{\mathrm{cap}}(d_y, r_s^I).
                        \label{eq:eta_interior_corner_disjoint}
                    \end{equation}
              \item If the corner vertex lies inside the disk ($d_x^2 + d_y^2 < (r_s^I)^2$), the two exterior half-planes overlap in the exterior corner quadrant $\{x - x_I \le -d_x, \, y - y_I \le -d_y\}$. By inclusion-exclusion:
                    \begin{equation}
                        A_{\mathrm{active}}^I = \pi (r_s^I)^2 - A_{\mathrm{cap}}(d_x, r_s^I) - A_{\mathrm{cap}}(d_y, r_s^I) + A_{\mathrm{overlap}}^I,
                        \label{eq:eta_interior_corner_overlap}
                    \end{equation}
                    where the overlap area $A_{\mathrm{overlap}}^I$ is computed in exact closed form via Green's circulation theorem along the circular arc connecting the horizontal-wall intersection angle $\theta_x = \pi + \arcsin(d_y / r_s^I)$ to the vertical-wall intersection angle $\theta_y = \pi + \arccos(d_x / r_s^I)$:
                    \begin{align}
                        A_{\mathrm{overlap}}^I & = \frac{1}{2} (r_s^I)^2 (\theta_y - \theta_x) \nonumber                              \\
                                               & \quad + \frac{1}{2} r_s^I \left[ d_x (\sin\theta_y - \sin\theta_x) \right. \nonumber \\
                                               & \qquad \qquad \left. - d_y (\cos\theta_y - \cos\theta_x) \right].
                    \end{align}
          \end{itemize}
\end{enumerate}

\subsection{First- and Second-Order Geometric Moments for Truncated Mass Matrices}
To evaluate the local mass matrix $\mathbf{M}_I^{\mathrm{trunc}}$ in Eq.~\eqref{eq:m_truncated_moments}, the first-order moments ($S_{1x}^I, S_{1y}^I$) and second-order moments ($S_{2xx}^I, S_{2yy}^I, S_{2xy}^I$) of $\Omega_s^I \cap \Omega$ relative to $\mathbf{x}_I$ are given analytically by:
\begin{enumerate}
    \item \textbf{Flat Boundary Half-Disk (Case 3):} For a boundary node $\mathbf{x}_B$ on a straight wall with inward domain unit normal $\vec{n}_{\mathrm{in}} = [n_x, n_y]^T \in \{[\pm 1, 0]^T, [0, \pm 1]^T\}$:
          \begin{equation}
              \begin{aligned}
                  S_{1x}^B  & = \frac{2}{3} n_x (r_s^B)^3, \qquad           & S_{1y}^B  & = \frac{2}{3} n_y (r_s^B)^3, \\
                  S_{2xx}^B & = S_{2yy}^B = \frac{\pi}{8} (r_s^B)^4, \qquad & S_{2xy}^B & = 0.
              \end{aligned}
          \end{equation}
    \item \textbf{Orthogonal Convex Corner Quarter-Disk (Case 3 at Vertex):} For a corner vertex node $\mathbf{x}_B$ with inward quadrant signs $\sigma_x, \sigma_y \in \{-1, +1\}$:
          \begin{equation}
              \begin{aligned}
                  S_{1x}^B  & = \frac{1}{3} \sigma_x (r_s^B)^3, \qquad       & S_{1y}^B  & = \frac{1}{3} \sigma_y (r_s^B)^3,           \\
                  S_{2xx}^B & = S_{2yy}^B = \frac{\pi}{16} (r_s^B)^4, \qquad & S_{2xy}^B & = \frac{1}{8} \sigma_x \sigma_y (r_s^B)^4.
              \end{aligned}
          \end{equation}
    \item \textbf{Single-Wall Circular Cap Truncation (Case 2):} For an interior node $\mathbf{x}_I$ at normal distance $d_I \in (0, r_s^I)$ from a straight wall with inward normal $\vec{n}_{\mathrm{in}}$, let $a_I = \sqrt{(r_s^I)^2 - d_I^2}$ denote the half-chord length and $\theta_0 = \arccos(d_I / r_s^I)$. Subtracting the exterior circular cap from the full disk yields the exact linear moments along the tangential ($\parallel$) and inward normal ($\perp$) axes:
          \begin{equation}
              S_{1\parallel}^I = 0, \qquad S_{1\perp}^I = +\frac{2}{3} \left((r_s^I)^2 - d_I^2\right)^{3/2} = \frac{2}{3} a_I^3,
          \end{equation}
          and the second-order moments:
          \begin{align}
              S_{2\parallel}^I      & = \frac{\pi}{4}(r_s^I)^4 - \left[ \frac{(r_s^I)^4}{4}\theta_0 - \frac{d_I a_I}{12}\left(5(r_s^I)^2 - 2d_I^2\right) \right], \\
              S_{2\perp}^I          & = \frac{\pi}{4}(r_s^I)^4 - \left[ \frac{(r_s^I)^4}{4}\theta_0 + \frac{d_I a_I}{4}\left((r_s^I)^2 - 2d_I^2\right) \right],  \\
              S_{2\parallel\perp}^I & = 0,
          \end{align}
          which smoothly interpolate between the flat-wall half-disk ($d_I \to 0$) and the uncut interior disk ($d_I \to r_s^I$).
\end{enumerate}

%% =========================================================================
%% REFERENCES
%% =========================================================================
\bibliographystyle{elsarticle-num-names}
\bibliography{references}

\end{document}
